\documentclass[11pt,a4paper]{article}
\usepackage[margin=25mm]{geometry}
\usepackage{amsmath,amssymb,amsthm,mathtools}
\usepackage{booktabs,longtable,array,enumitem}
\usepackage[expansion=false]{microtype}
\usepackage{xcolor}
\usepackage[colorlinks=true,linkcolor=blue!45!black,citecolor=blue!45!black,urlcolor=blue!45!black]{hyperref}
\usepackage{fancyhdr}
\pdfmapfile{+cm.map}
\pdfmapfile{+cmextra.map}
\pdfmapfile{+symbols.map}
\pdfmapfile{+latxfont.map}
\setlength{\emergencystretch}{3em}
\setlength{\parskip}{4pt}
\setlength{\parindent}{0pt}
\setlist{itemsep=3pt,topsep=5pt,leftmargin=*}
\pagestyle{fancy}\fancyhf{}\fancyhead[L]{\small GCT\_NSRH: checked reductions and open gates}\fancyfoot[C]{\thepage}
\setlength{\headheight}{14pt}
\newcommand{\Q}{\mathbb Q}\newcommand{\Z}{\mathbb Z}\newcommand{\C}{\mathbb C}\newcommand{\N}{\mathbb N}\newcommand{\F}{\mathbb F}
\newcommand{\R}{\mathbb R}
\newcommand{\A}{\mathbb A}\newcommand{\PP}{\mathbb P}\newcommand{\GL}{\mathrm{GL}}
\newcommand{\Hom}{\operatorname{Hom}}\newcommand{\Spec}{\operatorname{Spec}}\newcommand{\Tr}{\operatorname{Tr}}\newcommand{\rank}{\operatorname{rank}}\newcommand{\Sym}{\operatorname{Sym}}\newcommand{\Res}{\operatorname{Res}}\newcommand{\IC}{\operatorname{IC}}\newcommand{\Gr}{\operatorname{Gr}}\newcommand{\Rep}{\operatorname{Rep}}\newcommand{\Frob}{\operatorname{Frob}}\newcommand{\End}{\operatorname{End}}
\newcommand{\id}{\operatorname{id}}\newcommand{\et}{\mathrm{\acute et}}\newcommand{\boxt}{\boxtimes}
\newcommand{\mathscr}{\mathcal}
\renewcommand{\S}{\ensuremath{\mathsection}}
\renewcommand{\textbullet}{\ensuremath{\bullet}}
\newcommand{\fid}[1]{\textsf{#1}}\newcommand{\gate}[1]{\textsf{#1}}
\newtheoremstyle{named}{8pt}{5pt}{\itshape}{} {\bfseries}{.}{.5em}{}
\theoremstyle{named}\newtheorem{finding}{Finding}[section]
\theoremstyle{definition}\newtheorem{definition}[finding]{Definition}
\theoremstyle{remark}\newtheorem{remark}[finding]{Remark}
\title{Purity from path sums and the nonstandard GCT positivity problem\\[4pt]\large Checked reductions, explicit barriers, and the missing realization}
\author{Research run \texttt{GCT\_NSRH\_20261121}}
\date{21 November 2026 (Europe/Berlin)}
\begin{document}\maketitle
\begin{abstract}
We investigate whether the path-quadric sign mechanism in Alp\"oge's supplied \emph{Remark} transfers to the nonstandard GCT positivity problem. Iterated Parseval gives an exact error bound for path sums over $\A^d/\F_q$. Tensor powers and a finite-spectrum argument then give the upper weight $a+d$; visible Tate duality gives equality. The stronger criterion F24 allows signed Euler multiplicities on the input at even path lengths, with an exact visibility condition when finite extensions vary. Counterexamples delimit the hypotheses.

The supplied GCT versions give a nonstandard-RH program without specifying its finite-field objects or weights. We test concrete ways of filling that gap. The actual rank-three algebra has an explicit finite-flat coherent center; regular matrices realize its actual 64-dimensional tensor action. Its order-two collision, coherent extensions and commutant defect are computed. Constant nearby-cycle purity loses that collision order. More sharply, one normalized actual projector-word trace is non-Laurent and cannot be a fixed finite Frobenius trace under extension-by-extension substitution of the quantum parameter. An explicit Artin stack realizes that all-extension scalar trace, but its cohomology is unbounded and its actual dual has unbounded spectrum. Reciprocal trace realizations do not give tensor inverses in the usual derived Weil category. The actual tensor lattice admits unit coherent Frobenius over the parameter base and on its normalization, while descent through the singular center has an explicit torsion defect. The direct collision algebra fails the smoothness hypothesis of the specified noncommutative Weil framework, but an explicit Auslander enlargement retains the actual algebra as a corner and its tensor module as a Hom space; its category invariant has four unit Frobenius eigenvalues. Derived lifts and a four-character model give complementary, carefully scoped results in rank one.

Ordinary geometric and motivic Satake already provide pure Tate coefficient objects for classical plethysm and Kronecker multiplicities. Their projector ranks and designated bases remain additional data: even a small scalar-Frobenius multiplicity space has no signed-permutation basis. Explicit orbit and interlacing criteria identify sufficient positive rules, while lattice and basis counterexamples exclude abbreviated deductions. Primitive Grassmannian motives and compressed marked-tree witnesses give concrete results for two-row rectangular coefficients. The known Kahle--Micha{\l}ek counterexample already refutes the supplied homogeneous Ehrhart PH1; a separate compact-input growth argument rules out its uniform dimension bound even on a positive simplex family. The exact known rectangular and no-occurrence theorems are retained. Every checked finding is translated into schemes, representation $K_0$ or coherent $G_0$, $\lambda$ operations, Weil realizations and motives, with exact source and archival locators.

\end{abstract}
\begin{center}\fbox{\parbox{.91\linewidth}{\textbf{Research status: work in progress.} The conditional theorems and counterexamples below are proved and checked. No nonstandard RH for the designated GCT objects, general canonical-basis positivity, general positive plethysm rule, or complexity separation is proved. The final part records the exact open gates.}}
\end{center}
\newpage\setcounter{tocdepth}{1}{\small\tableofcontents}\newpage

\part{The exact question and its sources}
\section{Scope, premises, and conventions}\label{sec:scope}
The main question is a transfer question. We accept the RH and Dirichlet-GRH results in the supplied Alp\"oge material as stipulated premises, as requested. We do not re-referee those results. Their acceptance does not supply a theorem relating an automorphic or finite-field trace system to a nonstandard quantum canonical basis.

The relevant finite-field description is in \cite[printed p.~3, footnote~3, attached to the paragraph before the timeline; also step~3]{Alpoge}. Timeline step~5 itself concerns a von Mangoldt functional and continues onto p.~4. We use ``the step-5 mechanism'' in the prompt's sense, but retain this more precise locator. The note names a path-quadric identity, middle-cohomological signs, iterated Parseval, and the absence of cancellation of inverse roots. It does not state a standalone theorem with all the geometric hypotheses used below.

The supplied \texttt{ALPOGE\_Main.lean} consists of five exported wrappers, including the RH wrapper \texttt{riemannHypothesis}. Its body is \texttt{rhOfAll}; the imported module is \texttt{Supplements.Kernels.\allowbreak MainTheoremsAll}. The supplied \texttt{ALPOGE\_Statements.lean} defines the zeta/Dirichlet zero and nonvanishing propositions. Neither file contains a nonstandard quantum-group, path-concentration, or theta/GCT transfer theorem. The imported development is not in this ZIP. The supplied scan records zero occurrences for the listed prohibited constructs; this run did not repeat that scan or execute the missing development. These are access and locator facts, not objections to the stipulated theorem. No Masullo result is used: none is identified or supplied in this prompt.

We write $v$ for the \emph{formal quantum parameter}, $q$ for a fixed prime power, and $Q=q^n$ for an extension-field cardinality. In the GCT sources the quantum parameter is usually called $q$. No substitution $v=Q$ or $v=Q^{1/2}$ constructs a Weil structure. All weights in the purity conclusion are integers. The Frobenius convention is geometric Frobenius, so $\Q_\ell(-1)$ has eigenvalue $q$ over $\F_q$.

``Finding'' means a deduction established for this investigation, including elementary lemmas in the precise form needed. It is not a claim of priority over the mathematical literature. Known research results are cited, not presented as new findings. A source statement checked in context is not thereby an independently checked proof of the entire cited theorem.

\section{What IV, VI, VII and VIII actually specify}\label{sec:sources}
\subsection{Versions and the missing RH statement}
The attached IV is the revised Blasiak--Mulmuley--Sohoni paper, arXiv cs/0703110v4 (2013), not merely the 2007 version referred to by VII and VIII. It proves nonstandard Schur--Weyl duality and the two-row crystal result. VI is arXiv 0704.0229v4 (2009); VII and VIII are 0709.0749v2 and 0709.0751v2 (2008). Rendered title-page dates produced by an undated \verb|\maketitle| are not used as version dates.

\textbf{Source finding A01.} In these four supplied versions there is no formulated nonstandard RH specifying finite-field objects, Frobenius actions, cohomological degrees and weights. This is an exact-source conclusion, not a theorem about what other work might contain. The relevant locators are:

{\small\begin{longtable}{p{.13\linewidth}p{.35\linewidth}p{.44\linewidth}}
\toprule Work&Exact locator&Content and status\\\midrule\endhead
IV&\S1.1, Problem 1.1; \S\S1.2--1.5; \S19&Positive-rule and basis program; the revised text has no finite-field purity statement.\\
VI&Introduction, Figure 1.1; \S1.7, Observation 1.7.1 and Figure 1.2&The proposed RH/extensions connection is marked by question marks; standard canonical positivity motivates the program.\\
VI&\S4.2, Proposition 4.2.1; \S4.2.2, Question 4.2.6 (PH0??); \S4.3&Asks for compatible canonical bases and analogous positivity. No nonstandard weight theorem is stated.\\
VII&\S1 after Conjectures 1.3--1.4; \S4.1, equations (9)--(10), Theorem 4.4; \S4.2, Theorem 4.5&Nonstandard minors and Laplace coefficients need interpretations analogous to perverse-sheaf constructions.\\
VIII&\S1, pp.~2--3; \S2.3 after Conjecture 2.5&An extension of perverse sheaves and finite-field RH is proposed as a way to understand minors and globalization.\\
VIII&\S2.4.2 after Conjecture 2.7; \S2.4.3 after Conjecture 2.9&Topological interpretations, Hard Lefschetz, and a theory for noncommutative, nonflat objects are suggested; the general case includes correction terms.\\\bottomrule
\end{longtable}}
These statements are supported by complete reading of their contexts and matching numbered PDF statements, not merely a search for the phrase ``Riemann hypothesis.'' Their appropriate conclusion is that a precise realization/purity package must first be specified.

\subsection{The algebraic objects}
Let $H$ be a complex connected classical reductive group, $X=V_\mu(H)$, $G=\GL(X)$, and $\rho:H\to G$ the representation homomorphism. The generalized plethysm multiplicity is
\begin{equation}\label{eq:plethysm}
a_{\lambda,\mu}^{\pi}=\dim\Hom_H\bigl(V_\pi,\Res_\rho V_\lambda(G)\bigr).
\end{equation}
For Kronecker coefficients, $H=\GL(V)\times\GL(W)$, $X=V\otimes W$, and $\rho(g,h)=g\otimes h$.

VII, Problem~1.1 and Theorems~1.2, 2.1 and 3.1, constructs a nonstandard matrix-coordinate bialgebra
\[
 S=\mathcal O(M_v^H(X)),\qquad \mathcal A_r^H=(S_r)^*,
\]
with a quantized map from $H_v$ to the associated nonstandard quantum object. Its matrix generators satisfy the projector relation
\[
 P^{+,H}_{X,X}(\mathbf u\otimes\mathbf u)
 =(\mathbf u\otimes\mathbf u)P^{+,H}_{X,X}
\]
(VII, \S2, equations (4)--(6)). The projector is constructed from the braiding on $X_v\otimes X_v$. The determinant may vanish, so the strict term may be quantum semigroup; an antipode is not automatic. The dual nonstandard symmetric-group deformation is $\mathcal B_r^H(v)$ (VII, \S5).

These are noncommutative algebras over fields containing $\Q(v)$, sometimes an algebraic extension. Their bar involution is not arithmetic Frobenius. Complete reducibility or a compact real form is not a Weil structure. Nonflatness and noncommutativity have different roles: nonflat commutative schemes still have \(\ell\)-adic sheaves, whereas a noncommutative coordinate algebra is not the coordinate ring of a scheme in the EGA/SGA sense.

General plethysm duality and reciprocity are VII, Conjectures~1.3--1.4 and refined Conjecture~6.1. They would give
\begin{equation}\label{eq:comparison}
a_{\lambda,\mu}^{\pi}=\sum_\alpha m_\lambda^\alpha n_\pi^\alpha,
\end{equation}
where $n_\pi^\alpha$ is an $H_v$ constituent multiplicity and $m_\lambda^\alpha$ is a Specht multiplicity after specialization at $v=1$. In revised IV, Theorems~1.5/12.1 establish the Kronecker duality, and \S1.4, equation~(11), gives the corresponding Kronecker decomposition. It would be incorrect to leave that particular comparison conjectural merely because VII predates the revision.

\subsection{The actual positivity targets}\label{sec:targets}
The basis constructions themselves have separate prerequisites. VIII, Conjecture~2.2, concerns compatible local crystal bases; Question~2.3's ordinary balanced-triple condition fails in examples; Conjecture~2.5 proposes a unique minimum-degree globalization. Conjectures~2.6 and 3.3 concern cells, Conjectures~3.1--3.2 the dual globalization, and Conjecture~3.7 the \emph{opposite} basis's quasi-subcells and specialization. VIII's Proposition~2.4 contains an omitted proof and assumes the local-crystal conjecture; we do not record it as independently proved.

\paragraph{Kronecker, coordinate side.}
VIII, Conjecture~2.7, asks that expansion coefficients of $CB(\mathcal A_r^H)$ in the external canonical basis $CB(Z)$, $Z=\End(X_v^{\otimes r})$, be polynomials with nonnegative \emph{rational} coefficients (equation~(9)). Multiplication and comultiplication constants have the form
\begin{equation}\label{eq:kroneckerpositive}
 \pm(v-v^{-1})^a f(v),\qquad a\in\Z_{\ge0},
\end{equation}
where $f$ is bar-invariant, coefficientwise nonnegative rational and unimodal. The same is proposed for $CB(S)$.

\paragraph{Kronecker, dual side.}
VIII, Conjecture~3.4, applies to structure constants and external expansion coefficients of $E(r)$ and $E^{\mathrm{opp}}(r)$. Its form is
\begin{equation}\label{eq:dualpositive}
 \pm(v^{1/2}-v^{-1/2})^a f(v),
\end{equation}
with $f$ nonnegative rational, bar-invariant and unimodal in the half-power variable. Thus raw coefficientwise positivity is not the stated assertion even in the Kronecker case.

\paragraph{General plethysm.}
VIII, Conjectures~2.8 and 3.5, concern normalized coefficients of minimal polynomials of algebraic structure/expansion coefficients. The displayed normalizations are equations~(10) and (14):
\[
 (-1)^e(v-v^{-1})^{e'}c(v),\qquad
 (-1)^e(v^{1/2}-v^{-1/2})^{e'}g(v).
\]
The nonzero normalized polynomial is required to be positive at $v=1$ (called saturated there). Conjectures~2.9 and 3.6 ask for a dominant positive unimodal part and a small correction, for example
\[
 c=c^0+c^1,\qquad
 \frac{\|c^1\|_2}{\|c^0\|_2}
 \le\frac1{\operatorname{poly}(\langle\mu\rangle,\langle\pi\rangle,\langle r\rangle)},
 \qquad c^0(1)\gg |c^1(1)|.
\]
Conjecture~3.6 is explicitly informal. A theorem target must fix the quantitative meaning of the polynomial and the asymptotic comparison. VIII, \S5.4.1, Figure~30, displays negative entries in a normalized example; these are cited source computations, not a recomputation of the whole algebra in this run.

There is also a sign-convention issue in ``coefficients of the minimal polynomial'': the usual monic polynomial of 1 is $T-1$, while the source later refers to trace and norm without the conventional sign on trace. We record this ambiguity rather than use it to disprove the intended conjecture. ``Frobenius automorphisms'' in VIII, \S5.4.1, refers to algebraic conjugations over $\Q(v)$, not to finite-field Weil Frobenius.

\subsection{Positivity hypotheses in the flip}
VI, Hypothesis~1.6.4 (PH1), is substantially stronger than numerical nonnegativity. For each triple there should be a rational polytope $P_{\lambda,\mu}^{\pi}$ whose Ehrhart function is
\[
 \widetilde a_{\lambda,\mu}^{\pi}(t)=a_{t\lambda,\mu}^{t\pi},
\]
with $\mu$ fixed (VI, \S1.6, definitions preceding Theorem~1.6.1). It must admit a description $Ax\le b$ with $A$ independent of $\lambda,\pi$ and $b$ homogeneous linear in them; ambient dimension and entry bitlengths are polynomially bounded as specified there; membership and a separating hyperplane are computable in polynomial time. The number of inequalities need not be polynomial.


\paragraph{The attached homogeneous PH1 is already false.}
Kahle--Micha{\l}ek \cite[Theorem~5 and Remark~6, p.~4]{KM}
prove that the fixed-inner-weight stretching function
\[
 \phi(s)=a_{(4s),(3)}^{(7s,5s)}
       =\frac{s+r(s)}3,\qquad
 (r(0),\ldots,r(5))=(3,-1,1,0,2,-2),
\]
with $r$ periodic modulo six, is not the Ehrhart function of a
rational polytope. Its degree is one, but $\phi(1)=0$ and
$-\phi(-1)=1$, contradicting the interior-versus-total inequality
in Ehrhart reciprocity. Their p.~2 explicitly identifies the target
as VI, Hypothesis~1.6.4. By the rectangular identification on their
p.~6, the same function is
\[
 g((4s,4s,4s),(4s,4s,4s),(7s,5s)).
\]
This is a cited prior refutation, not a new finding of this run.
Their pp.~3--4 distinguish the later 2011 formulation allowing
inhomogeneous families $Ax\le sb+c$. The cited theorem does not
exclude every dimension of that later formulation. We therefore
do not leave the old homogeneous assertion as an open gate, or
equate its failure with a refutation of every later positivity
program. Finding~\fid{F34} separately tests the compact dimension
requirement on a family with a genuine positive simplex model.

VI, Hypothesis~1.6.5 (SH), asks for a polynomial bound on the saturation index in $\dim G+\operatorname{ht}\lambda+\operatorname{ht}\pi$ and asymptotic generic strict saturation. Hypothesis~1.6.6 (PH2) bounds the positivity index; Hypothesis~1.6.7 (PH3) bounds the modular index of a positive rational Hilbert-series form. Strict saturation in every case is not the attached revised hypothesis.

The conditional implication to relaxed PHflip is VI, Theorem~1.6.8. A printed discrepancy must be retained: its part~(a) displays $a_{b\lambda,b\mu}^{b\pi}$ although the stretching definition and its stated proof through PH1 keep $\mu$ fixed. We use only the fixed-$\mu$ consequence directly justified by the displayed Ehrhart model and record the all-three-scaling formulation as requiring clarification. This editorial point is not needed for our main obstruction. More fundamentally, VIII, \S6, explicitly asks for additional effective-label properties. An abstract canonical basis or an abstract positive count does not prove PH1, SH, or the existence of a lower-bound obstruction.

\section{The Kazhdan--Lusztig template, with weights fixed}\label{sec:KL}
Let $G/\F_Q$ be split reductive, $B$ a Borel subgroup, $C_w=BwB/B$ and $X_w=\overline{C_w}$. Put
\[
 \IC_w^u=j_{!*}\Q_\ell[\ell(w)][-\ell(w)]
\]
with the usual extension to $X_w$; the notation records the unshifted convention, constant in degree zero on $C_w$. Then
\begin{equation}\label{eq:KL}
 P_{x,w}(T)=\sum_{j\ge0}\dim\mathcal H^{2j}(\IC_w^u)_x\,T^j.
\end{equation}
Odd stalk groups vanish, and each eigenvalue in degree $2j$ is $Q^j$. This is the Tate/parity statement, stronger than an absolute-value bound. Exact primary locators for this form are Haines \cite[Theorems 2.0.1, 5.0.4; Proposition 3.0.2; \S4]{Haines}. The original KL geometric identification is \cite[Theorem 4.3]{KL}.

In the modern proof, a Bott--Samelson resolution is proper, its fibers admit affine pavings, and the decomposition theorem isolates the intersection complex. The affine paving provides even/Tate cohomology. The compatibility with Weil structures must be checked: a geometric splitting need not itself be Frobenius-equivariant. Haines \S4 discusses precisely this issue. His introduction also notes that the original KL proof does not require the decomposition theorem; the latter is the requested modern template, not a claim about historical logical necessity.

Deligne's upper-weight theorem is \cite[Theorem 3.3.1; Corollaries 3.3.4--3.3.6]{WeilII}. BBD intermediate extension preserves purity (5.3.2), pure perverse sheaves are geometrically semisimple (5.3.8), and pure complexes split geometrically (5.4.5); relative Hard Lefschetz is 5.4.10 \cite{BBD}. None of these statements identifies an arbitrary pure trace with the dimension polynomial \eqref{eq:KL}. Arithmetic Frobenius semisimplicity is distinct: BBD, pp.~138--139, explicitly discusses pure unipotent Jordan-block sheaves.

{\small\begin{longtable}{p{.25\linewidth}p{.29\linewidth}p{.38\linewidth}}
\toprule Standard ingredient&Nonstandard datum supplied&Unfilled role\\\midrule\endhead
Flag schemes, Schubert strata&$H_v$, nonstandard FRT bialgebra; support-defined cells&An actual space/category with those strata and stalk functors.\\
Distinguished $\IC_w$ and its normalized class&Candidate minimum-degree bases&Objects whose classes are the \emph{specified} bases; existence and uniqueness.\\
Weil structure and closed-point Frobenius&Algebraic bar involution&Arithmetic models, traces, degrees and integer weights.\\
Proper resolutions/convolution&Algebra multiplication and comultiplication&Correspondences or functors with a decomposition theorem.\\
Tate parity and effective multiplicities&Signed rational forms in VIII&Positive multiplicity objects, signs, factors and denominators.\\
Verdier duality and Hard Lefschetz&Bar symmetry/unimodality conjectures&Compatible duality and Lefschetz isomorphisms.\\
Crystal/path labels&Conjectural cells; proved two-row tableaux&Effective labels, specialization and polynomial-time verification in general.\\\bottomrule
\end{longtable}}
For standard quantum groups, quiver and flag constructions supply related geometric realizations; it would be misleading to assert that every quantum canonical-basis constant is a stalk of one Schubert variety.

\part{Checked findings}
\section{The sign mechanism and its exact strength}\label{sec:main}
This section contains the main proved reduction. Its hypotheses explicitly retain the information that the phrase ``middle cohomology'' compresses.

\begin{finding}[\fid{F01}: path identity and Parseval estimate]\label{f:path}
Let $V=\F_Q^d$, $d\ge1$, with a nondegenerate bilinear pairing, and let $\psi$ be a nontrivial additive character. For $k\ge2$ and any function $t:V\to\C$, put
\[
 M=\sum_{x\in V}t(x),\quad
 P_k(x_1,\ldots,x_k)=\sum_{i=1}^{k-1}\langle x_i,x_{i+1}\rangle,
 \quad N_k=\sum_{P_k=0}\prod_{i=1}^k t(x_i).
\]
Then $QN_k=M^k+R_k$, where
\[
 R_k=\sum_{h\in\F_Q^\times}\sum_{x_1,\ldots,x_k\in V}
       \prod_i t(x_i)\,\psi(hP_k),
\]
and
\begin{equation}\label{eq:parseval}
 |R_k|\le (Q-1)Q^{d(k-1)/2}\|t\|_2^2\|t\|_\infty^{k-2}.
\end{equation}
No nonnegativity assumption on $t$ is required.
\end{finding}
\begin{proof}
Orthogonality gives $\sum_h\psi(hz)=Q\,1_{z=0}$; separating $h=0$ gives the identity. For $h\ne0$, set $K_h(x,y)=\psi(h\langle x,y\rangle)$ and $D=\operatorname{diag}(t)$. Nondegeneracy gives
\[
 (K_hK_h^*)_{x,x'}=\sum_y\psi(h\langle x-x',y\rangle)
 =Q^d1_{x=x'},
\]
so $\|K_h\|=Q^{d/2}$. The inner sum is
$t^{\mathsf T}K_hD K_h\cdots D K_ht$, with $k-1$ kernels and $k-2$ diagonal factors. The norm of the row functional $t^{\mathsf T}$ is $\|t\|_2$, even for complex $t$ (it is the Hermitian dual of $\bar t$). Cauchy--Schwarz and submultiplicativity give the bound for one $h$; sum the $Q-1$ terms. For $k=2$ there are no diagonal factors.
\end{proof}
This refines Fable suggestion~3's scalar toy in the exact form needed here. It is an open-path matrix coefficient, not in general a trace moment: tracing the power closes the path into a cycle.

\begin{finding}[\fid{F02}: finite spectral no-cancellation]\label{f:spectral}
For distinct nonzero complex numbers $\alpha_j$ and nonzero constants $c_j$, let $p_n=\sum_jc_j\alpha_j^n$. For a nonempty finite spectrum,
\[
 \limsup_{n\to\infty}|p_n|^{1/n}=\max_j|\alpha_j|.
\]
Consequently $|p_n|\le Cn^A R^n$ for all $n\ge1$, $R>0$, implies $|\alpha_j|\le R$. The same holds for any subexponential prefactor.
\end{finding}
\begin{proof}
Near zero,
\[
 \sum_{n\ge1}p_nz^n=\sum_j\frac{c_j\alpha_jz}{1-\alpha_jz}.
\]
Every $1/\alpha_j$ is a genuine pole, since the other summands are regular there. The radius of convergence is $1/\max|\alpha_j|$. Cauchy--Hadamard proves the equality; the bound gives radius at least $1/R$. The empty case is vacuous. Coincident spectral values must first be grouped; the conclusion concerns their \emph{nonzero} total coefficients.
\end{proof}
Positive multiplicities are one sufficient way to ensure that regrouping does not erase a root. Eigenvalues themselves need not be positive, or real. The traces of powers do not see nilpotent Jordan off-diagonal entries, so semisimplicity is not required for this lemma.

\begin{finding}[\fid{F03}: tensor-path purity criterion]\label{f:criterion}
Fix a prime power $q>1$, an integer $a$, $d\ge1$, and $B>0$. For every $n\ge1$ let $Q=q^n$ and $t_n:\F_Q^d\to\C$ satisfy $|t_n(x)|\le BQ^{a/2}$. Form $M_n,N_{k,n},R_{k,n}$ as in \fid{F01}. Assume:
\begin{enumerate}[label=(S\arabic*)]
\item A fixed finite multiset $\mathcal A$ of nonzero algebraic numbers, with positive integer multiplicities, and $\epsilon\in\{1,-1\}$ satisfy
\[
 M_n=\epsilon\sum_{\alpha\in\mathcal A}\alpha^n\qquad(n\ge1).
\]
\item For an unbounded set of integers $k\ge2$ there is, for each $k$, a fixed finite multiset $\mathcal B_k$ of nonzero algebraic numbers, possibly empty, with positive integer multiplicities such that
\[
 q^nN_{k,n}=-\epsilon^k\sum_{\beta\in\mathcal B_k}\beta^n\qquad(n\ge1).
\]
\end{enumerate}
Then $|\alpha|\le q^{(a+d)/2}$ for every $\alpha\in\mathcal A$.

For the all-conjugate assertion, require for every embedding $\sigma$ of the algebraic spectral data a corresponding trace system $t_{n,\sigma}$ satisfying these hypotheses with spectra $\sigma(\mathcal A),\sigma(\mathcal B_k)$ and a constant $B_\sigma$ independent of $n$. If in addition
\begin{equation}\label{eq:duality}
 \mathcal A=q^{a+d}/\mathcal A
 \quad\text{as multisets},
\end{equation}
then every eigenvalue has every complex absolute value $q^{(a+d)/2}$: the spectrum is pure of weight $a+d$.
\end{finding}
\begin{proof}
The pointwise bound gives $\|t_n\|_2^2\le B^2Q^{a+d}$. Equation~\eqref{eq:parseval} therefore gives
\begin{equation}\label{eq:overhead}
 |R_{k,n}|\le B^kQ^{k(a+d)/2+d/2+1}.
\end{equation}
By the exact path identity and (S1)--(S2),
\begin{equation}\label{eq:positivespectrum}
 -\epsilon^kR_{k,n}=
 \sum_{\beta\in\mathcal B_k}\beta^n+
 \sum_{(\alpha_1,\ldots,\alpha_k)\in\mathcal A^k}
 (\alpha_1\cdots\alpha_k)^n.
\end{equation}
After grouping equal roots all coefficients are positive integers. In particular the diagonal tuple contributes $\alpha^k$ and cannot be cancelled. For each \emph{fixed} $k$, \fid{F02} and \eqref{eq:overhead} give
\[
 |\alpha|^k\le q^{k(a+d)/2+d/2+1},\qquad
 |\alpha|\le q^{(a+d)/2+(d/2+1)/k}.
\]
Let $k$ tend to infinity through the unbounded set. This order of limits is essential: constants such as $B^k$ disappear in the $n$th-root limit before $k$ varies. Applying the same argument to each $\sigma$ gives all-conjugate upper bounds. Since $a+d$ is an integer, $q^{a+d}$ is rational and fixed by every $\sigma$. Apply the upper bound to $q^{a+d}/\sigma(\alpha)$, which occurs by \eqref{eq:duality}, to obtain the reverse inequality.
\end{proof}
The purity conclusion is conditional on a \emph{specified Tate duality}. An early draft allowed an arbitrary algebraic dualizing scalar; the independent check found that its conjugates could have different absolute values. The final theorem uses the integer-weight scalar $q^{a+d}$ and the correction is recorded in the audit.

\begin{finding}[\fid{F04}: controlled background terms]\label{f:errors}
Replace (S2) by
\[
 q^nN_{k,n}=-\epsilon^k\sum_{\beta\in\mathcal B_k}\beta^n+E_{k,n}.
\]
Suppose $|E_{k,n}|\le C_kn^{A_k}q^{n(k(a+d)/2+\eta_k)}$ for all $n$, with $\eta_k/k\to0$ along the unbounded set of $k$. Then \fid{F03}'s upper bound, and its purity conclusion under the same all-embedding duality hypotheses, still hold, provided that the modified identity and error estimate also hold for every corresponding trace system $t_{n,\sigma}$, with constants $C_{k,\sigma},A_{k,\sigma}$ and $\eta_{k,\sigma}/k\to0$ allowed to depend on $\sigma$.
\end{finding}
\begin{proof}
The positive spectral sum in \eqref{eq:positivespectrum} is now $\epsilon^k(E_{k,n}-R_{k,n})$. Its exponential rate for fixed $k$ is at most
$q^{k(a+d)/2+\max(d/2+1,\eta_k)}$. Apply \fid{F02} and take the $k$th root. Since $\max(d/2+1,\eta_k)/k\to0$, the same limit follows. Duality supplies the reverse bound as before.
\end{proof}
This is a sufficient error budget, not a necessity theorem for every possible proof. It permits suitable removed Tate or boundary terms; it does not justify discarding uncontrolled terms.

\begin{finding}[\fid{F17}: finite extensions and a separate dual system]\label{f:extensiondual}
The upper-weight conclusion of \fid{F03}--\fid{F04} still holds if, for each allowed $k$, the opposite-sign identity and its error estimate are known only for $n=e_kr$, all $r\ge1$, where the positive integer $e_k$ may depend on $k$.

Self-duality can also be replaced by a separate dual coefficient system. Suppose the first spectrum $\mathcal A$ has the proved upper weight $a+d$, a second spectrum $\mathcal A^D$ has the proved upper weight $d-a$, and a perfect Frobenius-equivariant pairing identifies
\[
 \mathcal A^D=\{q^d/\alpha:\alpha\in\mathcal A\}
\]
with multiplicities. Then $\mathcal A$ is pure of weight $a+d$, provided both upper bounds hold under every embedding.
\end{finding}
\begin{proof}
For fixed $k$, apply the spectral lemma in the variable $r$. The positive multiset now contains $\alpha^{e_kk}$; equal powered roots still have positive multiplicities, so none disappears. Its exponential bound is
\[
 (q^{e_k})^{k(a+d)/2+d/2+1}
\]
(or the corresponding maximum with $\eta_k$). Taking the $(e_kk)$th root gives exactly the same bound on $|\alpha|$ as before, and the unbounded-$k$ limit is unchanged. Constants may depend on $k$ and $e_k$. This uses full multiples of $e_k$, not an arbitrary infinite subsequence such as the odd degrees in \fid{F05}(d).

For the second assertion apply the dual upper bound to $q^d/\alpha$:
\[
 |\alpha|\ge\frac{q^d}{q^{(d-a)/2}}=q^{(d+a)/2}.
\]
Combine it with the first upper bound. Since $q^d$ is rational, the same calculation holds under every embedding.
\end{proof}
This refinement allows a geometric construction to choose a finite field of definition separately for each path length. The upper weight $d-a$ for the dual system is an independent hypothesis: the original pointwise upper weight $a$ alone does not give the required upper estimate for the dual coefficient system. The refinement also distinguishes Verdier duality with a dual coefficient system from an assertion that a single compact-support module is self-dual. To use the latter option geometrically, the necessary compact-support/ordinary-cohomology comparison still has to be proved; it is not supplied by dimension alone.

\begin{finding}[\fid{F24}: signed Euler spectra at even path lengths]\label{f:signed}
Keep the pointwise setup and path identity of \fid{F03}, and put $w=a+d$.
Replace its input sign hypothesis by
\[
 M_n=\sum_{\alpha\in A}c_\alpha\alpha^n\quad(n\ge1),
 \qquad c_\alpha\in\mathbb R\setminus\{0\},
\]
where $A$ is finite, consists of distinct nonzero roots, and is nonempty.
Thus $A$ is the \emph{visible spectrum}, after grouping equal roots and deleting
zero coefficients. Suppose that for an unbounded set of \emph{even} $k\ge2$,
\[
 q^nN_{k,n}=-\sum_{\beta\in B_k}b_{k,\beta}\beta^n+E_{k,n},
 \qquad b_{k,\beta}>0,
\]
for every $n\ge1$, with the error bound of \fid{F04}. Then every
$\alpha\in A$ satisfies $|\alpha|\le q^{w/2}$. The positive coefficients
$b_{k,\beta}$ may be real; positive integer multiplicities suffice.

The all-embedding conclusion holds when each embedding has the corresponding
trace system, real transformed input coefficients, and error bounds. Integer
Euler multiplicities meet the coefficient condition. A reciprocal relation
$A=q^w/A$ among the visible supports, or a separately bounded visible dual
system as in \fid{F17}, supplies purity of the visible spectrum.

If the identities are available only at $n=e_kr$, the same conclusion holds
provided that the largest original radius remains visible after powering
by $e_k$, for an unbounded set of the allowed even $k$. This extra condition
cannot be omitted for signed spectra.
\end{finding}
\begin{proof}
Let $R=\max_{\alpha\in A}|\alpha|$ and
$A_R=\{\alpha:|\alpha|=R\}$. In the group algebra of the multiplicative group
generated by $A_R$, form
\[
 F=\sum_{\alpha\in A_R}c_\alpha[\alpha],\qquad
 F^k=\sum_\rho d_{k,\rho}[\rho],
\]
grouping equal products. The coefficients $d_{k,\rho}$ are real and
\[
 \sum_\rho d_{k,\rho}
   =\left(\sum_{\alpha\in A_R}c_\alpha\right)^k\ge0
\]
because $k$ is even. Moreover $F^k\ne0$. Otherwise evaluating
$[\alpha]\mapsto\alpha^n$ would give
$(\sum_{A_R}c_\alpha\alpha^n)^k=0$ for every $n\ge1$.
The Vandermonde matrix for $n=1,\ldots,|A_R|$ is invertible, contradicting
the nonzero coefficients on distinct nonzero roots. This argument uses
nonnilpotence of $F$, not a false assertion that a group algebra with
torsion has no zero divisors. Since a nonzero finite collection of
nonpositive real numbers has negative sum, at least one $d_{k,\rho}$ is
strictly positive.

Every root in $F^k$ has modulus $R^k$. A product containing a smaller-radius
root has smaller modulus, so it cannot cancel the chosen positive
coefficient. The path identity gives
\[
 M_n^k+\sum_{\beta\in B_k}b_{k,\beta}\beta^n=E_{k,n}-R_{k,n}.
\]
Adding the positive $B_k$ spectrum cannot cancel the chosen top-radius
coefficient either. The left side therefore has a surviving root of modulus
$R^k$. Its other coefficients need not be positive. Applying \fid{F02}
and \eqref{eq:overhead} for fixed $k$ gives
\[
 R^k\le q^{kw/2+\max(d/2+1,\eta_k)}.
\]
Take the $k$th root and the unbounded even-$k$ limit. The all-embedding and
dual arguments are the same as in \fid{F03} and \fid{F17}, but only visible
roots can be used in a reciprocal pairing.

For the finite-extension assertion, group first:
\[
 M_{er}=\sum_\gamma d_{e,\gamma}\gamma^r,
 \qquad d_{e,\gamma}=\sum_{\alpha^e=\gamma}c_\alpha.
\]
If the nonzero grouped spectrum is nonempty, put
$R_e=\max_{d_{e,\gamma}\ne0}|\gamma|^{1/e}$. Applying the preceding
argument in $r$ with $e=e_k$ gives
\[
 R_{e_k}\le q^{w/2+\max(d/2+1,\eta_k)/k}.
\]
Thus $R_{e_k}=R$ along an unbounded allowed set suffices. This equality is
automatic if the top-radius coefficients have one common sign, or if no
ratio of two distinct top-radius roots is a root of unity. It is not
automatic for general signed coefficients. Each embedding needs its own
top-radius visibility condition.
\end{proof}

\paragraph{Two exact limits of the refinement.}
First, take $a=0$, $d=1$, and let $t_n$ be the constant function $0$ for even
$n$ and $2$ for odd $n$. Then $|t_n|\le2$ and
\[
 M_n=q^n-(-q)^n.
\]
Both visible roots have modulus $q>q^{1/2}$. For every even $k$ choose
$e_k=2$. At every required degree $n=2r$ all counts vanish, so the
opposite-sign hypothesis holds with empty $B_k$ and zero error. The top
spectrum disappeared under powering. This refutes the earlier draft
extension without visibility, while leaving the positive-multiset
statement \fid{F17} intact. Geometrically the pointwise sequence itself
is the alternating trace of the constant line minus the unramified line
with Frobenius $-1$; its failure is not an exotic nonalgebraic sequence.

Second, the real-coefficient condition cannot simply be dropped for the
prescribed unbounded even set. Let $a=0$, $d=1$, and $t_n(x)=i$ for every
$x,n$. Then $M_n=iq^n$ has the root $q>q^{1/2}$ with coefficient $i$.
For $k=2m\ge6$ with $k\equiv2\pmod4$, put $a_j=x_{2j-1}$ and $b_j=x_{2j}$. The path
polynomial is
\[
 P_k=\sum_{j=1}^{m}a_jb_j+
       \sum_{j=1}^{m-1}a_{j+1}b_j.
\]
Its coefficient matrix is triangular with diagonal entries one, so it is
equivalent to a nondegenerate bilinear pairing. Its zero count over
$\F_Q$ is $Q^m+(Q^m-1)Q^{m-1}$, including characteristic two.
Since $i^k=-1$,
\[
 QN_{k,n}=-Q^k-Q^{k/2+1}+Q^{k/2}.
\]
Take $B_k=\{q^k,q^{k/2+1}\}$ and $E_{k,n}=Q^{k/2}$; the error has
$\eta_k=0$. Every other stated hypothesis holds for all $n$. This is a
counterexample to that deleted hypothesis, not to the integer-Euler
version or to possible other criteria with phase control.

\paragraph{What the stronger theorem removes.}
For a bounded constructible Weil complex, alternating cohomological
multiplicities are integers. The input need not be supported in one parity
if the even-path output has the prescribed negative visible spectrum and
the pointwise bound holds. Identical roots that already cancel between
degrees remain invisible: for example
$\Q_\ell(-1)\oplus\Q_\ell(-1)[-1]$ has zero trace in every extension.
Consequently \fid{F24} strictly weakens the input parity gate; it does not
infer bounds for erased cohomology or provide the path-output sign.


\subsection{Countable trace completions: an upper bound and a duality obstruction}
The finite-spectrum hypothesis has a useful extension, but its duality
requirement becomes restrictive. Call a visible weighted spectrum
$\mathcal A=\{(\alpha,c_\alpha):\alpha\in A\}$ \emph{admissible} if
$A\subset\C^*$ is bounded, its roots are distinct, every
$\{\alpha:|\alpha|\ge\epsilon\}$ is finite, all $c_\alpha\ne0$, and
for some nonnegative integer $N$,
\[
 S_N=\sum_{\alpha\in A}|c_\alpha|\,|\alpha|^N<\infty.
\]
Equal roots are grouped and zero coefficients removed before imposing
these conditions. The empty spectrum represents zero.

\begin{finding}[\fid{F32}: admissible countable spectra and visible reciprocal duality]\label{f:countable}
For a nonempty admissible spectrum, the maximum radius $R$ is attained
and $u_n=\sum c_\alpha\alpha^n$ converges absolutely for $n\ge N$.
It satisfies
\[
 \limsup_{n\to\infty}|u_n|^{1/n}=R.
\]
An eventual sequence has at most one admissible visible expansion.
Finite sums, tensor products, and integer power substitutions preserve
admissibility after grouping, although power substitution may erase the
entire largest-radius part.

The signed even-path criterion of \fid{F24} remains true when its
input visible spectrum and each positive path-output spectrum are
admissible countable spectra. The identities and bounds need hold only
for all sufficiently large extension degrees, with thresholds depending
on the path length. The input coefficients must still be real, and the
finite-extension and all-embedding versions retain the same explicit
visibility and transformed-hypothesis requirements.

If $A$ is admissible and $\kappa\ne0$, and every $\kappa/\alpha$ is
visible in a bounded dual support, then $A$ is finite. In particular,
an infinite completion of this kind cannot have bounded visible
reciprocal duality. If the nonzero grouped coefficients are integers,
absolute convergence of one nonnegative moment already implies the
boundedness and local-finiteness conditions above.
\end{finding}
\begin{proof}
Choose a root $\alpha_0$. The finite set of roots of modulus at least
$|\alpha_0|$ contains a largest one; it is a global maximum $R>0$.
For $n\ge N$, absolute summation gives $|u_n|\le S_NR^{n-N}$.
The generating function of the shifted sequence is
\[
 H_N(T)=\sum_{m\ge0}u_{N+m}T^m
       =\sum_{\alpha\in A}\frac{c_\alpha\alpha^N}{1-\alpha T}
       \quad (|T|<1/R).
\]
The equality follows from the double absolute bound
$S_N/(1-R|T|)$. Its right side converges locally normally away from
$\{1/\alpha\}$. Indeed, on a compact set with $|T|\le L$, roots
of modulus at most $1/(2\max(1,L))$ have denominator modulus at
least $1/2$, giving a normally convergent tail by $2S_N$; only
finitely many other roots remain. The pole set is locally finite.
At $T=1/\alpha$ the residue is
$-c_\alpha\alpha^{N-1}\ne0$; all other terms are holomorphic in a
neighborhood. Thus the Taylor series has radius exactly $1/R$.
Cauchy--Hadamard, and the harmless fixed shift by $N$, give the
asserted exponential rate. To prove uniqueness, subtract two proposed
expansions at a common sufficiently large moment index. Their grouped
difference is admissible. If nonempty its exponential rate is positive,
contradicting an eventually zero sequence.

For a tensor power of order $k$, a product of modulus at least
$\epsilon$ has every factor of modulus at least
$\epsilon/R^{k-1}$. Only finitely many factor tuples can therefore
produce any nonzero product, or any product outside a prescribed disk.
After grouping to coefficients $d_{k,\rho}$,
\[
 \sum_\rho |d_{k,\rho}|\,|\rho|^N\le S_N^k.
\]
This proves tensor closure, including all multiplicative relations
among roots. Finite sums use a common moment index and the triangle
inequality. For the substitution $n=er$, group the at most $e$ roots
with a given $e$th power and choose $N'$ with $eN'\ge N$; boundedness
and the moment estimate prove admissibility. Cancellation on the top
circle is possible, as $q^n-(-q)^n$ already shows.

For completeness, consider the countable signed path argument directly.
Write $w=a+d$ and let $A_R$ be the finite top circle. For every allowed
even $k$, the finite group-algebra element
$F=\sum_{\alpha\in A_R}c_\alpha[\alpha]$ has
$F^k\ne0$: otherwise evaluation $[\alpha]\mapsto\alpha^n$ for
all positive integers $n$, followed by Vandermonde invertibility,
would force every coefficient to vanish. Its coefficient sum is
$(\sum c_\alpha)^k\ge0$, so at least one of its real grouped
coefficients is positive. The corresponding root has modulus $R^k$.
No product using a smaller-radius root reaches that modulus, and the
positive path-output spectrum cannot cancel its coefficient.

Write the assumed path output as
$q^nN_{k,n}=-\sum_\beta b_{k,\beta}\beta^n+E_{k,n}$, with
$b_{k,\beta}>0$ and
$|E_{k,n}|\le C_kn^{s_k}q^{n(kw/2+\eta_k)}$,
$\eta_k/k\to0$. The exact path identity gives
\[
 M_n^k+\sum_\beta b_{k,\beta}\beta^n=E_{k,n}-R_{k,n},\qquad
 |R_{k,n}|\le B^kq^{n(kw/2+d/2+1)}.
\]
The left side is admissible by closure and visibly contains the root
just found. Applying the proved spectral-radius formula at fixed $k$
gives
\[
 R^k\le q^{kw/2+\max(\eta_k,d/2+1)}.
\]
Taking $k$th roots and then unbounded even $k$ proves
$R\le q^{w/2}$. Constants depending on $k$ and polynomial factors
in $n$ disappeared in the preceding fixed-$k$ limit. If only
$n=e_kr$ are used, the top circle must survive grouping for an unbounded
allowed set of $k$; then the same estimate has the factor $e_k$ on
both sides and yields the same conclusion. For each field embedding,
admissibility, real coefficients, estimates and visibility must be
verified for the transformed system separately.

For the reciprocal assertion let the dual support have radius bound
$R_D<\infty$. It is positive when $A\ne\varnothing$, and
$|\alpha|\ge |\kappa|/R_D$ for every $\alpha\in A$.
Local finiteness away from zero makes $A$ finite. No convergence
assumption on the dual support is needed for this deduction.
Consequently the two independent bounds used in \fid{F17}, together
with visible reciprocity $q^d/A$, still give purity, but now also
force the original visible support to be finite.

Finally suppose the coefficients are nonzero integers and the moment
sum converges. For $N=0$ every term is at least one, so the support is
finite. For $N\ge1$,
\[
 \#\{\alpha:|\alpha|\ge\epsilon\}\le S_N\epsilon^{-N},\qquad
 |\alpha|\le S_N^{1/N}.
\]
These prove the two support conditions. The requirement $N\ge0$
matters if boundedness was not separately given: roots $2^j$ with
coefficients one have a convergent moment at index $-1$.
\end{proof}

\paragraph{The natural completion of the actual normalized trace.}
The trace in \fid{F26}, after $v=q^{\pm n/2}$, is
\[
 h_n=24+\frac{16q^{-2n}}{(1+q^{-n})^4}
 =24+16\sum_{j\ge0}(-1)^j\binom{j+3}{3}
                   (q^{-(j+2)})^n.
\]
The binomial expansion is absolutely convergent for every $n\ge1$.
It gives an admissible integer Euler spectrum
$A=\{1,q^{-2},q^{-3},\ldots\}$, with
$S_1=24+16q^{-2}/(1-q^{-1})^4$. Its radius is one, consistent with
an upper weight-zero bound. This observation does not construct the
pointwise and path-output identities required by the theorem.
For every fixed $\kappa\ne0$, the reciprocal support
$\{\kappa,\kappa q^2,\kappa q^3,\ldots\}$ is unbounded.
With these nonzero integer coefficients its nonnegative moments do
not converge. Uniqueness excludes replacing the same eventual trace
by a different admissible finite or pure reciprocal expansion.
Thus this completion repairs the finite-spectrum requirement for an
upper-bound argument but cannot supply bounded visible reciprocal
duality for this actual word trace. Regularized traces, non-atomic
spectra, duals outside this convergence class, and other comparison
dictionaries are not excluded.

\subsection{What the hypotheses mean geometrically}\label{sec:tracegeom}
Let $\mathcal F$ be a constructible Weil sheaf on $\A^d/\F_q$. Its trace functions $t_n$ have the required pointwise estimate if stalk ranks are bounded by $B$ and all stalk eigenvalues have weight at most $a$. Let
\[
 H_k=\{(x_1,\ldots,x_k):P_k(x_1,\ldots,x_k)=0\}
 \subset(\A^d)^k.
\]
The character kernel is represented by an Artin--Schreier sheaf pulled back along $hP_k$; the coefficient system on $H_k$ is the restriction of $\mathcal F^{\boxt k}$. The Grothendieck trace formula identifies the sums with alternating traces of $R\Gamma_c$.

If $R\Gamma_c(\A^d_{\bar{\F}_q},\mathcal F)$ is concentrated in degree $d$, then $\epsilon=(-1)^d$. If the relevant cohomology of $H_k$ is concentrated in degree $kd-1$, its sign is
\[
 (-1)^{kd-1}=-\bigl((-1)^d\bigr)^k.
\]
Multiplication by $q^n$ is the twist $(-1)$. These facts give precisely (S1)--(S2), or \fid{F04} when background terms have been controlled. A suitable duality of the coefficient system and cohomology supplies \eqref{eq:duality}.

Being a hypersurface of dimension $kd-1$ does \emph{not} prove this concentration. Singularities, exceptional fibers, compact-support boundary groups and Tate contributions are separate obligations. An affine variety also need not identify $H_c^d$ with its dual without an additional comparison. This is the information the proposed sign route still needs. Finding~\ref{f:signed} removes the input common-parity requirement at even lengths, but retains an explicit negative visible path-output spectrum. Its spectral core can be supplied geometrically or by a rigorous algebraic/category substitute; our argument does not prove that schemes are the only possible source of it.

\subsection{Countermodels to omitted hypotheses}
\begin{finding}[\fid{F05}: precise failures of abbreviated sign arguments]\label{f:countermodels}
The following general implications are false.
\begin{enumerate}[label=(\alph*)]
\item Path identity plus Parseval implies the expected middle weight.
\item Bounds for alternating traces control every individual cohomological eigenvalue without non-cancellation.
\item An upper-weight bound implies purity of that weight.
\item Bounds on an arbitrary infinite set of extension degrees suffice.
\item Nonnegative weights and Parseval make the Fourier operator positive.
\item Punctual purity implies arithmetic Frobenius semisimplicity.
\end{enumerate}
\end{finding}
\begin{proof}
(a) Take $t_n=1$ on $\F_{q^n}$, $a=0,d=1$. Then $M_n=q^n$, with eigenvalue $q>\sqrt q$, while \fid{F01} holds for all $k,n$. The actual cohomology is $H_c^2(\A^1)=\Q_\ell(-1)$, not the assumed middle degree.

(b) On $\Spec\F_q$, the complex $\Q_\ell(-1)\oplus\Q_\ell(-1)[-1]$ has a $q$-eigenvalue in degrees 0 and 1. All alternating traces vanish, and its virtual zeta function is 1.

(c) The eigenvalue 1 satisfies the weight-at-most-1 bound but has actual weight zero.

(d) The spectrum $\{R,-R\}$ has trace zero for every odd $n$, with $R$ arbitrarily large. Our lemma uses every extension degree; we assert no general necessity theorem excluding other suitably controlled subsequences.

(e) For $t=1$ on $\F_3$, the normalized Fourier matrix has diagonal entry $e^{2\pi i/3}/\sqrt3$ at $x=1$ and hence is not Hermitian. Over $\F_2$ the normalized matrix is $2^{-1/2}\left(\begin{smallmatrix}1&1\\1&-1\end{smallmatrix}\right)$, with eigenvalues $1,-1$. Both preserve the Euclidean norm.

(f) The matrix $\left(\begin{smallmatrix}1&1\\0&1\end{smallmatrix}\right)$ has pure weight-zero eigenvalues and trace 2 for every power, but is not diagonalizable. It defines a continuous $\ell$-adic representation through $\widehat\Z\to\Z_\ell$ and the associated unipotent matrices. This is the standard arithmetic-versus-geometric distinction also discussed in BBD after 5.3.8.
\end{proof}
In particular Fable suggestion~3's proposed identification of Fourier-after-multiplication with a positive operator is false in that literal generality. This does not refute the valid norm estimate or an independently constructed different positive operator.

\section{From purity to positivity: what is additionally required}\label{sec:positivity}
\begin{finding}[\fid{F06}: weights, positive values, and unimodality are different]\label{f:phase}
Exact purity does not imply a positive Frobenius trace. Even positivity at every positive numerical specialization, together with bar invariance, does not imply coefficientwise positivity. Pure graded pieces with a duality need not have unimodal dimensions.
\end{finding}
\begin{proof}
On $\Spec\F_q$, the rank-one Weil sheaf with Frobenius $-1$ is pure of weight zero but has trace $(-1)^n$ over $\F_{q^n}$. Next
\[
 h(v)=v^2-1+v^{-2}\ge1\quad(v\in\mathbb R\setminus\{0\})
\]
by $v^2+v^{-2}\ge2$, but its central coefficient is negative and $h(v^{-1})=h(v)$. Since $h(1)=1$, a factorization $\pm(v-v^{-1})^a f(v)$ with $a\ge0$ and nonnegative Laurent $f$ would force $a=0$, positive sign and $f=h$, which is impossible. Finally take $M_0=\Q_\ell$, $M_1=0$, $M_2=\Q_\ell(-2)$, with grade $j$ assigned weight $2j$ and the perfect pairing $M_0\otimes M_2\to\Q_\ell(-2)$. Its dimension sequence is $(1,0,1)$, not unimodal. Every degree-one operator $L$ has $L^2=0$, so it cannot satisfy the isomorphisms of \fid{F07} for $D=2$. This fixes the grading convention; it is not the usual degree-two Lefschetz grading on the cohomology of $\PP^1$.
\end{proof}
The last example is an abstract graded Weil object, not the cohomology of a smooth projective variety to which Hard Lefschetz already applies. It isolates precisely the extra theorem that such geometry would supply.

\begin{finding}[\fid{F07}: effective categorification and Lefschetz suffice]\label{f:category}
Let $\mathscr C$ be a split semisimple additive category with grading shift $\langle1\rangle$ and an additive product. Suppose its designated simple representatives satisfy finite decompositions
\[
 L_a*L_b\simeq\bigoplus_{c,j}L_c\langle j\rangle^{\oplus m_{ab}^{c,j}}.
\]
If an isomorphism from its graded split Grothendieck ring to the target algebra sends $[L_a]$ to the \emph{specified} basis element $b_a$ and grading shift to $v$, then its structure constants are
$\sum_jm_{ab}^{c,j}v^j\in\N[v,v^{-1}]$.

For a normalized multiplicity space $M=\bigoplus_{j=0}^D M_j$, suppose there is an operator $L:M_j\to M_{j+1}$ such that
$L^{D-2j}:M_j\to M_{D-j}$ is an isomorphism for $j\le D/2$. Then $\sum_j\dim M_jv^j$ is symmetric and unimodal.
\end{finding}
\begin{proof}
Apply the split Grothendieck relation to the displayed direct sum and compare coefficients under the basis-preserving isomorphism. Multiplicities of actual direct summands are nonnegative integers. For the second assertion the displayed isomorphisms give symmetry. If $j<D/2$ and $Lx=0$, then $L^{D-2j}x=0$, so $x=0$. Thus $\dim M_j\le\dim M_{j+1}$ up to the middle; symmetry gives the other half.
\end{proof}
For VIII this must be applied to the correctly \emph{normalized} multiplicity polynomial, with the signs, $(v-v^{-1})$ factors and positive rational denominators justified separately. The grading step must match the variable in which unimodality is asserted. An arbitrary derived $K_0$ realization is insufficient because $[K[1]]=-[K]$.

\begin{finding}[\fid{F08}: pure geometry and positive Gram determinants do not fix a basis]\label{f:twopoints}
There is a smooth projective finite \(\F_q\)-scheme, an integral basis of effective $K_0$ classes, and a positive definite Gram form, such that multiplication has a negative structure constant in that basis.
\end{finding}
\begin{proof}
Take $X=\Spec\F_q\sqcup\Spec\F_q$. Its $K_0$ ring is $\Z^2$ with componentwise product. The realization $H^0_{\et}(X_{\bar{\F}_q},\Q_\ell)=\Q_\ell^2$ has identity Frobenius and is pure of weight zero. Put
\[
 b_1=(1,1),\qquad b_2=(1,2).
\]
They form an integral basis, since the determinant is 1, and both are actual vector bundles, with the displayed ranks on the two components. The natural real pairing has Gram matrix
\[
 \begin{pmatrix}2&3\\3&5\end{pmatrix},
\]
whose principal minors are $2,5,1>0$. Nevertheless
\[
 b_2^2=(1,4)=3b_2-2b_1.
\]
All assertions follow directly. Exterior powers are also defined componentwise: $\lambda_t(r,s)=((1+t)^r,(1+t)^s)$. Thus passing to the ordinary geometric $\lambda$-ring framework does not remove this counterexample.
\end{proof}
This disproves a general bridge from purity/determinants/effectivity to coefficients of a chosen basis. It does not disprove the positivity of any designated GCT canonical basis.

\section{An obstruction in the actual nonstandard Hecke inclusion}\label{sec:ambient}
\begin{finding}[\fid{F09}: the inherited positive cone has rank one]\label{f:hecke}
Let $R=\Z[v,v^{-1}]$, $R_+=\N[v,v^{-1}]$, $\delta=v+v^{-1}$, and let $H=R1\oplus Rc$ with $c^2=\delta c$. In $H\otimes_RH$, put
\[
 P=c\otimes c+(c-\delta)\otimes(c-\delta),\qquad A=R[P].
\]
Equip the ambient tensor algebra with its positive cone $C_+$ in the basis
$1,c\otimes1,1\otimes c,c\otimes c$. Then
\begin{equation}\label{eq:cone}
 A=R1\oplus RP,\qquad P^2=\delta^2P,\qquad A\cap C_+=R_+1.
\end{equation}
Consequently no basis of the rank-two algebra $A$ consists entirely of ambient-positive elements, although its internal basis $\{1,P\}$ has positive multiplication.
\end{finding}
\begin{proof}
Set $b=c-\delta$. Then $cb=0$ and $b^2=-\delta b$, so the cross terms vanish and
$P^2=\delta^2(c\otimes c+b\otimes b)=\delta^2P$. Also
\[
 P=\delta^2\,1-\delta(c\otimes1+1\otimes c)+2c\otimes c.
\]
This proves that $1,P$ span $A$ and are independent. For $a+bP$, its four ambient coefficients are
\[
 a+\delta^2b,\quad-\delta b,\quad-\delta b,\quad2b.
\]
If all lie in $R_+$, then $2b\in R_+$ implies coefficientwise $b\ge0$. A nonzero such $b$ makes $\delta b$ nonzero and nonnegative, contradicting $-\delta b\ge0$. Thus $b=0$, and then $a\in R_+$. The converse is immediate. Finally $\delta^2=v^2+2+v^{-2}$ is positive, proving the intrinsic positivity assertion.
\end{proof}
This is the actual rank-one nonstandard generator of revised IV, \S1.4, equation~(8), and Definition~11.1, equation~(93). We fix the consistent earlier \S3.1 convention
$(T-v)(T+v^{-1})=0$, $c=C'_s=T+v^{-1}$ and $C_s=T-v=c-\delta$.
The opening of \S11.1 interchanges the exponents while retaining the idempotent assertion; this printed normalization slip is recorded rather than silently used. The proof above is unambiguous and also works over $\Q[v,v^{-1}]$ with the rational positive cone.

\paragraph{Categorical consequence.}
Suppose the usual graded tensor Hecke category has split $K_0$ identified with $H\otimes H$, with actual objects giving the ambient positive cone. Every object in an additive graded subcategory whose character lies in $A$ has class in $A\cap C_+=R_+1$. Its Grothendieck image therefore spans at most $R1$, not $A$. Hence the full nonstandard algebra cannot be obtained as such an effective subcategory through this fixed inclusion.

This is a proved obstruction to a specific natural route. It allows a different category, a different heart or character map, or a derived realization with signs. It is not a counterexample to VIII's signed-normalized conjecture and does not assert that $P$ must be a designated canonical-basis vector. The contrast between intrinsic positivity and failure of inherited positivity is exactly why a new realization is needed for this approach.

\begin{finding}[\fid{F10}: an algebraic replacement and its limits]\label{f:pivots}
Let $R=\Q[v,v^{-1}]$ and $R_+=\Q_{\ge0}[v,v^{-1}]$. Suppose an $R$-algebra $B$ has a positive basis $(e_j)$ and an $R$-subalgebra $A$ has a basis $(b_i)$ with expansions $b_i=\sum_jp_{ji}e_j$, $p_{ji}\in R_+$. If each $b_i$ has an exclusive coordinate $j_i$ with coefficient $r_iv^{m_i}$, $r_i>0$, and all other $b_h$ have zero coefficient there, then $A$ has positive structure constants in this basis. Disjoint positive supports, without this pivot normalization, do not suffice, even with bar invariance and generic semisimplicity. Disjoint supports only after specialization are weaker still.
\end{finding}
\begin{proof}
Write $b_sb_t=\sum_i m_{st}^i b_i$. Its ambient coefficients are positive. The exclusive coordinate is $r_iv^{m_i}m_{st}^i$, so division by the positive monomial preserves the cone and proves $m_{st}^i\in R_+$.

For the second assertion put
\[
 p=v+v^{-1},\quad h=v^2-1+v^{-2},\quad c=ph=v^3+v^{-3},
\]
and take $B=R\times R[t]/(t^2-c)$ with basis
$1=(1,1)$, $x=(0,t)$, $y=(0,1)$. It has positive multiplication
$x^2=cy$, $xy=x$, $y^2=y$. The free subalgebra with basis $1,x,z=py$ has disjoint positive ambient supports but
\[
 x^2=hz,\quad xz=px,\quad z^2=pz.
\]
The coefficient $h$ is negative in its middle term, and \fid{F06} shows it also fails the stated signed-factor positivity form. The product presentation proves associativity, and $t^2-c$ is separable over $\Q(v)$, so the generic algebra is semisimple. Since $p$ is invertible over $\Q(v)$, the generic subalgebra equals the generic ambient algebra and is also semisimple. All the displayed data are bar-invariant.

Finally in the positive algebra $\Q[v][x]/(x^2-x)$, the basis $b_0=1$, $b_1=x+v$ has nonnegative external coefficients and disjoint singleton supports modulo $v$, but direct expansion gives
\[
 b_1^2=(1+2v)b_1-v(1+v)b_0.
\]
This proves the last assertion.
\end{proof}
The statement provides a possible algebraic route that does not require new finite-field geometry. VIII, Conjecture~2.2, asserts disjoint supports in $L_\alpha/vL_\alpha$, not the global exclusive monomial pivots of \fid{F10}. Nor is raw positivity necessarily appropriate for all of VIII's signed targets. The missing pivot/normalization theorem cannot be obtained merely by repeating the local support assertion.

\begin{finding}[\fid{F11}: two limits of the Grothendieck-ring shortcut]\label{f:lambda}
Neither virtual $\lambda$-ring identities nor reduction of a noncommutative algebra modulo $p$ supplies the missing effective positive cone and arithmetic Frobenius.
\end{finding}
\begin{proof}
In $K_0(\Rep\mathbb G_m)=\Z[z,z^{-1}]$, the line character $z$ satisfies
\[
 \lambda_t(z-1)=\frac{1+zt}{1+t}
 =1+(z-1)t+(1-z)t^2+\cdots.
\]
Thus $\lambda^2(z-1)=1-z$, outside the effective cone $\N[z,z^{-1}]$. Also
$\gamma_t(z-1)=\lambda_{t/(1-t)}(z-1)=1+(z-1)t$. The augmentation/\(\gamma\)-filtration bookkeeping supplies no additional positivity order.

For the Frobenius assertion take $A=e_{12}$ and $B=e_{21}$ in $M_2(\F_p)$. Their $p$th powers vanish, whereas $(A+B)^2=I$, hence $(A+B)^p=I$ for $p=2$ and $A+B$ for odd $p$. The map $x\mapsto x^p$ is not additive. Therefore the automatic absolute Frobenius of a commutative scheme is not obtained by the same formula on a noncommutative coordinate algebra.
\end{proof}
A separately constructed quantum or categorical Frobenius remains possible. Also braided tensor categories do not automatically carry the symmetric exterior-power structure required by SGA~6.

\section{The theta route: exact positive consequences}\label{sec:theta}
\subsection{Periods and multiplicities are different data}
Alp\"oge's step~4 (p.~3) asserts positivity of a functional via a determinant. Step~7 (p.~4) describes determinants containing GGP periods, interlacing values 0 or 1, and controlled failures of interlacing. Step~8 describes iterated lifts between totally definite unitary groups and a gain of multiplicity. Page~1 and the introductory paragraph on p.~3 explain this as the analogue of Stepanov order of vanishing, not an assertion about the dimension of a Kronecker Hom space. The supplied Lean wrappers contain no theorem identifying the two kinds of multiplicity.

The following two findings state exactly what a positive projector or a chain of branching rules can give. They also show why merely finding a positive determinant is not the missing representation-theoretic formula.

\begin{finding}[\fid{F12}: positive projectors already realize numerical multiplicities]\label{f:projector}
For a finite group $H$ and a finite-dimensional complex unitary representation $U$, put
\[
 P_H=\frac1{|H|}\sum_{h\in H}\rho(h).
\]
It is the orthogonal projector onto $U^H$, with $\Tr P_H=\rank P_H=\dim U^H$ and nonnegative principal minors in an orthonormal basis. For partitions of $D$,
\[
 g(\lambda,\mu,\nu)=\Tr P_{S_D}
 \quad\text{on }[\lambda]\otimes[\mu]\otimes[\nu].
\]
For $\lambda\vdash dn$ and $\dim V\ge\ell(\lambda)$,
\[
 a_\lambda(d[n])=\Tr P_{H_{d,n}}\quad\text{on }[\lambda],\qquad
 H_{d,n}=(S_n)^d\rtimes S_d\subset S_{dn},
\]
where $a_\lambda(d[n])$ is the multiplicity in $\Sym^d(\Sym^nV)$.
\end{finding}
\begin{proof}
Inversion gives $P_H^*=P_H$, and counting each element $|H|$ times in the double sum gives $P_H^2=P_H$. Its image is fixed by $H$, and it is the identity on invariant vectors. Thus it is the claimed orthogonal projector. In an orthonormal basis its principal submatrices are Gram matrices of the vectors $P_He_i$, so their determinants are squared exterior norms. Its spectrum consists of $\dim U^H$ ones and the remaining zeros.

Specht modules are self-dual and can be made unitary by averaging a form; the triple invariant-space description gives the Kronecker assertion. For plethysm, invariants under the within-block permutations give $(\Sym^n V)^{\otimes d}$, and under block permutations give $\Sym^d(\Sym^nV)$. Apply invariants to Schur--Weyl decomposition
\[
 V^{\otimes dn}=\bigoplus_\lambda S_\lambda V\otimes[\lambda]
\]
and compare the $S_\lambda V$ multiplicity. This is also the projector explicitly used in BIP, equation~(4.1), with Proposition~3.3.
\end{proof}
This is a proof of the useful reformulation, not a new numerical nonnegativity result. BIP, Proposition~4.5, already gives projected tableaux spanning the highest-weight space; its Theorem~4.7 expresses evaluations as sums of signed column contributions. Independence and a compatible basis remain separate.

For any spanning vectors $v_i$ of a multiplicity space, their Gram matrix satisfies
\[
 \sum_{|I|=r}\det G_I=\sum_{|I|=r}\bigl\|\bigwedge_{i\in I}v_i\bigr\|^2.
\]
This sum is positive exactly when the rank is at least $r$. Determinant positivity therefore restates the rank question; it does not by itself determine which minors are nonzero or provide an efficient basis-selection rule. We make no claim that every such rank question is intrinsically hard.

\begin{finding}[\fid{F13}: when an interlacing chain counts a target multiplicity]\label{f:branching}
Let $\mathscr C_r,\ldots,\mathscr C_0$ be split semisimple linear categories of finite-length objects and let $F_i:\mathscr C_i\to\mathscr C_{i-1}$ be additive exact functors. Suppose
\[
 F_i S_i(a)\simeq\bigoplus_b S_{i-1}(b)^{\oplus b_i(a,b)},\qquad b_i(a,b)\in\{0,1\},
\]
with finite sums. Then the multiplicity of $S_0(c)$ in $F_1\cdots F_rS_r(a)$ is the number of allowed paths from $a$ to $c$:
\[
 \sum_{a_{r-1},\ldots,a_1}
 b_r(a,a_{r-1})\cdots b_1(a_1,c).
\]
If the allowed steps are specified by integral interlacing inequalities and every intermediate coordinate is bounded by the fixed endpoints, these paths are the integer points of the resulting polytope.
\end{finding}
\begin{proof}
Apply each additive functor to the previous finite direct-sum decomposition. The multiplicities multiply along paths and add over intermediate simples. Semisimplicity identifies the final coefficient with the desired multiplicity. Every product of 0/1 entries is 1 exactly for an allowed path. Sending a path to its array of intermediate integer weights is a bijection with the integer solutions of the inequalities. Boundedness makes the polyhedron a polytope.
\end{proof}
Ordinary restriction along consecutive block inclusions $\GL_r\supset\GL_{r-1}$ has the interlacing rule $\lambda_j\ge\mu_j\ge\lambda_{j+1}$; see \cite[Theorem 2.1, equation (2.2), and the patterns (2.4)]{Molev}. The Kronecker map $(g,h)\mapsto g\otimes h$ and the plethysm map $g\mapsto S_\mu(g)$ are different maps. A theta lift between dual pairs is not automatically their composite restriction functor.

A small object mismatch is instructive. $S_{(d^r)}(\C^r)=\det^d$ is one-dimensional and has a unique usual interlacing pattern, whereas the rectangular Specht module is a different representation. For example
\[
 [(2,2)]\otimes[(2,2)]=[(4)]\oplus[(2,2)]\oplus[(1^4)].
\]
To check this without a table lookup, the permutation action of $S_4$ on the three pairings of four letters is the direct sum of the trivial representation and $[(2,2)]$. The latter has character $(2,0,2,-1,0)$ on cycle types $(1^4),(2,1,1),(2,2),(3,1),(4)$. Its pointwise square equals this character plus the trivial and sign characters, proving the identity. A correctly identified chain could still count higher multiplicities; the example only prevents substitution of a determinant character for a rectangular Specht module.

\subsection{A lattice obstruction to ordinary interlacing alone}
\begin{finding}[\fid{F18}: homogeneous difference constraints and a rectangular parity ray]\label{f:saturation}
Let $u$ range over an integral label set closed under positive-integer scaling. Fix a number of coordinates and a finite list of inequalities
\[
 P(u)=\{x\in\mathbb R^s:x_i-x_j\le b_{ij}(u)\},\qquad x_0=0,
\]
where every specified $b_{ij}(u)$ is integral and $b_{ij}(Nu)=Nb_{ij}(u)$ for every $N\ge1$. Assume $P(u)$ is bounded, and put $h(u)=\#(P(u)\cap\Z^s)$. Then
\[
 h(Nu)>0\quad\Longleftrightarrow\quad h(u)>0.
\]
The same property holds for a fixed finite sum of such counting functions. Thus these models cannot count all two-row rectangular Kronecker coefficients. On the particular ray with all three partitions $(d,d)$, the established two-row formula gives
\[
 g((d,d),(d,d),(d,d))=\boldsymbol1_{2\mid d}
 =\#\bigl(d\{1/2\}\cap\Z\bigr)\qquad(d\ge1).
\]
This ray has an exact rational Ehrhart model despite its failure of the displayed ordinary support property.
\end{finding}
\begin{proof}
Homogeneity gives $P(Nu)=NP(u)$. If $x_i-x_j\le b$ with $b\in\Z$, then
$\lfloor x_i\rfloor\le\lfloor x_j+b\rfloor=\lfloor x_j\rfloor+b$.
Hence coordinatewise flooring, including $x_0=0$, takes every real point of $P(u)$ to an integral one. If $z\in P(Nu)\cap\Z^s$, apply this to $z/N$. Conversely an integral $x\in P(u)$ gives $Nx\in P(Nu)\cap\Z^s$. Positivity of a fixed finite sum is positivity of at least one summand, so the same argument applies. Boundedness is only needed to make these finite counts.

For $u=((1,1),(1,1),(1,1))$, the triple tensor product is the sign representation of $S_2$, so $g(u)=0$. The explicit $S_4$ character calculation following \fid{F13} gives $g(2u)=1$: the class sizes are $(1,6,3,8,6)$ and the character $(2,0,2,-1,0)$ has cubic average $(8+3\cdot8-8)/24=1$. Thus an exact model of the stated form on all two-row rectangles is impossible.

For the whole ray use IV, \S18.3, equation~(172), printed p.~128. In its column convention, $(d,d)=[d,0]$ as a two-row label and $[0,0,d,0]$ as a four-row label. The coefficient for the second rectangle is the constant term of that equation's generating function. Only $k=0$ and $j=0$ in its first sum survive; the second sum is empty. The remaining condition is that $n_3,n_2,n_1$ are all even and nonnegative. Substituting $(n_3,n_2,n_1)=(0,d,0)$ gives exactly the displayed parity formula. The single rational point $\{1/2\}$ has an integral point after dilation by $d$ exactly when $d$ is even.
\end{proof}
The difference-constraint result includes ordinary interlacing with integral endpoint rows. Its exclusions are essential: extra congruences, sum equations, nonhomogeneous endpoint shifts, a changing constraint list or a changed lattice can evade it. For example $2z=d$ retains precisely the parity that ordinary integral differences lose. The example concerns all two-row rectangles, not a model restricted only to IP's large-parameter region.

Nor does it contradict VI's definition of ``strictly saturated'' (Definition~1.2.4, printed p.~14). That definition tests each nonzero residue polynomial separately; here $f_{\rm even}=1$ and $f_{\rm odd}=0$, so its saturation index is zero. We use the displayed support property to avoid conflating the two notions. IV, Theorem~18.2 and equations~(163)--(164), explicitly retain congruences and remove columns in pairs. A successful theta comparison must preserve these arithmetic data or equivalent information; the source does not identify them with a particular GGP local sign.

In Grothendieck's language, the tensor group-scheme map $\GL_{2,\Z}\times\GL_{2,\Z}\to\GL_{4,\Z}$ induces, over $\Q$, a positive restriction map on representation $K_0$. The coefficient of $[\det^d\boxtimes\det^d]$ in the restriction of $[S_{(d,d)}(V\otimes W)]$ is this parity function. Thus an effective $K_0$ homomorphism does not by itself give the ordinary support property. The proof above is already the exact integral-lattice dictionary; its group-scheme framework is SGA~3 I, 4.7.2--3, with cote~90, scan~17, as an action/comodule anchor.

\paragraph{Exact remaining theta gate.}
Construct functors and endpoint identifications making Alp\"oge's lift chain the target restriction, or identify its period spaces with the projector images of \fid{F12}. To obtain canonical-basis positivity, additionally identify grading, operations, and the specified basis as in \fid{F07}. For a \(\#P\) rule the path encodings and admissibility tests must have polynomial bounds; for PH1 the stronger polyhedral specifications in VI must also hold. None of these identifications is supplied by steps~4, 7 or~8.

% Add bibliography key BlasiakRep for Jonah Blasiak,
% Representation theory of the nonstandard Hecke algebra,
% arXiv:1201.2209v2 (2012), https://arxiv.org/abs/1201.2209v2.
\subsection{An actual nonstandard tower: specialization fails on lines and succeeds after mixing}
There is already a nonstandard tower to which the generic path mechanism
applies. Blasiak \cite[Theorem~5.1, p.~14; Proposition~5.2, pp.~14--15;
Definition~6.1, p.~23]{BlasiakRep} proves that the two-row nonstandard
Temperley--Lieb algebra over $K=\Q(v)$ is split semisimple, that its
successive restrictions are multiplicity free, and hence that its simple
modules have seminormal path lines. These are established source results,
not new consequences of purity. The source distinguishes these lines
from a basis compatible with classical decomposition at $v=1$; its
Remark~6.2 also gives a non-multiplicity-free example outside the
two-row quotient. In particular, its simple module at \emph{tower rank
four}
\[
 M_K=M_{(3,1),K}\otimes M_{(2,2),K}
\]
has dimension six and restriction
\begin{equation}\label{eq:towerrestriction}
 M_K|_{A_{3,K}}=N(\delta)\oplus N(D-2)\oplus\epsilon_+
                  \oplus\epsilon_-.
\end{equation}
Here the rank-three notation is that of \fid{F19}/\fid{F21}.
The following comparison computes exactly what happens to its path lines
in the natural ordinary Specht tensor lattice.

\begin{finding}[\fid{F30}: regular path mixing and a projector base-change defect]\label{f:tower}
Let $R=\Q[v,v^{-1}]_{(v-1)}$, $k=R/(v-1)=\Q$, and fix
\[
 M_R=\mathsf S_{(3,1),R}\otimes_R\mathsf S_{(2,2),R},
\]
with the specified nonstandard action. Write $A_R$ for the actual
nonstandard algebra acting at tower rank four. Its specialization has
\[
 M_1\simeq[(3,1)]\oplus[(2,1,1)].
\]
The two generic path lines through the penultimate summand $N(\delta)$
lie in a regular saturated rank-two plane $E_R$. No vector on either
line which is diagonally rescaled to be regular with nonzero reduction
can specialize into one of these two classical constituents.

Nevertheless a regular non-diagonal change of the six path representatives
in this same lattice does produce a specialization-compatible basis.
It mixes only two pairs of corresponding rank-two path labels, by
unit rescaling followed by the blocks
\[
 \begin{pmatrix}1&1\\1&-1\end{pmatrix}.
\]
Finally
\[
 \End_{A_R}(M_R)=R,\qquad \End_{A_1}(M_1)=k\times k,
\]
and the natural base-change map is $a\mapsto(a,a)$. Thus the two
classical projectors do not lift to endomorphisms commuting with the
whole nonstandard deformation, although the stated regular mixed basis
exists.
\end{finding}
\begin{proof}
\par\smallskip\noindent\textbf{A fixed saturated plane.}\quad
The ordinary parabolic Hecke algebra $H_3\subset H_4$ has trivial
projector
\[
 p_+=\frac{\sum_{w\in S_3}v^{\ell(w)}T_w}
              {\sum_{w\in S_3}v^{2\ell(w)}}.
\]
The denominator specializes to six and is a unit of $R$. The numerator
is multiplied by $v$ by each simple generator on either side; squaring
it therefore multiplies it by the displayed denominator. Hence $p_+$
is a regular central idempotent of $H_3$. Ordinary branching gives
\[
 M_{(3,1)}|_{H_3}=M_{(3)}\oplus M_{(2,1)},\qquad
 M_{(2,2)}|_{H_3}=M_{(2,1)}.
\]
It follows that
\[
 E_R=M_R(p_+\otimes1)
\]
is a free rank-two direct summand, stable under $A_{3,R}$, and its
generic fiber is exactly the $N(\delta)$ summand in
\eqref{eq:towerrestriction}. Because the projector is regular,
\begin{equation}\label{eq:towersaturated}
 E_R=E_K\cap M_R.
\end{equation}
Indeed an element of the intersection is fixed by $p_+\otimes1$
and so lies in its $R$-image. Every regular nonzero reduction of a
rescaled generic path vector in $E_K$ therefore lies in $E_1$.

\par\smallskip\noindent\textbf{The special decomposition and its class sum.}\quad
At $v=1$, $P_i=2(1+s_i)$ for the diagonal symmetric-group action,
so the specialized algebra image generates that action. On the five
cycle types $(1^4),(2,1,1),(2,2),(3,1),(4)$, the characters are
\[
 \chi_{31}=(3,1,-1,0,-1),\qquad
 \chi_{22}=(2,0,2,-1,0),\qquad
 \chi_{211}=(3,-1,-1,0,1).
\]
Their product satisfies
\[
 \chi_{31}\chi_{22}=(6,0,-2,0,0)=\chi_{31}+\chi_{211},
\]
which proves the stated classical decomposition in characteristic zero.
Let $U$ be the sum-zero subspace of $\mathbb R^4$, and $W$ the
sum-zero subspace of the permutation representation on the three
pairings of four letters. These are orthogonal models for $[31]$
and $[22]$. The $S_3$ fixing the fourth letter has invariant line
in $U$ generated by
\[
 w=(1,1,1,-3),\qquad \langle w,w\rangle=12.
\]
Thus $E_1=\mathbb Rw\otimes W$. Let $p$ project orthogonally onto
$\mathbb Rw$, put $\Pi=p\otimes1$, and let
\[
 Z=\sum_{1\le i<j\le4}(ij)
\]
act diagonally. For $i,j\le3$, the transposition fixes $w$ and
$p\rho_U((ij))p=p$. For an external transposition $(i4)$ its inner
product with $w$ is $-4$, giving
\[
 p\rho_U((i4))p=-\tfrac13p.
\]
The action on the three pairings is the quotient $S_4\to S_3$ with
Klein-four kernel. The three internal transpositions map bijectively
to its three transpositions, and so do the three external ones.
The sum of those three matrices in the natural three-point
permutation module is the all-ones matrix; it vanishes on $W$.
Consequently both transposition sums on $W$ are zero and
\begin{equation}\label{eq:towercompression}
 \Pi Z\Pi=p\otimes(0-\tfrac13\,0)=0.
\end{equation}

The class sum $Z$ acts on $[31]$ and $[211]$ by $+2$ and $-2$:
its scalar is six times the transposition character divided by the
dimension, namely $6(\pm1)/3$. If a nonzero $x\in E_1$ belonged
to either constituent, then $\Pi Z\Pi x=\pm2x$, contradicting
\eqref{eq:towercompression}. Thus $E_1$ meets each constituent only
in zero. Together with \eqref{eq:towersaturated}, this proves the
diagonal-normalization obstruction. All identities in this real
model have rational coefficients, so the conclusion holds over $k$.

\par\smallskip\noindent\textbf{Regular complementary path planes.}\quad
Put $J_R=M_R((1-p_+)\otimes1)$. It is the complementary regular
rank-four summand. The three central block idempotents
$z_+,z_-,z_{\rm mat}$ of the actual rank-three algebra in \fid{F19}
are defined over $R$, without dividing by the colliding-root
difference. They commute with the ordinary parabolic projector.
Set
\[
 L_{+,R}=J_Rz_+,\quad L_{-,R}=J_Rz_-,\quad F_R=J_Rz_{\rm mat}.
\]
These direct summands are free over the local ring $R$, with generic
types $\epsilon_+,\epsilon_-,N(D-2)$ and ranks $1,1,2$. Hence
\[
 M_R=L_{+,R}\oplus L_{-,R}\oplus E_R\oplus F_R.
\]
This separates $E_R$ from $J_R$ by the ordinary projector first;
it does not try to separate the two colliding matrix branches by a
singular nonstandard central projector.

At specialization the four $S_3$ types are respectively trivial,
sign, standard, standard. The first-factor orthogonal complement
$U_0=w^\perp\subset U$ is standard, so
$J_1=U_0\otimes W=1\oplus\mathrm{sgn}\oplus[21]$ and $F_1$
is its unique standard summand. Write these specialized spaces as
$L_+,L_-,E,F$. Since $Z$ commutes with $S_3$, it sends $E$ into
$E\oplus F$; there is no homomorphism from the standard module
to either scalar type. Equation~\eqref{eq:towercompression} removes
the $E$ component. Also $Z^2=4$ by the classical decomposition.
It follows that
\[
 Z(E)=F,\qquad Z(F)=E,\qquad J:=Z/2:E\longrightarrow F
\]
is an isomorphism with inverse $Z/2$. Its graph planes
\[
 G_+=\{x+Jx:x\in E\},\qquad G_-=\{x-Jx:x\in E\}
\]
have $Z$-eigenvalues $+2$ and $-2$. The classical restrictions
$[31]|_{S_3}=1\oplus[21]$ and
$[211]|_{S_3}=\mathrm{sgn}\oplus[21]$ therefore identify the
actual constituents as
\begin{equation}\label{eq:towergraphs}
 [31]=L_+\oplus G_+,\qquad [211]=L_-\oplus G_-.
\end{equation}

\par\smallskip\noindent\textbf{A two-pair mixing repair.}\quad
The idempotent $e=P_1/D$ from tower rank two is regular. Each of $E_R,F_R$
splits by $e,1-e$ into two rank-one free direct summands. Choose
bases $x_{+,R},x_{-,R}$ of these two lines in $E_R$, and
$f_{+,R},f_{-,R}$ in $F_R$. Their generic lines are exactly the
four tower paths through the two standard penultimate labels.
At specialization $e=(1+s_1)/2$. Since $J$ commutes with $s_1$,
it carries each line of $E$ to the matching line of $F$. Rescale
each $f_{\sigma,R}$ by a constant unit of $\Q$ so that its
reduction is $Jx_\sigma$.
For each $\sigma=+,-$, replace the two regular vectors by
\[
 g_{\sigma,+}=x_{\sigma,R}+f_{\sigma,R},\qquad
 g_{\sigma,-}=x_{\sigma,R}-f_{\sigma,R}.
\]
The displayed two-by-two block has determinant $-2$, a unit in
$R$, so these four vectors, together with bases of $L_{+,R}$
and $L_{-,R}$, form an $R$-basis. Their reductions lie in the
two graph planes in \eqref{eq:towergraphs}. This proves the
regular specialization-compatible basis claim in the original
fixed lattice. It uses non-diagonal mixing across distinct generic
penultimate summands, precisely the freedom left open by the first
part of the finding.

\par\smallskip\noindent\textbf{Why the classical projectors still do not lift.}\quad
By the cited split simplicity,
$\End_{A_K}(M_K)=K\,\id$. Any element of $\End_{A_R}(M_R)$
therefore becomes a scalar $\lambda\id$ generically. In an
$R$-basis its diagonal entries belong to $R$, forcing $\lambda\in R$.
Conversely all such scalars commute with the action. This proves
$\End_{A_R}(M_R)=R$. At specialization the image is the diagonal
$k[S_4]$ action with two pairwise nonisomorphic absolutely simple
constituents, each once, so its endomorphism ring is $k\times k$.
The base-change map sends a scalar to $(a,a)$; it is injective with
one-dimensional cokernel, and neither nontrivial classical projector
lies in its image.

The three-vector groups in the mixed basis span direct summands of
the underlying free $R$-module, whose reductions are the two
classical constituents. They are not $A_R$-stable generically,
which would contradict generic simplicity. Thus the constructive
basis comparison makes no false claim of a nonstandard-equivariant
projector lift.
\end{proof}

\paragraph{Exact scope and the coherent dictionary.}
The excluded operation is diagonal normalization of this tower's
path lines in the specified ordinary Specht tensor lattice; the
positive operation is the explicit regular mixing just constructed.
No arbitrary-lattice obstruction, VIII canonical-basis identification,
uniform general specialization algorithm or positive structure-constant
rule is inferred. The known path labels count generic simple dimensions,
not classical Kronecker multiplicities by themselves. Their verification
is polynomial in the expanded tower rank, not automatically in its
binary logarithm.


\section{Two further geometric routes and their limits}\label{sec:other}
\begin{finding}[\fid{F14}: what an ambient pure realization would prove]\label{f:sub}
A Frobenius-stable subquotient of a pure algebraic Weil module of weight $w$ is pure of weight $w$. A stable subquotient of its endomorphism module, with Frobenius acting by conjugation, is pure of weight zero.
\end{finding}
\begin{proof}
A basis adapted to a stable subspace makes Frobenius block upper triangular. The characteristic polynomial factors into those of the subspace and quotient, so their eigenvalues are selected from the original spectrum. This works under every embedding. Eigenvalues on $V\otimes V^*$ are ratios $\alpha/\beta$, of absolute value 1; apply the first assertion. No diagonalizability is required.
\end{proof}
Thus if the desired nonstandard multiplicity objects were already realized as such stable subquotients, a new spectral RH would not be needed for them. Stability, grading, effective classes, and a basis-preserving realization would still need proof. Declaring an arbitrary vector space to have identity Frobenius proves an artificial weight-zero assertion and identifies no GCT coefficient.

\begin{finding}[\fid{F15}: orbit models exist, but their ordinary purity does not determine Hilbert multiplicities]\label{f:orbit}
An integral polynomial tensor $f$ in a finite free representation $W_\Z$ of $\GL_{N,\Z}$ has a finite-type affine orbit-closure model over $\Z$ whose complex fiber is the usual reduced orbit closure. Ordinary pure cohomology of a projective scheme does not determine its polarized coherent Hilbert function or its equivariant coherent character.
\end{finding}
\begin{proof}
The orbit map induces $\Z[W]\to\Z[\GL_N]$. Its kernel $I$ defines the finite-type scheme $\Spec(\Z[W]/I)$. The quotient is a subring of $\Z[\GL_N]$ and is torsion-free, hence flat over $\Z$. Since $\C$ is flat over $\Z$, tensoring preserves the kernel. The resulting image ring is a subring of the reduced domain $\C[\GL_N]$, so the complex fiber is the reduced scheme-theoretic orbit closure. No corresponding equality after arbitrary reduction modulo $p$ is asserted; kernel formation can fail to commute with nonflat base change. For an explicit example, let $\mathbb G_m$ scale $f=p$ in $\A^1_\Z$. The map $\Z[x]\to\Z[t,t^{-1}]$, $x\mapsto pt$, is injective, so this model is $\A^1_\Z$. Reducing $f$ first gives zero and an orbit closure consisting of one point, whereas the special fiber of the model is $\A^1_{\F_p}$.

For the second assertion fix $X=\PP^1_{\F_q}$ and $L_r=\mathcal O(r)$, $r>0$. Monomials of degree $rd$ in two variables form a basis of $H^0(X,L_r^d)$, so its dimension is $rd+1$. The underlying scheme and its ordinary/intersection cohomology are independent of $r$: they are $\Q_\ell$ in degree 0 and $\Q_\ell(-1)$ in degree 2 (the projective-space calculation, \cite[Example 16.3, p.~109]{Milne}). Thus the pure Frobenius data do not determine the polarized Hilbert function. Finally, on a point with trivial $\mathbb G_m$ action, a weight-$r$ linearization of $\mathcal O$ gives character $z^{rd}$ on its $d$th tensor power, while the underlying Weil cohomology remains the same. This proves the character assertion.
\end{proof}
A projective integral model of the homogeneous orbit closure can likewise be formed by projectivization. For a reduced equidimensional projective fiber $X/\F_q$ of dimension $e$ with a dense smooth open $j:U\hookrightarrow X$, the unshifted complex $\IC_X^u=j_{!*}\Q_\ell[e][-e]$ is pure of weight zero (BBD~5.3.2), and $IH^i(X)=H^i(X_{\bar{\F}_q},\IC_X^u)$ is pure of weight $i$ by proper pushforward and duality (BBD, Remark~5.4.9, printed p.~143). For a nonreduced fiber pass to its reduction for this \(\ell\)-adic assertion. This is already ordinary finite-field mathematics. The GCT multiplicities instead belong to the graded coherent representation $\C[X]_d$, or to coherent sections with the required saturation and degree hypotheses. The polarization, linearization and comparison object cannot be omitted. Riemann--Roch is consistent with the example: it explicitly retains the varying Chern character of $L_r^d$.

Hilbert/Quot schemes represent the relevant coherent moduli under their usual hypotheses (FGA, Expos\'e~221, Theorems~3.1--3.2); they do not identify their intersection cohomology with the desired multiplicity spaces. Constructing an affine space from an already known multiplicity module encodes its unknown dimension in a scheme, but does not compute that dimension or produce a canonical indexing set. This route was pursued and reduces to an additional coherent-to-constructible comparison, not to an unproved ordinary Weil theorem.

\begin{finding}[\fid{F16}: a positive rank-one model exists with a different action]\label{f:rankonemodel}
The abstract rank-one algebra of \fid{F09} has a positive realization inside the same ambient algebra after changing the inclusion. Put $E=c\otimes c$. Then $R1\oplus RE$ is a subalgebra, and $P\mapsto E$ gives an $R$-algebra isomorphism from $R1\oplus RP$. The coefficient $\delta^2$ is the normalized Tate cohomology polynomial of $\PP^1\times\PP^1$. This isomorphism does not preserve the original ambient action relevant to GCT.
\end{finding}
\begin{proof}
We have $E^2=\delta^2E$, and $1,E$ are independent in the ambient tensor basis. Hence both algebras have the presentation $R[z]/(z^2-\delta^2z)$, with the respective chosen generators. The new basis $1,E$ is ambient-positive. The even cohomology of $\PP^1\times\PP^1$ is $\Q_\ell$, $\Q_\ell(-1)^{\oplus2}$ and $\Q_\ell(-2)$ in degrees 0, 2 and 4, by the projective-space calculation and K\"unneth. Its polynomial, normalized by $v^{-2}$ with $v^2$ counting a Tate degree, is
\[
 v^{-2}(1+2v^2+v^4)=v^{-2}+2+v^2=\delta^2.
\]
At $v=1$, let $s$ be the nontrivial element of $S_2$, so $c=1+s$. Then
\[
 P=2(1+s\otimes s),\qquad E=(1+s)\otimes(1+s).
\]
On the sign tensor sign representation, $P$ acts by 4 and $E$ by 0. Therefore the isomorphism changes the given ambient representation and cannot be substituted into the Kronecker comparison without a new compatibility theorem.
\end{proof}
This is a positive result for the abstract rank-one Hecke case and a precise limit of that result. Purity here is degree by degree: the total ungraded cohomology has several weights, not one weight-zero Weil module. The Tate-grade dimensions are $(1,2,1)$, whereas consecutive integer powers of $v$ give $(1,0,2,0,1)$; no assertion about VIII's precise unimodality normalization follows without fixing the grading step. This is neither an identification of VIII's full specified bases nor a new consequence of the sign-only theorem. In motive language the coefficient is realized by the Tate object $\mathbf1\oplus\mathbb L^{\oplus2}\oplus\mathbb L^2$, with a grading shift; the missing datum is the action comparison, not its purity.

\subsection{Derived lifts and a positive model of the actual rank-one action}
\begin{finding}[\fid{F27}: what changing parity or the effective cone accomplishes]\label{f:derived}
Let $R=\Z[v,v^{-1}]$, $\delta=v+v^{-1}$,
$H=R[c]/(c^2-\delta c)$, and let $P$ be the actual rank-one element
of \fid{F09}. In a graded additive monoidal categorification whose
split Grothendieck ring has the named effective basis
$1,c_1,c_2,c_1c_2$, the following hold.
\begin{enumerate}
\item $P$ has an explicit bounded two-term lift, but no bounded complex
with Euler class $P^k$ ($k\ge1$) has all its nonzero terms in a single
homological parity. An overall sign or internal grading shift does not
alter this obstruction.
\item The explicit zero-differential lift has tensor-power homotopy
amplitude $k$. Separately, if a nonzero bounded complex in an abelian tensor category
with faithful exact vector-space fiber functor has extreme nonzero
cohomology degrees $p,r$, its $k$th tensor power has amplitude $k(r-p)$.
Mixed cohomological parity persists in every tensor power.
\item A different positive cone does preserve the actual rank-one
character actions. The faithful map
\[
 H\longrightarrow R^2,\qquad a+bc\longmapsto(a+b\delta,a)
\]
becomes an isomorphism after inverting $\delta$. Its tensor square sends
\[
 P\longmapsto(\delta^2,0,0,\delta^2)
\]
on the four labels $++,+-,-+,--$. Thus the missing sign-tensor-sign
action in \fid{F16} is restored. In the even form $u=v^2$, the two
nonzero values are actual positive Tate classes
$u^{-1}+2+u$.
\end{enumerate}
These statements do not identify the new character cone with VIII's
canonical basis or construct the negative path-output spectrum of
\fid{F24}.
\end{finding}
\begin{proof}
The already proved quadratic relation gives
\[
 P^k=\delta^{2k-2}P
 =\delta^{2k}\,1-\delta^{2k-1}(c_1+c_2)
                  +2\delta^{2k-2}c_1c_2.
\]
Its four named coordinates have both signs for every $k$.
Let $B_0,B_1,B_2,B_3$ represent $1,c_1,c_2,c_1c_2$, with $B_0$ the
tensor unit. Internal shift $(1)$ represents $v$ and is distinct
from homological shift $[1]$. Set
\[
 X_+=B_3^{\oplus2}\oplus B_0(2)\oplus B_0^{\oplus2}\oplus B_0(-2),
\]
\[
 X_-=B_1(1)\oplus B_1(-1)\oplus B_2(1)\oplus B_2(-1).
\]
The complex $K^0=X_+$, $K^1=X_-$ with zero differential has class
$P$. Any complex with terms only in even degrees has effective Euler
class, and one with terms only in odd degrees has the negative of an
effective class. Neither can equal the displayed mixed-sign class of
$P^k$. Homotopy equivalence preserves Euler class, so it cannot repair
this; overall signs and grading monomials preserve the mixed pattern.

The zero-differential $K^{\otimes k}$ has nonzero terms in all degrees
$0,\ldots,k$. Use unit summands from $X_+$ and a $B_1$ summand from
each chosen $X_-$; $c_1^j=\delta^{j-1}c_1\ne0$ proves nonvanishing
when $j\ge1$. If a zero-differential complex is homotopy equivalent
to $D$, the composite maps back to it equal the identity degreewise,
because its homotopy correction $dh+hd$ is zero. Hence its degree-$i$
term is a retract of $D^i$; no nonzero extreme can disappear. This
proves the amplitude claim for this particular lift, not for every
complex with its Euler class.

For the separate cohomological assertion, apply the faithful exact
tensor fiber functor. Over a field the K\"unneth decomposition of
$H^j(K^{\otimes k})$ is the direct sum of
$H^{i_1}(K)\otimes\cdots\otimes H^{i_k}(K)$ with sum of degrees $j$.
The extreme summands are the nonzero powers of $H^p$ and $H^r$.
Opposite-parity nonzero groups in degrees $i,j$ give nonzero summands
in degrees $ki$ and $(k-1)i+j$, of opposite parity. Exactness and
faithfulness detect these assertions in the original category.
Shifts, Tate twists and finite base-field extensions do not reduce
the amplitude. This argument does not assume an arbitrary convolution
or pushforward is exact for a chosen heart.

For the character model, the map evaluates $c$ at $\delta$ and zero.
Its kernel is zero since $R$ is a domain and $\delta\ne0$; its image
is the pairs whose difference is in $\delta R$. Localizing makes it
surjective. The four-coordinate tensor map is injective by successive
evaluation of $1,c_1,c_2,c_1c_2$, and direct evaluation of $P$ gives
the stated tuple. At $v=1$, $c_i=1+s_i$ and $P=2(1+s_1s_2)$,
so both equal-sign characters have eigenvalue four and both others
have eigenvalue zero. Every module over the localized algebra splits
by its four actual orthogonal character idempotents. It becomes a
quasi-coherent sheaf on four copies of the localized base, a vector
bundle when its corresponding modules are finite projective.

For an integer-Tate formulation set $b=vc$ and use
\[
 H_{\rm even}=\Z[u,u^{-1}][b]/(b^2-(1+u)b),\qquad
 P=u^{-1}\{b_1b_2+(b_1-1-u)(b_2-1-u)\}.
\]
Scalar extension recovers the original algebra. On each rational
character point, represent $b_i$ by $\mathbf1\oplus\mathbf1(-1)$
if its sign is $+$, and by zero otherwise. The element $P$ is then
represented on each diagonal point by
\[
 T=\mathbf1(1)\oplus\mathbf1^{\oplus2}\oplus\mathbf1(-1),
 \qquad \Tr(F^n\mid T)=q^{-n}+2+q^n,
\]
and by zero off the diagonal. Tensoring with this object on the
two supported points gives the corresponding positive operation on
the character Grothendieck group. It has weights $-2,0,0,2$;
it is not a single pure weight-zero representation.
\end{proof}

There is an actual Hecke category satisfying the first hypothesis:
the product of the two rank-one Soergel categories. For one factor
take $A=\Q[x]$, $\deg x=2$, $s(x)=-x$, and
$B=A\otimes_{A^s}A(1)$. The rank-two free decomposition over
$A^s=\Q[x^2]$ gives $B\otimes_A B=B(1)\oplus B(-1)$.
The split Grothendieck/KL identification is the known result
\cite[\S3.5, Theorem~3.5, pp.~1107--1108; Theorem~1.1, p.~1093]{EW}.
The bounded category used above is its homotopy category; no
unspecified perverse $t$-exactness was invoked.

The character map is also a finite normalization of affine schemes:
$R^2$ is the integral closure of $H$ in its total fraction ring,
and $R^4$ is that of the tensor-square algebra. Indeed the separate
branch idempotents are integral; conversely an element integral over
these finite $R$-algebras is integral over the normal ring $R$ and
has each coordinate in $R$. The map is an isomorphism on
$\delta\ne0$, but it is not an integral identification of the named
KL effective basis with the full character basis.

The finite-spectrum lemma of \fid{F24} shows why tensoring the
derived input does not itself manufacture the required opposite sign:
every nonempty visible real Euler spectrum has a positive top-radius
coefficient in each even tensor power. Input shifts contribute sign
$(-1)^{kr}=1$ there. An extra odd shift of the output reverses a sign
formally, but its identification with path-hypersurface cohomology
still requires the trace identity. Full scalar localization also
cannot be silently realized by finite Weil traces:
$(1+q^n)^{-1}$ is excluded by \fid{F26}. A particular localized
idempotent may have its denominator cancel in character coordinates,
as $b/(1+u)\mapsto(1,0)$ does. The obstruction concerns the specified
non-Laurent all-extension trace, not every object in a localized
category.


\subsection{The actual rank-three algebra has a coherent scheme model}
\begin{finding}[\fid{F19}: a finite-flat center, its collision, and the lost radical]\label{f:rankthree}
Put $\delta=v+v^{-1}$, $D=\delta^2$, and take either
\[
 R=\Q[v,v^{-1}]_{(v-1)}
 \quad\text{or}\quad
 R_0=\Z[1/2,v,v^{-1},\delta^{-1},(D-1)^{-1},(D-2)^{-1}].
\]
Write $R$ for the chosen ring in the formulas below. Let $A$ be the actual nonstandard Hecke algebra $\check H_3$ over this ring, with its specified generators $P_1,P_2$. Define
\[
 z_1=\frac{D-1}{D},\qquad z_2=\frac{4(D-1)}{D^2},\qquad
 C=R[z]/((z-z_1)(z-z_2)).
\]
Then there is an explicit algebra isomorphism
\begin{equation}\label{eq:r3model}
 A\simeq R\times R\times M_2(C).
\end{equation}
On the two scalar blocks $(P_1/D,P_2/D)$ is respectively $(1,1)$ and $(0,0)$; on the matrix block it is
\[
 e=\begin{pmatrix}1&0\\0&0\end{pmatrix},\qquad
 f=\begin{pmatrix}1-z&1\\z(1-z)&z\end{pmatrix}.
\]
In particular $A$ is free of rank ten. At $v=1$ over a field $k$ of characteristic different from 2 and 3,
\[
 A_1\simeq k\times k\times M_2(k[\varepsilon]/\varepsilon^2).
\]
The actual specialization map to the ambient tensor group algebra has kernel the four-dimensional square-zero Jacobson radical and image the diagonal $k[S_3]$.

Over the local characteristic-zero ring, the generic center branches collide to order two:
\begin{equation}\label{eq:r3collision}
 z_1-z_2=\frac{D-1}{D^2}(v-v^{-1})^2.
\end{equation}
The algebra is also the global coherent endomorphism algebra of an explicitly given vector bundle on the ordinary commutative scheme
$S=\Spec R\sqcup\Spec R\sqcup\Spec C$.
\end{finding}
\begin{proof}
The source input is the complete presentation in IV, Theorem~11.16, equations~(111)--(112), printed p.~65. With $x=P_1$, $y=P_2$, $a_1=\delta$ and $a_2=D-2$, it is
\[
 x^2=Dx,\quad y^2=Dy,\quad
 x(yx-a_1^2)(yx-a_2^2)=y(xy-a_1^2)(xy-a_2^2).
\]
Set $e=x/D$, $f=y/D$. These are idempotents, and $z=(e-f)^2=e+f-ef-fe$ is central: multiplying on either side by $e$ gives $e-efe$, and multiplying on either side by $f$ gives $f-fef$. Consequently
\[
 efe=(1-z)e,\qquad fef=(1-z)f.
\]
Expansion of the displayed braid relation gives
\[
 D^5(z-z_1)(z-z_2)(e-f)=0.
\]
For example, its two sides expand using $xyx=D^3(1-z)e$ and $xyxyx=D^5(1-z)^2e$, and the corresponding expressions in $f$; the two factors are $D^2(1-z)-a_i^2$. Thus the presentation is equivalent to the idempotent relations and $g(z)(e-f)=0$, where $g(z)=(z-z_1)(z-z_2)$.

Multiplication by $e-f$ gives $zg(z)=0$. Since $g(0)=z_1z_2$ is a unit, the central Chinese remainder decomposition splits the algebra into its $z=0$ and $g(z)=0$ blocks. In the first, $g(0)(e-f)=0$ forces $e=f$, giving $R[e]/(e^2-e)=R\times R$.

In the second block the central ring is $C$. Every word is a $C$-linear combination of $1,e,f,ef$, by removing repeated idempotents, shortening alternating triples with the preceding identities, and replacing $fe$ by $e+f-ef-z$. The matrices in the statement satisfy those identities. Both $z$ and $1-z$ are units in $C$: their evaluations at the two roots are units, since
\[
 1-z_1=1/D,\qquad 1-z_2=(D-2)^2/D^2.
\]
Equivalently their norms in the quadratic algebra are units, which gives their inverses. Now
\[
 ef(1-e)=E_{12},\qquad (1-e)fe=z(1-z)E_{21}.
\]
Together with $e$ and $1-e$ these generate every matrix unit. The matrices $1,e,f,ef$ are also independent: the lower-left coordinate of a relation first kills its $f$ coefficient; the upper-right then kills its $ef$ coefficient; the lower-right kills its scalar coefficient; the upper-left kills the remaining $e$ coefficient. This proves the isomorphism and freeness.

At $v=1$, $D=4$ and $z_1=z_2=3/4$, giving the dual-number description. Its ideal
$J=0\times0\times M_2(\varepsilon k)$ has square zero and semisimple quotient $k\times k\times M_2(k)$, so it is exactly the Jacobson radical and has dimension four. In the original ambient inclusion,
\[
 P_i=2(1+s_i\otimes s_i).
\]
These generate the diagonal group algebra, since $s_i\otimes s_i=P_i/2-1$. Its dimension is six and it is semisimple in the stated characteristics. The image of $J$ is a nilpotent ideal of that semisimple image and hence is zero; dimension then makes $J$ the entire kernel.

Equation~\eqref{eq:r3collision} follows from $D-4=(v-v^{-1})^2$. Evaluation at the two roots identifies
\[
 C=\{(u_1,u_2)\in R^2:u_1-u_2\in(z_1-z_2)R\}.
\]
Indeed $a+bz$ evaluates to a pair whose difference is $b(z_1-z_2)$; conversely such a pair gives this $b$ and $a=u_1-bz_1$. Injectivity follows because $R$ is a domain and the two generic roots differ. At $v=1$, the factor multiplying $(v-1)^2$ in \eqref{eq:r3collision} is a unit over the local characteristic-zero ring. The generic central idempotents divide by $z_1-z_2$ and do not extend through this collision; the finite algebra $C$ over this local ring has local special fiber and is itself local, so has no nontrivial idempotent.

Finally take the trivial vector bundle $E$ of ranks 1, 1 and 2 on the three components of $S$. Its algebra of global coherent endomorphisms is exactly $R\times R\times M_2(C)$. The displayed matrices retain the named nonstandard generators. This proves the coherent realization.
\end{proof}
This strengthens the construction beyond an abstract matching-dimension model. It also corrects a possible flatness misreading: the family remains rank ten; the \emph{specialized inclusion} loses injectivity. IV, Remark~11.9 (p.~61), already distinguishes abstract specialization from ambient image; Proposition~11.14 (pp.~64--65) gives the generic simple dimensions $1,1,2,2$ and the merging of the two two-dimensional representations.

Over a finite field of good characteristic at $v=1$, $S$ has two reduced rational points and one double rational point. Its geometric \(\ell\)-adic cohomology is that of three rational points, since a nilpotent thickening does not change the \emph{\'etale topos} (\cite[Proposition~59.45.4, tag~03SI]{StacksNil}). Thus it is pure in degree zero and has dimension three; this purity does not see the coherent radical. Replacing the three vector-bundle ranks by constant \(\ell\)-adic sheaves gives a chosen endomorphism algebra of dimension six on the reduced points, not a faithful realization of the ten-dimensional coherent algebra.

A full endomorphism algebra of a finite semisimple object over an algebraically closed coefficient field is a product of matrix algebras; it cannot be $A_1$. An algebra acting as a \emph{subalgebra} of such an endomorphism algebra can be nonsemisimple, so that broader construction is not excluded. The coherent bundle above models the algebra; \fid{F21} below supplies the actual right tensor module. Neither construction identifies VIII's based left/right structure or formal quantum grading with Frobenius weights. Those remain precise comparison tasks. In Grothendieck's language the result lives in finite-flat schemes and coherent $\End(E)$, not in an unproved RH; EGA~I's affine scheme formalism, SGA~3's module/comodule language, and cote~90, scan~17, supply the framework, while the complete algebra proof is given here.

\subsection{The actual tensor action and its coherent deformation}
\begin{finding}[\fid{F21}: the specified rank-three tensor lattice is realized]\label{f:r3tensor}
Use the ring $R_0$, algebra $A$, center $C$, roots $z_1,z_2$ and named
generators of \fid{F19}. Let $V=W=R_0^2$ and give
$M=(V\otimes W)^{\otimes3}$ the actual right tensor-monomial action of
GCT~IV, \S6.1, p.~33, and Definition~11.1, equation~(93). Write
$\epsilon_+$ for $P_i=D$, $\epsilon_-$ for $P_i=0$, and
\[
 N(a):\quad P_1=\begin{pmatrix}D&0\\a&0\end{pmatrix},\qquad
 P_2=\begin{pmatrix}0&a\\0&D\end{pmatrix}.
\]
These are right row matrices, as in IV Proposition~11.14,
equation~(110). There is an action-preserving isomorphism over $R_0$,
\begin{equation}\label{eq:r3tensor}
 M\simeq\epsilon_+^{20}\oplus\epsilon_-^4
       \oplus N(\delta)^{16}\oplus N(D-2)^4.
\end{equation}
In particular every displayed change of basis is regular at $v=1$ in
characteristic different from $2,3$; no colliding-root difference is inverted.

Put $L_i=C/(z-z_i)$ and
$S=\Spec R_0\sqcup\Spec R_0\sqcup\Spec C$. For
\[
 E=(\mathcal O,\mathcal O,\mathcal O_C^2),\qquad
 \mathcal F=(\mathcal O^{20},\mathcal O^4,
             \widetilde L_1^{16}\oplus\widetilde L_2^4),
\]
the coherent model realizes this same action as
\[
 M\simeq\Gamma(S,\mathcal Hom(E,\mathcal F)),\qquad A=\Gamma(S,\End E),
\]
with right action by precomposition. It is faithful before specialization
and kills the four-dimensional radical after $v=1$.
\end{finding}
\begin{proof}
Let $L=R_0^2$. The cited ordinary tensor action, with the consistent
normalization $c_i=T_i+v^{-1}$, is
\[
 \mathbf v_{\mathbf k}c_i=
 \begin{cases}
 v\mathbf v_{\mathbf k}+\mathbf v_{\mathbf k s_i},&k_i<k_{i+1},\\
 \delta\mathbf v_{\mathbf k},&k_i=k_{i+1},\\
 v^{-1}\mathbf v_{\mathbf k}+\mathbf v_{\mathbf k s_i},&k_i>k_{i+1}.
 \end{cases}
\]
On $(112,121,211)$ it gives
\[
 c_1=\begin{pmatrix}\delta&0&0\\0&v&1\\0&1&v^{-1}\end{pmatrix},\qquad
 c_2=\begin{pmatrix}v&1&0\\1&v^{-1}&0\\0&0&\delta\end{pmatrix}.
\]
The row-basis matrix
\[
 B_3=\begin{pmatrix}v^2&v&1\\-1-v^{-2}&v&1\\
                         -v^{-1}&-v^{-2}&\delta\end{pmatrix}
\]
has determinant $(D-1)^2$. Multiplying gives
\[
 B_3c_1B_3^{-1}=\begin{pmatrix}\delta&0&0\\0&\delta&0\\0&1&0\end{pmatrix},
 \quad
 B_3c_2B_3^{-1}=\begin{pmatrix}\delta&0&0\\0&0&1\\0&0&\delta\end{pmatrix}.
\]
Thus this block is $\epsilon\oplus U$, with
\[
 c_1|_U=\begin{pmatrix}\delta&0\\1&0\end{pmatrix},\qquad
 c_2|_U=\begin{pmatrix}0&1\\0&\delta\end{pmatrix}.
\]
On $(122,212,221)$ the original $c_1,c_2$ matrices interchange; applying
$B_3$ and swapping its last two rows again gives $\epsilon\oplus U$.
The weight spaces $111$ and $222$ are each $\epsilon$. Since the determinants
are units, this proves $L^{\otimes3}=\epsilon^4\oplus U^2$ over $R_0$,
not just over its fraction field.

By $P_i=c_i\otimes c_i+(c_i-\delta)\otimes(c_i-\delta)$,
$\epsilon\otimes\epsilon=\epsilon_+$, and each mixed tensor is
$N(\delta)$. On $U\otimes U$ in the order $(11,12,21,22)$,
\[
 P_1=\begin{pmatrix}D&0&0&0\\\delta&0&0&0\\\delta&0&0&0\\
                         2&-\delta&-\delta&D\end{pmatrix},\quad
 P_2=\begin{pmatrix}D&-\delta&-\delta&2\\0&0&0&\delta\\
                         0&0&0&\delta\\0&0&0&D\end{pmatrix}.
\]
Take
\[
 B_4=\begin{pmatrix}
 \delta&-1&-1&\delta\\0&-1&1&0\\
 2-D&-\delta&-\delta&D\\-D&\delta&\delta&D-2
 \end{pmatrix},\qquad \det B_4=8(D-1)^2.
\]
Direct multiplication yields
\[
 B_4P_1B_4^{-1}=\begin{pmatrix}D&0&0&0\\0&0&0&0\\0&0&D&0\\0&0&D-2&0\end{pmatrix},
\quad
 B_4P_2B_4^{-1}=\begin{pmatrix}D&0&0&0\\0&0&0&0\\0&0&0&D-2\\0&0&0&D\end{pmatrix}.
\]
This is $\epsilon_+\oplus\epsilon_-\oplus N(D-2)$, again over $R_0$.
In $(\epsilon^4\oplus U^2)\otimes(\epsilon^4\oplus U^2)$ the unmixed
trivial contribution has multiplicity $16$, the two mixed contributions
have total multiplicity $16$, and $U\otimes U$ occurs four times.
This proves \eqref{eq:r3tensor}, whose rank is $20+4+32+8=64$.

It remains to match the named action to the center model, rather than
merely match dimensions. For $a=\delta,D-2$, put $q_a=a/D$, a unit, and
\[
 S_a=\begin{pmatrix}1&0\\-q_a^2&q_a\end{pmatrix}.
\]
The matrices of $S_a(P_1/D)S_a^{-1}$ and $S_a(P_2/D)S_a^{-1}$ are
respectively
\[
 \begin{pmatrix}1&0\\0&0\end{pmatrix},\qquad
 \begin{pmatrix}1-z_i&1\\z_i(1-z_i)&z_i\end{pmatrix},
 \quad z_i=1-q_a^2.
\]
They are exactly the standard row module for $M_2(C)$ evaluated on branch
$i$. The right-module formula with $\mathcal Hom$ follows. The two branch
evaluations jointly embed $C$ into $R_0^2$, so their modules have joint
annihilator zero; both scalar blocks also occur. At $v=1$ the branch
modules become the same residue module of $k[\varepsilon]/\varepsilon^2$,
and the central nilpotent acts by zero. Its matrix ideal is the radical
of \fid{F19}, proving the last assertion.
\end{proof}
The generic multiplicities themselves are already in IV \S12.3; the
additional assertion here is the displayed regular action isomorphism.
The coherent sheaf $\mathcal F$ is not locally free on the collision
component: its two generic branch ranks are $16$ and $4$. This result
does not identify a simultaneous canonical left quantum-group action
or VIII's based tensor lattice with the monomial lattice.

\begin{finding}[\fid{F22}: constant nearby-cycle purity does not detect collision order]\label{f:collision}
Let $R=k[[t]]$ and, for $m\ge1$, set
\[
 C_m=\{(u,w)\in R^2:u\equiv w\pmod{t^m}\}
     \simeq R[x]/(x(x-t^m)).
\]
For $m\ge n\ge1$, $C_m\subset C_n$ induces a finite universal
homeomorphism $\Spec C_n\to\Spec C_m$ over $\Spec R$, although
\[
 \operatorname{length}_R(R^2/C_m)=m,\qquad
 \operatorname{cond}(C_m,R^2)=t^mR^2.
\]
Their corresponding ordinary \'etale coefficient categories and
cohomology are equivalent. With $k=\F_q$ and unshifted constant
$\Q_\ell$ coefficients, the unique special-point nearby and vanishing
cycles are
\[
 R\Psi\Q_\ell=\Q_\ell^2[0],\qquad
 R\Phi\Q_\ell=\Q_\ell[0],
\]
with trivial inertia and Frobenius. These constant data, including every
finite-extension trace and tensor power, are independent of $m$.
The completed center of \fid{F19} in good characteristic is $C_2$.
\end{finding}
\begin{proof}
Map $x$ to $(0,t^m)$ to obtain the displayed description. For $m\ge n$,
$C_n$ is a finite $C_m$-module since it is already generated by two
elements over the subring $R$. After inverting $t$, the map is the
identity on $k((t))^2$. Each special fiber has one point, with residue
field $k$, and the map induces the identity on that field. It is
therefore integral, surjective and radicial, hence a universal
homeomorphism. The difference map
$(u,w)\mapsto w-u\bmod t^m$ has kernel $C_m$ and gives
$R^2/C_m\simeq R/(t^m)$. The integral closure is $R^2$; a pair multiplies
every element of $R^2$ back into $C_m$ precisely when both coordinates
are divisible by $t^m$. This proves the coherent formulas.

Topological invariance of the \'etale site gives the claimed equivalence,
compatible with the base trait; precise locators are Stacks
Theorem~59.45.2, tag~04DZ, and Proposition~59.45.4, tag~03SI,
the former citing EGA~IV, 18.1.2 \cite{StacksNil}.
This is an equivalence of coefficient-sheaf topoi, not an identification
of coherent structure sheaves as ringed objects.

After strict henselization at the special point, the geometric generic
fiber is two points. Indeed the finite local algebra over the strictly
henselian trait is strictly henselian, and its two branches are already
split over the fraction field. The defining nearby-cycle stalk is
therefore the cohomology of two geometric points: $\Q_\ell^2$ in degree
zero, with trivial inertia and Frobenius. The specialization map from
$i^*\Q_\ell=\Q_\ell[0]$ is diagonal. The defining triangle
\[
 i^*\Q_\ell\longrightarrow R\Psi\Q_\ell
 \longrightarrow R\Phi\Q_\ell\longrightarrow i^*\Q_\ell[1]
\]
thus gives the stated cokernel in degree zero and no other groups.
This also follows by finite base change, Stacks Lemma~59.91.9,
tag~0A4E. Every object and map just described is independent of $m$.
The traces of tensor powers $r$ of $R\Psi$ and $R\Phi$ are respectively
$2^r$ and $1$ over every finite extension.

Finally \eqref{eq:r3collision} says that the actual root difference is
a unit times $(v-1)^2$. Translate $z$ by a root and rescale the translated
coordinate by that unit after completion. The relation becomes
$x(x-t^2)=0$, where $t=v-1$. This identifies the center with $C_2$.
\end{proof}
This is an information limit for the specified constant coefficient
families. A nilpotent algebra action supplied separately, or an
$m$-dependent coefficient construction, is additional information and
is not ruled out. A sign criterion applied only to these common traces
can establish their correct weight-zero purity but cannot determine
$m$. The next finding exhibits coherent information that does retain it.

\begin{finding}[\fid{F25}: the commutant defect, coherent extensions, and a regular basis]\label{f:r3deformation}
Use \fid{F21} over $R=\Q[v,v^{-1}]_{(v-1)}$ and write
$\Delta=z_1-z_2$. Then
\[
 \End_A(M)=M_{20}(R)\times M_4(R)\times M_{16}(R)\times M_4(R).
\]
This is free of rank $688$. Its base-change map to the full special
commutant
\[
 \End_{A_1}(M_1)=M_{20}(k)\times M_4(k)\times M_{20}(k)
\]
is injective with cokernel of dimension $128$; the latter commutant has
dimension $816$. Moreover
\[
 \operatorname{Ext}^1_C(L_1,L_2)
 =\operatorname{Ext}^1_C(L_2,L_1)=R/(\Delta),\qquad
 \operatorname{Ext}^1_C(L_i,L_i)=0,
\]
and $\operatorname{Ext}^1_A(M,M)=(R/(\Delta))^{128}$, of $R$-length $256$.

There is also a regular center basis $1,y$, with
\begin{equation}\label{eq:r3basis}
 y=D^2(z_1-z),\qquad
 y^2=\lambda y,\qquad
 \lambda=(v-v^{-1})^2(v^2+1+v^{-2}).
\end{equation}
Together with the two scalar-block units, $E_{ij},yE_{ij}$ form a basis
of the actual rank-ten algebra over $R_0$, with structure constants
among $0,1,\lambda$. This elementary basis is not asserted to be VIII's
designated canonical basis.
\end{finding}
\begin{proof}
Central blocks separate the two scalar sectors and the matrix sector.
Morita equivalence identifies the last sector's multiplicity module
with $L_1^{16}\oplus L_2^4$. A homomorphism $L_1\to L_2$ is determined
by $r\in R$ satisfying $(z_2-z_1)r=0$. Since $\Delta$ is a nonzero
element of the domain $R$, this Hom group is zero, as is the reverse
one; $\End_C(L_i)=R$. This gives the asserted commutant. At $v=1$ the
two branch modules become the same simple residue module. Their
multiplicity spaces merge into $k^{20}$. The map on these factors is
the block-diagonal embedding $M_{16}(k)\times M_4(k)\hookrightarrow
M_{20}(k)$; the missing off-diagonal blocks each have dimension $16\cdot4$.

For the extension calculation put $f_i=z-z_i$. The following is an
exact free resolution of $L_1$:
\[
 \cdots\xrightarrow{f_1}C\xrightarrow{f_2}C
             \xrightarrow{f_1}C\longrightarrow L_1\longrightarrow0.
\]
To check exactness, lift to the domain $R[z]$ and cancel $f_1$ or $f_2$
from divisibility by $f_1f_2$. Applying $\Hom_C(-,L_2)$ gives
$R\xrightarrow{-\Delta}R\xrightarrow0R\xrightarrow{-\Delta}R\to\cdots$,
so degree one is $R/(\Delta)$. Exchanging indices gives the reverse
group. Applying $\Hom_C(-,L_1)$ instead begins
$R\xrightarrow0R\xrightarrow\Delta R$ and has zero degree one.
Morita equivalence and the multiplicities in \eqref{eq:r3tensor} give
$16\cdot4+4\cdot16=128$ copies; scalar blocks contribute no extensions
of their free $R$-modules. Since $\Delta$ is a unit times $(v-1)^2$,
each crossed quotient has length two. For the general $C_m$ the same
resolution gives crossed extension group $R/(t^m)$, retaining precisely
the collision order omitted by the constant \'etale data.

Finally $D$ is a unit and the affine change $z\mapsto y$ is invertible.
The two evaluations of $y$ are $0$ and $D^2(z_1-z_2)$. Their product
relation yields $y^2=D^2\Delta y$. Since
\[
 D^2\Delta=(D-1)(D-4),\quad
 D-4=(v-v^{-1})^2,\quad D-1=v^2+1+v^{-2},
\]
this is \eqref{eq:r3basis}. Matrix-unit multiplication proves the
structure-constant assertion. At $v=1$, $\lambda=0$ and the four
$yE_{ij}$ span the radical.
\end{proof}
IV \S12.3 already records the numbers $688$ and $816$. Here the free
rank-$688$ commutant family is kept distinct from the full
special-fiber commutant, and their comparison map is explicit.

The scalar in \eqref{eq:r3basis} has a precise ordinary geometric trace:
with $Q=v^2$,
\[
 \lambda=Q^{-2}(Q-1)^2(1+Q+Q^2)
       =Q^{-2}\#(\mathbb G_m^2\times\PP^2)(\F_Q).
\]
It is the compact-support alternating trace with Tate twist $(2)$.
The remaining factor $v^{-2}(1+v^2+v^4)$ is the centered Tate polynomial
of $\PP^2$. This mixed-cohomology trace interpretation of one scalar
does not supply functorial multiplication or associativity constraints,
and it does not assert VIII's consecutive-$v$ unimodality convention.
Indeed $\lambda=v^4-v^2-v^{-2}+v^{-4}$ has negative Laurent coefficients.


\subsection{Coherent Frobenius: the collision obstruction and an actual normalization repair}
Here Frobenius means absolute Frobenius on characteristic-$p$ coherent
objects. It is not yet geometric Frobenius on an $\ell$-adic Weil
realization. The distinction between pullback on the singular center
and pullback on the parameter trait gives different, computable answers.

\begin{finding}[\fid{F36}: unit Frobenius depends on the coherent base]\label{f:coherentfrobenius}
Let $k=\F_q$ have characteristic $p>3$, let $R=k[[t]]$, and put
\[
 \Delta=t^m u(t),\quad u(t)\in R^*,\quad m\ge1,\qquad
 C=R[x]/(x(x-\Delta)),\quad L_1=C/(x),\quad L_2=C/(x-\Delta).
\]
For absolute Frobenius $F_C(c)=c^p$,
\[
 F_C^*L_i\simeq R\oplus R/(\Delta^{p-1})
 \quad\text{as }R\text{-modules}.
\]
Thus no unit $C$-linear Frobenius isomorphism $F_C^*L_i\simeq L_i$
exists. Every positive Tor degree of the derived pullback is also
nonzero. In the actual algebra and 64-dimensional module of
\fid{F19}/\fid{F21}, the center has $m=2$ and underived algebra
pullback adds torsion of $R$-length $80(p-1)$.

Nevertheless the actual module has a unit \emph{parameter-base}
semilinear Frobenius structure, compatible with the intrinsic
entrywise Frobenius on $A=R\times R\times M_2(C)$ after transport
through its regular adapted basis. On the normalized center the
actual module has a unit coherent Frobenius model preserving its
original $A$-action under restriction. Pushing this model back through
the collision gives the nonunit surjection just computed.
\end{finding}
\begin{proof}
\par\smallskip\noindent\textbf{The two branches.}\quad
Since $x^2=\Delta x$, characteristic $p$ gives
\[
 x^p=\Delta^{p-1}x,\qquad
 (x-\Delta)^p=\Delta^{p-1}(x-\Delta).
\]
Tensoring the cyclic presentations along $F_C$ gives
\[
 F_C^*L_1=C/(\Delta^{p-1}x),\qquad
 F_C^*L_2=C/(\Delta^{p-1}(x-\Delta)).
\]
The ring $C$ is free over $R$ with basis $1,x$. The first ideal is
exactly $R\Delta^{p-1}x$, because
$x(a+bx)=(a+b\Delta)x$. This proves the stated $R$-module
decomposition for $L_1$; use the basis $1,x-\Delta$ for $L_2$.
The torsion has length $m(p-1)$, while each $L_i$ is $R$-free.
Hence no possible $C$-linear isomorphism can be a unit structure.
The decomposition asserted here is only over $R$, not a split
decomposition of $C$-modules.

There is a natural surjection $\theta_i:F_C^*L_i\to L_i$,
since $(x_i^p)\subset(x_i)$, where $x_1=x$ and $x_2=x-\Delta$.
It linearizes the branch map $r\mapsto r^p$. Its precise kernels are
\[
 \ker\theta_1\simeq L_2/\Delta^{p-1}L_2,
 \qquad \ker\theta_2\simeq L_1/\Delta^{p-1}L_1:
\]
the first is generated by the class of $x$, annihilated by $x-\Delta$,
and the second by the class of $x-\Delta$, annihilated by $x$.
For the $p^e$-power Frobenius all formulas replace $p-1$ by
$p^e-1$; in particular the coefficient-field-linear $q$th-power
version has length $m(q-1)$ per branch.

\par\smallskip\noindent\textbf{Derived pullback on the same object.}\quad
Put $f_1=x$, $f_2=x-\Delta$. Cancellation in the domain $R[X]$
gives $\operatorname{ann}_C(f_1)=(f_2)$ and
$\operatorname{ann}_C(f_2)=(f_1)$, so
\[
 \cdots\xrightarrow{f_1}C\xrightarrow{f_2}C
       \xrightarrow{f_1}C\longrightarrow L_1\longrightarrow0
\]
is an exact free resolution. Frobenius pullback replaces its maps by
$\Delta^{p-1}f_1$ and $\Delta^{p-1}f_2$.
The scalar $\Delta$ is a nonzerodivisor on the $R$-free ring $C$,
so their kernels are still $(f_2)$ and $(f_1)$, while the preceding
images are only $\Delta^{p-1}(f_2)$ and $\Delta^{p-1}(f_1)$.
Writing $C_{F_C}$ for the right module twisted by Frobenius, we obtain
\[
 \operatorname{Tor}^{C}_{2j+1}(C_{F_C},L_1)
     \simeq L_1/\Delta^{p-1}L_1\quad(j\ge0),
\]
\[
 \operatorname{Tor}^{C}_{2j}(C_{F_C},L_1)
     \simeq L_2/\Delta^{p-1}L_2\quad(j\ge1).
\]
Exchange the branches for $L_2$. Every positive degree has
underlying $R$-module $R/(\Delta^{p-1})$.
Thus $\mathbf L F_C^*L_i$ is unbounded below and cannot equal the
degree-zero branch object. Changing its resolution cannot repair
this; replacing the object or using an unbounded category is a
different construction.

\par\smallskip\noindent\textbf{The actual algebra action and multiplicities.}\quad
The completed good-characteristic model from \fid{F19} has
\[
 A=R\times R\times M_2(C).
\]
Define $\phi_A$ by $F_R$ on the scalar blocks and entrywise $F_C$
on the matrix block. It is a ring homomorphism because, in the
commutative coefficient ring $C$,
\[
 \left(\sum_j a_{ij}b_{j\ell}\right)^p
       =\sum_j a_{ij}^p b_{j\ell}^p.
\]
This is not the noncommutative matrix-power map of \fid{F11}.
It fixes the matrix units, so its extension-of-scalars functor
commutes with the standard Morita equivalence.

The actual right module is
\[
 M=R^{20}\oplus R^4\oplus
    (L_1\otimes_C C^{1\times2})^{16}
       \oplus(L_2\otimes_C C^{1\times2})^4,
\]
through the regular action-preserving basis of \fid{F21}. Under
$F_A^*M=M\otimes_{A,\phi_A}A$, each branch multiplicity module
becomes $F_C^*L_i$; the two scalar sectors stay free. Therefore
\begin{equation}\label{eq:frobtorsion}
 F_A^*M\simeq R^{64}\oplus
          (R/(\Delta^{p-1}))^{40}
 \quad\text{as }R\text{-modules}.
\end{equation}
The branch multiplicity sheaf has 20 copies and torsion length
$20m(p-1)$; the actual row module has 40 copies and length
$40m(p-1)$. The natural algebra-linearized surjection discards
exactly this torsion. Derived Morita equivalence transfers the
resolution above as well: every positive Tor degree has 40 copies
of that torsion module, and the scalar sectors contribute none.

For the original center take $t=v-1$, $x=z_1-z$ and
$\Delta=z_1-z_2$. Its formula is
\[
 \Delta=\frac{(D-1)(D-4)}{D^2}
        =\frac34t^2+O(t^3).
\]
In characteristic greater than three its order is exactly two,
giving the claimed length $80(p-1)$.

\par\smallskip\noindent\textbf{Parameter pullback and the regular unit repair.}\quad
Parameter-base pullback of the center gives
\[
 C^{(p)}=R\otimes_{F_R,R}C
       \simeq R[X]/(X(X-\Delta^p)).
\]
Relative Frobenius sends $X$ to $x^p=\Delta^{p-1}x$.
In the two free bases its matrix is
$\operatorname{diag}(1,\Delta^{p-1})$; it is injective with
cokernel of length $m(p-1)$. Its source has branch separation of
order $mp$, whereas the target has order $m$. Extension over this
relative map is the additional operation in center pullback which
causes the torsion.

In the adapted free $R$-basis of $M$, take coefficientwise
Frobenius on every scalar and branch row. Branch evaluation
$C\to R$ commutes with Frobenius, so this semilinear map satisfies
\[
 \tau_{\rm ad}(ma)=\tau_{\rm ad}(m)\phi_A(a).
\]
Its parameter linearization
$R\otimes_{F_R,R}M\to M$ is an isomorphism, sending each
$1\otimes e_j$ to $e_j$. This tensor product is different from
the algebra tensor product in \eqref{eq:frobtorsion}.

The transport to the original tensor basis is explicit. With the
right-row convention of \fid{F21}, write
\[
 \rho_{\rm ad}(a)=B\rho_{\rm old}(a)B^{-1},\qquad
 r_{\rm old}=r_{\rm ad}B,\qquad
 U=(B^{[p]})^{-1}B,
\]
where $B^{[p]}$ means entrywise Frobenius and $B$ is the regular
invertible adapted basis matrix. Then
\begin{equation}\label{eq:frobtransport}
 \tau_{\rm old}(r)=r^{[p]}U,\qquad
 \rho_{\rm old}(\phi_A(a))
       =U^{-1}\rho_{\rm old}(a)^{[p]}U.
\end{equation}
Indeed transform $r$ to $rB^{-1}$, apply entrywise Frobenius,
and multiply back by $B$; applying the same operation to the
action matrix proves the second equality. It also directly checks
$\tau_{\rm old}(r\rho(a))=\tau_{\rm old}(r)\rho(\phi_A(a))$.
The matrix $U$ is invertible, proving the unit parameter structure
on the actual lattice.

Naive entrywise Frobenius in the original tensor basis is not
assumed to preserve the original algebra image. Although its
generator entries have the numerical form $P_i(v)^{[p]}=P_i(v^p)$,
the intrinsic algebra map acts by the conjugated expression in
\eqref{eq:frobtransport}. The regular conjugation cannot be omitted
without an additional argument.

\par\smallskip\noindent\textbf{A unit model on the normalization.}\quad
The normalization is
\[
 C\hookrightarrow\bar C=R\times R,
 \qquad a+bx\longmapsto(a,a+b\Delta).
\]
Its image is the pairs congruent modulo $\Delta$, and its quotient
has length $m$. It is the full normalization: the branch idempotents
are integral, and each factor $R$ is integrally closed.
Extending the algebra to this normalized center gives
\[
 A\hookrightarrow\bar A
       =R\times R\times M_2(R)\times M_2(R),
 \qquad \operatorname{length}_R(\bar A/A)=4m.
\]
Let $f:\bar S\to S$ normalize the three-component scheme in
\fid{F21}; it has four components upstairs. Choose free bundles
\[
 \bar E=(R,R,R^2,R^2)=f^*E,\qquad
 \bar{\mathcal F}=(R^{20},R^4,R^{16},R^4).
\]
They satisfy exactly
\[
 f_*\bar{\mathcal F}=\mathcal F,\qquad
 f_*\mathcal Hom(f^*E,\bar{\mathcal F})
       \simeq\mathcal Hom(E,\mathcal F).
\]
The last is ordinary tensor--Hom adjunction for the finite locally
free $E$. Its right $A$ action is restriction of the $\bar A$
action along the displayed inclusion. It therefore realizes the
same actual module, not just the same dimension or character.

On the normalized free branches, the chosen bases give unit
coherent Frobenius structures on $\bar E$ and
$\bar{\mathcal F}$ and hence the corresponding row module.
Entrywise Frobenius on $\bar A$ commutes with its inclusion of
$A$, since branch evaluations commute with absolute Frobenius.
But the natural comparison
\[
 F_S^*f_*\bar{\mathcal F}
       \longrightarrow f_*F_{\bar S}^*\bar{\mathcal F}
       \simeq f_*\bar{\mathcal F}
\]
is, on each branch, precisely $C/(x_i^p)\twoheadrightarrow C/(x_i)$.
It has torsion length $20m(p-1)$ on the multiplicity sheaf and
$40m(p-1)$ on the actual row module. Thus the unit upstairs
becomes the computed nonunit map downstairs. After inverting $t$
the normalization is an isomorphism and the entire defect vanishes.
\end{proof}


Thus the failed claim is a unit center/algebra pullback on the
specified branch object; the successful claims are the unit
parameter-base structure and the action-compatible normalized model.
Replacing the object, permitting unbounded constructions, or passing
to an \'etale/Weil realization are separate operations. None here
identifies quantum degree with a Weil weight, chooses VIII's basis,
or proves the opposite path-output sign needed by the main criterion.

\subsection{Testing the smooth proper noncommutative Weil framework}
There is a precise finite-field noncommutative Weil conjecture in a
different framework: for a smooth proper dg category, cyclotomic
Frobenius acts on the two finite-dimensional groups $TP_0[1/p]$
and $TP_1[1/p]$, with expected conjugate moduli $1$ and $\sqrt q$.
Tabuada \cite[\S1, pp.~2--3; \S5.1, pp.~13--14; \S6.1,
pp.~15--16]{TabuadaWeil} specifies these objects and proves the
comparison with smooth proper schemes in Theorem~1.5.
This cyclotomic Frobenius is not defined by raising matrices to
their $p$th powers. Its periodicity variable is also not the GCT
quantum parameter. The following test concerns the actual algebra
already constructed, without identifying the two conjectural programs.

\begin{finding}[\fid{F39}: the actual collision fiber fails homological smoothness]\label{f:ncweil}
In characteristic $p>3$, the abstract specialization of the
rank-three algebra in \fid{F19},
\[
 A_1=k\times k\times M_2(B),\qquad B=k[\epsilon]/(\epsilon^2),
\]
is proper but not smooth as a dg algebra concentrated in degree zero.
The same holds for its category of perfect complexes. Therefore this
particular category is not an input to the cited smooth proper
noncommutative Weil conjecture.

In contrast, at a split finite-field parameter point where the two
center roots are distinct and the stated denominators are units,
the algebra is $k^2\times M_2(k)^2$. Its perfect category is Morita
equivalent to four rational points. Its normalized even Frobenius
spectrum consists of four ones and its odd spectrum is empty.
Replacing $A_1$ by its semisimple quotient instead gives three
points and discards the nilpotent collision data.
\end{finding}
\begin{proof}
Properness of the one-object dg algebra means that its underlying
complex over $k$ is perfect; this holds because $A_1$ is finite
dimensional. Homological smoothness would make its diagonal bimodule
perfect over $A_1\otimes_k A_1^{\rm op}$.
We test the dual-number block directly.
Put
\[
 S=B\otimes_k B=k[x,y]/(x^2,y^2),\qquad
 u=x-y,\quad w=x+y.
\]
There is a free $S$-resolution of the diagonal $B$,
\[
 \cdots\xrightarrow{u}S\xrightarrow{w}S
       \xrightarrow{u}S\longrightarrow B\longrightarrow0.
\]
Here the rightmost positive differential is multiplication by $u$.
To check exactness in full, use the basis $1,x,y,xy$ of $S$.
The kernel of the augmentation $x,y\mapsto\epsilon$ is
$\langle x-y,xy\rangle=\operatorname{im}(u)$.
Multiplication by $u$ sends
$a+bx+cy+dxy$ to $a(x-y)+(c-b)xy$, so its kernel is
$\langle x+y,xy\rangle=\operatorname{im}(w)$.
Multiplication by $w$ has kernel
$\langle x-y,xy\rangle=\operatorname{im}(u)$.
Also $uw=0$. These identities prove exactness at every stage.

After tensoring over $S$ with the diagonal $B$, the differentials
alternate between zero and multiplication by $2\epsilon$. Hence
\[
 HH_0(B)=B,\qquad
 HH_{2j+1}(B)=B/(\epsilon)\ (j\ge0),\qquad
 HH_{2j}(B)=k\epsilon\ (j\ge1).
\]
In particular $\operatorname{Tor}^{S}_i(B,B)$ is nonzero for
every $i>0$. A perfect $S$-module would have bounded derived tensor
product with $B$, so $B$ is not homologically smooth.
If $A_1$ were smooth, restricting its diagonal along the central
matrix-block idempotent would make the diagonal of $M_2(B)$
perfect. The matrix Morita equivalence, obtained by the full matrix
unit $E_{11}$ on each side, is exact and preserves finite projective
resolutions; it would then make $B$ perfect over $S$, a contradiction.
This also proves the assertion for the Morita-invariant perfect dg
category, using the smoothness definition in the cited source.

At a split good parameter the distinct-root Chinese remainder
decomposition of \fid{F19} gives $k^2\times M_2(k)^2$.
Applying ordinary matrix Morita equivalence to its two matrix factors
identifies its perfect category with that of $\coprod_1^4\Spec k$.
The latter is smooth proper. The crystalline comparison
\cite[\S7, equations~(7.1)--(7.4), p.~17]{TabuadaWeil}
has only four degree-zero classes, each with Frobenius one; this
proves the stated normalized spectrum. The radical of $A_1$ is
the square-zero matrix ideal $M_2(k\epsilon)$, so its quotient is
$k^2\times M_2(k)$ and the same calculation gives three points.
\end{proof}

The dual-number resolution is the familiar elementary computation;
its use here is a test of the actual nonstandard collision fiber,
not a new claim about dual numbers or priority over noncommutative
motivic theory. A smooth categorical replacement or a relative theory
is not excluded. Such a replacement would need a specified comparison
retaining the coherent specialization and the designated based
operations. The pure spectrum at the good fibers does not supply
that comparison, and semisimplification expressly loses data measured
by \fid{F25}. The smoothness failure also does not refute any
conjecture whose domain already requires smoothness.

\subsection{A smooth category retaining the actual collision corner}
The direct smoothness failure in \fid{F39} has an explicit repair
at the level of a larger category. The following is an application
of the classical Auslander construction to the actual algebra and
module already identified in \fid{F19}/\fid{F21}; no priority claim
for the Auslander algebra is intended.

\begin{finding}[\fid{F40}: a smooth proper corner and an exact module lift]\label{f:auslander}
Let $k=\F_q$ have characteristic greater than three, and put
\[
 B=k[\epsilon]/(\epsilon^2),\qquad
 \Gamma=\End_B(B\oplus k),\qquad
 \Lambda=k\times k\times M_2(\Gamma).
\]
Then $\dim\Gamma=5$, $\dim\Lambda=22$, and both algebras are
smooth proper of global dimension two. For an explicit idempotent
$f\in\Lambda$,
\[
 f\Lambda f=A_1=k\times k\times M_2(B).
\]
There are projective right $\Lambda$-modules
$\mathcal P=f\Lambda$ and $\widetilde M$, with
$\dim_k\widetilde M=104$, such that
\[
 \End_\Lambda(\mathcal P)=A_1,\qquad
 \operatorname{Hom}_\Lambda(\mathcal P,\widetilde M)
   \simeq M_1
\]
as the actual right $A_1$-module of dimension 64 in \fid{F21}.
Thus the full algebra, including its nilpotents, and the specified
specialized tensor action are recovered.

The induced functor
$\operatorname{Perf}(A_1)\to\operatorname{Perf}(\Lambda)$ is fully
faithful but not right-admissible. The noncommutative Chow motive
of $\Lambda$ is a sum of four unit motives. Its normalized
cyclotomic Frobenius therefore has four eigenvalues one on
$TP_0[1/p]$, and $TP_1[1/p]=0$.
No relative deformation or based-operation comparison is asserted.
\end{finding}

\begin{proof}
Let $e_1,e_2$ be the idempotents of the summands $B,k$.
Let $a:k\to B$ send $1$ to $\epsilon$, let $b:B\to k$ be
the quotient, and let $c=ab$. Multiplication is composition,
with the right factor applied first. The five elements
$e_1,e_2,a,b,c$ form a basis and satisfy
\[
 a=e_1ae_2,\qquad b=e_2be_1,\qquad ba=0,\qquad
 c=ab\ne0,\qquad c^2=0.
\]
Indeed the four matrix Hom spaces have dimensions $2,1,1,1$,
and these maps span them. The radical is
$J=\langle a,b,c\rangle$, with $J^2=kc$, $J^3=0$ and
$\Gamma/J=k^2$.

Write
\[
 \mathsf P_1=e_1\Gamma=\langle e_1,a,c\rangle,\qquad
 \mathsf P_2=e_2\Gamma=\langle e_2,b\rangle
\]
and let $S_i$ be their simple tops.
Left multiplication by $a$ sends $e_2,b$ to $a,c$, and left
multiplication by $b$ sends $e_1$ to $b$ while killing $a,c$.
Thus the following are exact projective resolutions:
\[
 0\longrightarrow\mathsf P_2\xrightarrow{a}\mathsf P_1
   \longrightarrow S_1\longrightarrow0,
\]
\begin{equation}\label{eq:auslandersimple}
 0\longrightarrow\mathsf P_2\xrightarrow{a}\mathsf P_1
   \xrightarrow{b}\mathsf P_2\longrightarrow S_2
   \longrightarrow0.
\end{equation}
Every right module has a radical filtration of length at most
three, with semisimple quotients. Direct sums of these resolutions
and the extension bound for projective dimension show that every
module has projective dimension at most two. Applying
$\operatorname{Hom}_\Gamma(-,S_2)$ to
\eqref{eq:auslandersimple} gives
$\operatorname{Ext}^2_\Gamma(S_2,S_2)=k$, since
$\operatorname{Hom}(\mathsf P_1,S_2)=0$ and
$\operatorname{Hom}(\mathsf P_2,S_2)=k$. Hence the global
dimension is exactly two.

We prove smoothness directly. All tensor products in the following
diagonal resolution are over $k$. Set
\[
 Q_0=(\Gamma e_1\otimes e_1\Gamma)
          \oplus(\Gamma e_2\otimes e_2\Gamma),
\]
\[
 Q_1=(\Gamma e_1\otimes e_2\Gamma)_a
          \oplus(\Gamma e_2\otimes e_1\Gamma)_b,\qquad
 Q_2=\Gamma e_2\otimes e_2\Gamma.
\]
These are projective $\Gamma$-bimodules. Let $\mu:Q_0\to\Gamma$
be multiplication and define
\[
 d_1(x\otimes y)_a=xa\otimes y-x\otimes ay,\qquad
 d_1(x\otimes y)_b=xb\otimes y-x\otimes by,
\]
\[
 d_2(x\otimes y)=(xb\otimes y)_a+(x\otimes ay)_b.
\]
The subscripts specify the summands; the terms of $d_1$ lie in
the corresponding source and target vertex summands of $Q_0$.
These are bimodule maps, with $\mu d_1=0$ and
\[
 d_1d_2(x\otimes y)=xba\otimes y-x\otimes bay=0.
\]
The left ideals have bases
$\Gamma e_1=\langle e_1,b,c\rangle$ and
$\Gamma e_2=\langle e_2,a\rangle$. Consequently
\[
 \dim Q_2=4,\quad\dim Q_1=12,\quad\dim Q_0=13,\quad
 \dim\Gamma=5.
\]
The map $\mu$ is surjective. The image of $d_1$ contains the
eight independent vectors
\[
 \begin{gathered}
 a\otimes e_2-e_1\otimes a,\quad
 a\otimes b-e_1\otimes c,\quad
 b\otimes e_1-e_2\otimes b,\quad
 c\otimes e_1-a\otimes b,\\
 b\otimes a,\quad b\otimes c,\quad
 c\otimes a,\quad c\otimes c.
 \end{gathered}
\]
Independence follows from the tensor bases: the first four have
unique coordinates at $a\otimes e_2,e_1\otimes c,e_2\otimes b,
c\otimes e_1$, respectively, and the last four are distinct
remaining coordinates. Thus $\operatorname{im}d_1=\ker\mu$.
The $a$-component of $d_2$ is the tensor product of the injection
$x\mapsto xb:\Gamma e_2\to\Gamma e_1$, which sends $e_2,a$
to $b,c$, with the identity of $e_2\Gamma$. Hence $d_2$ is
injective. Its image has dimension four, equal to
$\dim\ker d_1$, proving exactness of
\begin{equation}\label{eq:auslanderdiagonal}
 0\longrightarrow Q_2\xrightarrow{d_2}Q_1
   \xrightarrow{d_1}Q_0\xrightarrow{\mu}\Gamma
   \longrightarrow0.
\end{equation}
This is a finite projective diagonal resolution. It proves
homological smoothness; finite dimensionality proves properness.
Matrix Morita equivalence and finite products give the same
properties and global dimension for $\Lambda$.

Now put $f=(1,1,I_2\otimes e_1)$. The corner equality follows
from $e_1\Gamma e_1=\langle e_1,c\rangle\simeq B$, with
$\epsilon\mapsto c\ne0$. In particular all named generators in
\fid{F19} are retained: its normalized matrix formulas use
$z=3/4+c$ at the collision. Left multiplication identifies
$\End_\Lambda(f\Lambda)=f\Lambda f=A_1$.

For the module recovery, write the actual specialization from
\fid{F21} as
\[
 M_1=\epsilon_+^{20}\oplus\epsilon_-^4\oplus S^{20},
\]
where $S=k^2$ is the right row module for $M_2(B)$, with
$\epsilon$ acting as zero. The projective right $M_2(\Gamma)$-module
\[
 Q=(E_{11}\otimes e_2)M_2(\Gamma)
     \simeq(\mathsf P_2,\mathsf P_2)
\]
has dimension four. Its corner is
\[
 Q(I_2\otimes e_1)=(\mathsf P_2e_1)^2
       =(kb)^2\simeq S,
\]
as an actual module, since $bc=bab=0$. Therefore
\[
 \widetilde M=\epsilon_+^{20}\oplus\epsilon_-^4\oplus Q^{20}
\]
is projective of dimension $104$, and
$\widetilde M f=M_1$. Evaluation at $f$ identifies
$\operatorname{Hom}_\Lambda(f\Lambda,\widetilde M)$ with this
corner. Precomposition by $\End_\Lambda(f\Lambda)$ is precisely
the right $A_1$-action, proving the claimed action-preserving
comparison.

For the embedding and its limitation, first consider
\[
 I:\operatorname{Perf}(B)\longrightarrow\operatorname{Perf}(\Gamma),
 \qquad X\longmapsto X\otimes_B^{\mathbf L}e_1\Gamma.
\]
It sends the generator $B$ to the projective $\mathsf P_1$.
Its right adjoint on ambient derived categories is corner
restriction $Y\mapsto Ye_1$. The unit is an isomorphism because
\[
 (X\otimes_B^{\mathbf L}e_1\Gamma)e_1
       =X\otimes_B^{\mathbf L}B=X.
\]
Thus $I$ is fully faithful. The matrix and scalar factors give
the assertion for $A_1,\Lambda$.

However $\mathsf P_2e_1=kb\simeq k$ is not perfect over $B$:
the exact resolution
\[
 \cdots\xrightarrow{\epsilon}B\xrightarrow{\epsilon}B
   \xrightarrow{\epsilon}B\longrightarrow k\longrightarrow0
\]
gives $\operatorname{Tor}^B_i(k,k)=k$ in every degree.
If $I$ had a right adjoint valued in $\operatorname{Perf}(B)$,
its value at $\mathsf P_2$ would represent
$X\mapsto\operatorname{Hom}_{D(B)}(X,k)$ on perfect $X$.
Evaluation at $B[n]$ would force the representing object to
have only the cohomology module $k$ in degree zero; naturality
for $\End_B(B)=B$ fixes its $B$-action. It would be isomorphic
to the nonperfect object $k$, a contradiction. This also
obstructs right admissibility of the actual matrix corner.

Finally $\langle S_1,\mathsf P_2\rangle$ is a full exceptional
pair in $\operatorname{Perf}(\Gamma)$.
The projective $\mathsf P_2$ has endomorphism ring $k$.
The first simple resolution gives
$\End(S_1)=k$ and no positive self-Ext, because
$\operatorname{Hom}(\mathsf P_2,S_1)=S_1e_2=0$.
Also $R\operatorname{Hom}(\mathsf P_2,S_1)=0$, so the pair
is semiorthogonal. Its triangle recovers $\mathsf P_1$, hence
both projective generators. Each exceptional subcategory has
both adjoints given by finite derived Hom and its dual, so this
is a semiorthogonal decomposition into two copies of
$\operatorname{Perf}(k)$.

Known semiorthogonal additivity
\cite[Example~12.3, p.~24]{TabuadaWeil} now gives
\[
 U(\Gamma)_{\Q}=U(k)_{\Q}^{\oplus2},\qquad
 U(\Lambda)_{\Q}=U(k)_{\Q}^{\oplus4}.
\]
The Frobenius-enriched TP functor
\cite[\S6.2, equation~(6.9), p.~16]{TabuadaWeil}
preserves these isomorphisms with Frobenius. The point comparison
\cite[\S7, equations~(7.1)--(7.4), p.~17]{TabuadaWeil}
then gives $TP_0(\Lambda)[1/p]=K^4$ with $F_0=1$ and
$TP_1(\Lambda)[1/p]=0$, where $K=\operatorname{Frac}W(k)$.
This is also an instance of the cited connective smooth proper
theorem and its motive computation
\cite[Theorem~2.12, p.~8; proof, equation~(12.6), p.~26]{TabuadaWeil}.
The source hypotheses hold by \eqref{eq:auslanderdiagonal} and
finite dimensionality. The special-case exceptional calculation
above avoids any assumption of nilinvariance for arbitrary
nonsmooth algebras.
\end{proof}

The original $M_1$ is not perfect over $A_1$. Its recovery here
is a corner restriction of a projective object, not the image of
$M_1$ under the perfect-module embedding. Its higher coherent
self-Ext need not agree with that of $\widetilde M$, and the
relative collision order in \fid{F25} has not been recovered
from this single fiber. The pure spectrum of the entire smooth
category does not identify a Frobenius realization of every
specified corner morphism, a positive nonstandard basis, or its
structure constants. The direct smoothness gate is nevertheless
strictly narrowed: this category retains the actual collision
algebra and tensor action as stated, without semisimplifying them.


\subsection{A normalization that cannot be a finite Frobenius trace}
\begin{finding}[\fid{F26}: rational normalization versus finite Weil spectra]\label{f:normalizedtrace}
Let $b>1$ be real and $f(x)\in\C(x)$. The eventual sequence $f(b^n)$ is
a finite exponential polynomial
\[
 \sum_{j=1}^s p_j(n)\alpha_j^n,
 \qquad p_j\in\C[n],\quad \alpha_j\in\C^*,
\]
if and only if $f$ is a Laurent polynomial.

On the \emph{actual} rank-three tensor module $M$ of \fid{F21}, set
$e_i=P_i/D$, the normalized algebraic idempotents. Then
\[
 \Tr_M(e_1e_2)=24+\frac{16}{D^2},\qquad
 \Tr_M(P_1P_2)=24D^2+16.
\]
For a prime power $q>1$ and either substitution $v=q^{n/2}$ or $v=q^{-n/2}$, the first trace
is not a finite exponential polynomial in $n$, even eventually.
It therefore cannot be the alternating trace of a fixed correspondence
against $F^n$ on any bounded finite-cohomology Weil realization.
This refutes the literal arithmetic word-trace comparison described
below; it does not assert that the GCT sources require that comparison.
\end{finding}
\begin{proof}
If $f=\sum a_jx^j$ is Laurent, then $f(b^n)=\sum a_j(b^j)^n$.
Conversely a finite exponential polynomial satisfies an eventual
constant-coefficient recurrence
\[
 \sum_{j=0}^r c_j f(b^{n+j})=0,\qquad c_0\ne0.
\]
For example its annihilator is a product of powers of the shift
operators $S-\alpha_j$, with nonzero constant term. The rational
function $R(x)=\sum_{j=0}^r c_j f(b^jx)$ vanishes at infinitely many
distinct $b^n$, so $R=0$. If $f$ has a finite nonzero pole $x_0$ of
maximal absolute value, none of $f(b^jx)$ for $j\ge1$ has a pole at
$x_0$: that would require the larger-modulus pole $b^jx_0$ of $f$.
Thus only $c_0f(x)$ has a pole there, a contradiction. A rational
function with poles only at $0$ and infinity is Laurent. Eventual
evaluation is always defined, since a rational function has only
finitely many poles. The same argument applies to $f(b^{-n})$ by
replacing $f(x)$ with $f(x^{-1})$.

In the source row matrices $N(a)$ from \fid{F21},
$\Tr(e_1e_2)=a^2/D^2$. The two scalar sectors contribute $20$ and zero.
Consequently
\[
 \Tr_M(e_1e_2)
 =20+\frac{16}{D}+\frac{4(D-2)^2}{D^2}
 =24+\frac{16}{D^2}.
\]
Multiplication by $D^2$ gives the unnormalized trace. Under the stated
substitutions, $D=q^n+2+q^{-n}$ and the sequence is
\[
 h_n=24+\frac{16q^{-2n}}{(1+q^{-n})^4}.
\]
As a rational function of $x=q^n$ this is
$24+16x^2/(x+1)^4$, which has a genuine pole at $x=-1$ and is not
Laurent. The criterion just proved excludes every finite exponential
polynomial, with arbitrary complex coefficients and Jordan-polynomial
terms allowed.

Finally let $T$ be any fixed endomorphism on a finite-dimensional
cohomology space with invertible Frobenius $F$. Its trace series is
\[
 \sum_{n\ge0}\Tr(TF^n)Z^n
 =\Tr\bigl(T(1-ZF)^{-1}\bigr),
\]
a rational function, equivalently a finite exponential polynomial in
$n$. A finite alternating sum over degrees has the same property.
Commutation of $T$ with $F$ is not needed for this obstruction; for a
Weil endomorphism it holds and only ordinary power sums are needed.
This proves the finite-realization assertion.
\end{proof}

\paragraph{The exact comparison excluded.}
Suppose a proposed finite realization assigns fixed Weil endomorphisms
$E_1,E_2$ to $e_1,e_2$ and requires, for every word $w$ and every
extension degree,
\[
 \sum_i(-1)^i\Tr\bigl(w(E_1,E_2)F^n\mid H^i\bigr)
 =\left.\Tr_M(w(e_1,e_2))\right|_{v=q^{n/2}}.
\]
The word $e_1e_2$ contradicts \fid{F26}; even imposing no idempotency
on the $E_i$ would not help. At one fixed specialization
$v=q^{1/2}$ the algebraic matrices can certainly act on a constant
Weil vector space. The resulting trace against $F^n$ is different
from substituting $v=q^{n/2}$ afresh at each extension degree.

The denominator issue is already visible in the finite Tate ring:
with $\mathbb L=v^2$, $D=\mathbb L+2+\mathbb L^{-1}$ is not a unit
of $\Q[\mathbb L,\mathbb L^{-1}]$. A bounded finite virtual sum of
Tate lines cannot have class $D^{-1}$. In contrast the actual
unnormalized trace $24D^2+16$ is a positive Laurent polynomial and
has a finite graded Tate realization. Individual matrix entries can
likewise be lifted to signed graded classes. Such entrywise lifts do
not construct compatible chain maps, multiplication functors or
categorical splittings of the normalized idempotents.

Localized motive rings, stacks with unbounded cohomology and completed
categories can accommodate rational expressions; they are not excluded
by this theorem. They require an additional finite extraction, or a
trace theorem beyond \fid{F02}/\fid{F24}, for the particular non-Laurent
family above. Indeed its convergent expansion is
\[
 h_n=24+16\sum_{j\ge0}(-1)^j\binom{j+3}{3}q^{-(j+2)n}.
\]
The independent analytic proof in the audit shows genuine generating
series poles at every $q^{j+2}$; the recurrence proof above avoids any
convergence assumption. Clearing denominators before taking traces is
another possible dictionary, but its compatibility with the designated
coefficient and operations is a separate claim.


\subsection{Fixed operators and extension-dependent quantum grading}
Clearing the normalization denominators in \fid{F26} removes
the obstruction for an individual word. It does not produce
one finite Frobenius space supporting all words under the same
extension-dependent substitution.

\begin{finding}[\fid{F42}: a common finite spectrum cannot encode all words]\label{f:commonwords}
On the actual rank-three tensor module of \fid{F21}, put
$D=v^2+2+v^{-2}$. For every $k\ge1$,
\[
 \Tr(P_1^k)=40D^k.
\]
Fix $q>1$ and set $D_n=q^n+2+q^{-n}$.
There is no finite-dimensional complex vector space $V$,
fixed invertible $F\in\End(V)$, and endomorphisms $T_k$ for
an unbounded set of positive integers $k$, such that
\[
 \Tr(T_kF^n)=40D_n^k
\]
for all sufficiently large $n$. The threshold may depend on $k$.
The assertion permits noncommuting $T_k,F$ and Jordan blocks,
and remains valid for finite Euler supertraces. It also remains
valid after multiplying each right-hand side by a nonzero scalar
$a_k$ and a single nonzero exponential factor $\beta_k^n$.

When $q$ is a prime power, the entire actual tensor family at the single
parameter $v_0=q^{1/2}$ has an exact action-compatible realization
by finite Weil objects with identity Frobenius. Its normalized
generators are genuine idempotent endomorphisms. Moreover each
individual sequence $40D_n^k$ has a positive finite Tate
coefficient object, with that object allowed to depend on $k$.
These constructions do not supply the opposite-sign path
comparison required in \fid{F24}.
\end{finding}

\begin{proof}
Use the consistent normalization of \fid{F19}: $c_i^2=\delta c_i$,
$\delta=v+v^{-1}$, and
\[
 P_i=c_i\otimes c_i+(c_i-\delta)\otimes(c_i-\delta).
\]
The two terms are orthogonal and their squares are $D$ times
themselves, so $P_i^2=DP_i$. This uses the symbolic definition
in GCT IV, equation~(93), with the section-3.1 convention
\cite[\S3.1 and Definition~11.1, equation~(93)]{IV},
rather than silently using the inconsistent section-11
assignments already recorded in the source audit.
In the decomposition of \fid{F21}, $P_1$ acts by $D$ on
the scalar positive block of multiplicity 20 and by zero on
the scalar negative block of multiplicity 4. Each of the
two-dimensional blocks has rank one for $P_1$ and trace $D$,
with multiplicities 16 and 4. Hence
\[
 \Tr(P_1^k)=40D^k\qquad(k\ge1).
\]
Both substitutions $v=q^{n/2}$ and $v=q^{-n/2}$ give
\begin{equation}\label{eq:wordpowers}
 40D_n^k
   =40\sum_{j=0}^{2k}\binom{2k}{j}q^{(k-j)n}.
\end{equation}
The $2k+1$ roots in this sum are distinct, and every coefficient
is nonzero.

For fixed $F,T_k$, the generating function beginning at $n=0$ is
\[
 \sum_{n\ge0}\Tr(T_kF^n)T^n
   =\Tr\bigl(T_k(1-TF)^{-1}\bigr).
\]
All its poles are among the inverses of the finitely many
eigenvalues of $F$. A Jordan block can increase a pole order
but cannot add a pole location. Removing finitely many initial
terms changes only a polynomial. On the other hand the
generating function of \eqref{eq:wordpowers} is
\[
 40\sum_{j=0}^{2k}\binom{2k}{j}
      \frac{1}{1-q^{k-j}T}.
\]
Every one of its $2k+1$ poles is genuine: only one summand has
that pole, and its coefficient is nonzero. Thus a common $F$
would require at least $2k+1$ distinct eigenvalues for unbounded
$k$, which is impossible.
For a finite graded object use the union of the even and odd
spectra; for a bounded complex use its finite-dimensional
cohomology. Multiplication by $a_k\beta_k^n$ preserves the
number of distinct nonzero roots, now
$\beta_kq^{k-j}$, proving the normalization variant.
This proof even permits unrelated $T_k$; requiring
$T_k=P^k$ or additional tensor relations cannot evade it.

We give the positive fixed-parameter construction with its
action comparison stated explicitly. Let $q$ be a prime power,
$\ell\nmid q$, and let a finite extension $E/\Q_\ell$
contain $v_0=q^{1/2}$.
Write
\[
 M_{r,q}=(V_q\otimes W_q)^{\otimes r}
\]
for the actual tensor module with the two ordinary Hecke
actions at $v_0$, and let $A_{r,q}$ be the image algebra of its
specified nonstandard generators. An action comparison consists
of finite Weil objects $H_r$, fiber identifications
$\omega(H_r)\simeq M_{r,q}$, and right $A_{r,q}$-actions
by Weil endomorphisms such that:
\begin{enumerate}
\item every named $P_i$ is intertwined by the fiber
identification;
\item the tensor identifications $H_r\otimes H_s\simeq H_{r+s}$
intertwine the local generators on disjoint consecutive blocks,
with the usual tensor associativity coherence.
\end{enumerate}
Extension of the finite field raises Frobenius to its $n$th
power while keeping $v_0$ and these operators fixed.

Give $V_q,W_q$ identity Frobenius and use the induced identity
Frobenius on their tensor powers. Every actual matrix is then
a Weil endomorphism. The local Hecke formula in
GCT IV, \S6.1 \cite[\S6.1, p.~33]{IV}, gives the actual operators, and locality proves
the disjoint-block compatibility. All their algebra relations
are unchanged. This proves the stated action comparison for
the whole actual image-algebra family.
For fixed dimensions two these are the two-row quotient actions;
no faithfulness assertion for the full abstract algebra at
every rank is required. The quantum Hecke operators are extra
morphisms, not an identification with the symmetry maps of the
ordinary Weil category.

Since $D_0=(v_0+v_0^{-1})^2\ne0$, scalar division defines
$e_i=P_i/D_0$ and gives $e_i^2=e_i$. Their images are genuine
Weil subobjects. This divides by a coefficient-field scalar;
it does not tensor-invert the object of \fid{F38}.
The resulting extension trace is
\[
 \Tr\bigl(w(P_i(v_0))F_r^n\bigr),
\]
with fixed operators. For identity Frobenius it is constant
in $n$, not the formal word trace retuned at $q^{n/2}$.
In particular the normalized mixed-word trace of \fid{F26}
is the constant $24+16/D_0^2$. This is an endomorphism trace,
not automatically the Euler multiplicity of an object;
$e_1e_2$ need not be idempotent.

There is also a valid coefficient-object construction for
\eqref{eq:wordpowers}. On a point let
\[
 \mathcal D=E(-1)\oplus E^{\oplus2}\oplus E(1),\qquad
 \mathcal C_k=E^{\oplus40}\otimes\mathcal D^{\otimes k}.
\]
Geometric Frobenius has trace $D_n$ on $\mathcal D$ and
$40D_n^k$ on $\mathcal C_k$ at every extension.
The canonical tensor identification
$\mathcal C_{k+1}\simeq\mathcal D\otimes\mathcal C_k$
realizes multiplication by this coefficient.
Its dimension is $40\cdot4^k$ and its weights range from
$-2k$ to $2k$, so positivity here does not assert purity.

The one-generator quadratic relation itself has a coherent
functor model. On the product of 64 finite Weil-representation
categories, let $\mathcal E_{40}$ project onto 40 chosen
components and set
$\mathcal P=\mathcal D\otimes\mathcal E_{40}$. Then
\[
 \mathcal P\circ\mathcal P
    \simeq\mathcal D\otimes\mathcal P
\]
by tensor associativity and
$\mathcal E_{40}^2=\mathcal E_{40}$, compatibly with further
powers. This models the split action of one generator and its
quadratic relation. It does not identify a simultaneous model
for the second generator, the original integral tensor lattice,
or the designated nonstandard bases.

Finally the action comparison alone cannot imply the main
sign hypothesis. Take a constant rank-$N$ sheaf with identity
Frobenius on affine $d$-space, with any of the preceding
fixed-parameter actions on its fibers. Its tensor and named
operator relations remain exact. Its pointwise trace is $N$,
so the pointwise hypothesis of \fid{F24} holds with $a=0$
and $B=N$, while its total compact trace is
\[
 M_n=Nq^{dn}.
\]
For $d>0$ the visible root $q^d$ exceeds $q^{d/2}$.
The proved \fid{F24} therefore rules out the required
opposite-sign path comparison with sublinear boundary budget
on an unbounded even set for this model. This proves the
claimed insufficiency with an explicit action-compatible
example.
\end{proof}

Thus a common finite operator realization, a fixed-parameter
action, and word-dependent coefficient objects impose different
comparison requirements. The latter objects may grow with the
word; the common-space obstruction does not exclude them.
For the main route one must additionally identify the actual
path complex and its tensor or idempotent construction by a
Weil-compatible comparison, with the required odd shift and a
controlled cone. Equality of fixed-operator traces alone does
not provide that comparison. Nor does fixed-$q$ numerical
positivity imply coefficientwise positivity in a formal
variable: $x^2-x+1$ is positive for every real $x$ but has a
negative coefficient.


\subsection{An actual stack realization of the normalized word trace}
The modern stack formalism changes the finiteness premise of
\fid{F26}. The primary inputs used here are Sun's trace formula
\cite[Theorem~1.1, p.~49; Definition~4.1, p.~84; Theorem~4.2,
p.~85]{SunStacks}, his explicit classifying-stack calculation
\cite[Theorem~4.8 and \S4.8.2, pp.~89--90; Example~7.4,
p.~109]{SunStacks}, and the published Laszlo--Olsson duality and
K\"unneth formulas \cite[\S9.1, Lemmas~9.1.1--9.1.2, p.~198;
Theorem~11.0.14, p.~204]{LO2}. These statements and their applicable
hypotheses were inspected; their full general proofs are cited inputs.

\begin{finding}[\fid{F35}: a convergent stack trace with an unbounded dual]\label{f:stacktrace}
Over $\F_q$ let
\[
 \mathcal Y=[(\A^6\times\mathbb G_m^8)/(\GL_2)^4]
   =\A^6\times\mathbb G_m^8\times(B\GL_2)^4,
\]
with trivial action, and let $\mathcal X$ be the disjoint union
of 24 rational points and 16 copies of $\mathcal Y$.
Then $\mathcal X$ is a finite-type Artin stack with affine stabilizers,
and for every $n\ge1$, with $Q=q^n$,
\begin{equation}\label{eq:stackwordtrace}
 \#\mathcal X(\F_Q)
 =24+\frac{16Q^2}{(Q+1)^4}
 =\sum_i(-1)^i\Tr(F^n\mid H_c^i(\mathcal X_{\bar{\F}_q},\Q_\ell)).
\end{equation}
Here $\#$ means groupoid mass and the series is absolutely convergent.
This realizes the complete scalar sequence of \fid{F26}.

The component $\mathcal Y$ is smooth of dimension $-2$.
Its compact cohomology is unbounded below and satisfies actual duality
\[
 H_c^{-4-i}(\mathcal Y)\simeq H^i(\mathcal Y)^\vee(2).
\]
Its visible compact spectrum is bounded and admissible as in
\fid{F32}; the corresponding ordinary dual spectrum is unbounded.
Thus actual duality holds, while the separate dual upper-bound
hypothesis of \fid{F17}/\fid{F32} fails in this model.
No nonstandard generator action or simultaneous word comparison
is asserted.
\end{finding}

\begin{proof}
Put $U=\A^6\times\mathbb G_m^8$ and $G=(\GL_2)^4$.
The atlas $U\to[U/G]$ is smooth of relative dimension 16, whereas
$\dim U=14$. This gives the stated dimension and smoothness of
$\mathcal Y$. The full $\mathcal X$ is not equidimensional: its
components have dimensions zero and minus two, and $\dim\mathcal X=0$.

A $\GL_2$-torsor over a field is the frame torsor of a rank-two
vector space, hence is trivial. For the trivial action each
$U(\F_Q)$ point therefore contributes $|G(\F_Q)|^{-1}$ to the mass.
Counting linearly independent columns gives
\[
 |\GL_2(\F_Q)|=(Q^2-1)(Q^2-Q)=Q(Q-1)^2(Q+1).
\]
Consequently
\[
 \#\mathcal Y(\F_Q)
 =\frac{Q^6(Q-1)^8}{Q^4(Q-1)^8(Q+1)^4}
 =\frac{Q^2}{(Q+1)^4}.
\]
Sun's Theorem~1.1 gives the asserted convergent trace for the constant
sheaf. Equivalently Theorem~4.2 applies: the constant sheaf is
bounded, geometrically stratifiable, pure of weight zero, and
convergent in Definition~4.1. The structural map is finite type.
The finite-field base also satisfies the excellence, invertibility
of $\ell$, and finite cohomological-dimension hypotheses of the
cited Laszlo--Olsson formalism.

We now compute the actual cohomology, suppressing the geometric
base change and coefficient field in the notation. Define
\[
 \sum_{s\ge0}a_s z^s=(1-z)^{-4}(1-z^2)^{-4}.
\]
For one classifying factor the known polynomial presentation is
$H^*(B\GL_2)=\Q_\ell[c_1,c_2]$, with $\deg c_i=2i$ and
$c_i$ a rational Chern class spanning $\Q_\ell(-i)$.
Sun's Example~7.4 explicitly gives
\[
 H_c^{-8-2s}(B\GL_2)
 =\Q_\ell(s+4)^{\oplus(\lfloor s/2\rfloor+1)},\qquad s\ge0.
\]
The ordinary cohomology of $\mathbb G_m^8$ is the exterior
algebra on eight copies of $\Q_\ell(-1)$ in degree one, while
$\A^6$ has only the unit in ordinary cohomology. K\"unneth gives
\begin{equation}\label{eq:stackordinarygroups}
 H^i(\mathcal Y)=
 \bigoplus_{\substack{0\le b\le8,\ s\ge0\\i=b+2s}}
 \Q_\ell(-(b+s))^{\oplus\binom8b a_s}.
\end{equation}
For a direct compact-support calculation, the affine factor has
$H_c^{12}(\A^6)=\Q_\ell(-6)$, the torus has
$H_c^{16-b}(\mathbb G_m^8)=\Q_\ell(-(8-b))^{\oplus\binom8b}$,
and the four classifying factors give
$H_c^{-32-2s}((B\GL_2)^4)=\Q_\ell(s+16)^{\oplus a_s}$.
The compact-support K\"unneth formula therefore yields
\begin{equation}\label{eq:stackcompactgroups}
 H_c^j(\mathcal Y)=
 \bigoplus_{\substack{0\le b\le8,\ s\ge0\\j=-4-b-2s}}
 \Q_\ell(b+s+2)^{\oplus\binom8b a_s}.
\end{equation}
The sum in each degree is finite. Since $a_s>0$ for every $s$,
these groups are unbounded below. The indicated weight is
$-2(b+s+2)=j-b$, consistent with Sun's affine-stabilizer bound
\cite[Theorem~1.4, p.~50]{SunStacks}. A fixed degree can be mixed:
$H_c^{-6}=\Q_\ell(3)^4\oplus\Q_\ell(4)^{28}$.

For $x=q^{-n}$ the sign in \eqref{eq:stackcompactgroups} is
$(-1)^b$, and the Euler series becomes
\[
 x^2\frac{(1-x)^8}{(1-x)^4(1-x^2)^4}
 =\frac{x^2}{(1+x)^4}.
\]
This grouping is legitimate: for every real $t>0$ the ungrouped
absolute eigenvalue moment is
\[
 \sum_{b,s}\binom8b a_s q^{-t(b+s+2)}
 =\frac{q^{-2t}(1+q^{-t})^4}{(1-q^{-t})^8}<\infty.
\]
The visible coefficient of $q^{-(j+2)}$ is consequently
\[
 \sum_{b=0}^{\min(8,j)}(-1)^b\binom8b a_{j-b}
   =(-1)^j\binom{j+3}{3}.
\]
For instance at $q^{-3}$ the actual degree-$-6$ multiplicity four
and degree-$-5$ multiplicity eight give the signed coefficient
$4-8=-4$. The full $\mathcal X$ therefore has exactly the
admissible visible spectrum
\[
 (1,24),\qquad
 \left(q^{-(j+2)},16(-1)^j\binom{j+3}{3}\right),\quad j\ge0.
\]
Its support and absolute moments meet \fid{F32}'s hypotheses, so
that finding's uniqueness applies in its stated convergence class.

Finally the smooth dualizing complex of $\mathcal Y$ is
$\Q_\ell(-2)[-4]$. The cited duality formula gives
$H_c^{-4-i}=H^i{}^\vee(2)$, also verified term by term by
\eqref{eq:stackordinarygroups}--\eqref{eq:stackcompactgroups}.
A compact root $\alpha=q^{-(b+s+2)}$ is paired with the
ordinary root $q^{-2}/\alpha=q^{b+s}$. After Euler grouping
the latter roots are all $q^j$, with nonzero coefficients
$(-1)^j\binom{j+3}{3}$. Their support is unbounded, and the
ordinary geometric-Frobenius trace terms do not tend to zero.
Thus this dual is outside the independently bounded visible-spectrum
class. For the nonequidimensional $\mathcal X$, Verdier duality
is computed componentwise; no common dimension-minus-two pairing
is imposed on all its components.
\end{proof}

This is a uniform geometric realization of one scalar word-trace
sequence, with unbounded cohomology explaining its compatibility
with \fid{F26}. It supplies neither operators for $e_1,e_2$ nor
their composition law, and no tensor-path sign or nonstandard basis
follows from its groupoid mass. Its failure of the independent dual
bound is a proved property of this model, not an exclusion of every
stack or completed realization.

\subsection{Finite weight windows and a fixed-parameter action}
\begin{finding}[\fid{F41}: finite windows lose the extension-wise normalization]\label{f:finitewindow}
Every graded Weil subquotient of the compact cohomology of the
stack $\mathcal Y$ in \fid{F35}, restricted to a fixed bounded
weight interval, has finite total dimension and bounded
cohomological support. No fixed endomorphism on such an object
has the alternating all-extension trace
\[
 h_n=24+\frac{16q^{2n}}{(q^n+1)^4}
\]
of the actual normalized word in \fid{F26}.

Nevertheless an actual 64-dimensional subquotient of a Tate
twist of $H_c^{-12}(\mathcal Y)$ supports, after choosing an
identification, the entire actual normalized tensor action
at the single parameter $v_0=q^{1/2}$ or $q^{-1/2}$.
This is a transport construction, not a restriction of an action
already constructed on the stack.
\end{finding}
\begin{proof}
The actual groups computed in \fid{F35} are
\[
 H_c^j(\mathcal Y)=
 \bigoplus_{\substack{0\le b\le8,\ s\ge0\\j=-4-b-2s}}
     \Q_\ell(b+s+2)^{\oplus\binom8b a_s},
 \qquad
 \sum_{s\ge0}a_s z^s=(1-z)^{-4}(1-z^2)^{-4}.
\]
Set $r=b+s+2$. Each summand has weight $-2r$ and scalar
Frobenius $q^{-r}$. A bounded weight interval permits only
finitely many integers $r$, and for each one there are only
finitely many pairs $(b,s)$, since $0\le b\le8$ and
$s=r-b-2\ge0$. Every multiplicity is finite. Also
$j=-2r+b$ takes only finitely many values. This proves the
finiteness statement and the same statement for subquotients.

The $q^{-r}$-eigenspaces are canonical in each degree.
Consequently a degree-preserving Frobenius-commuting action,
if supplied, preserves these windows; any invariant
subquotient retains it. This conditional statement does not
supply that action or make every chosen subspace invariant.
For any fixed endomorphism $T$ on a finite graded window,
its alternating trace with $F^n$ has the form
\[
 \sum_r c_r q^{-rn}
\]
with finitely many coefficients. These are the alternating
traces of the diagonal blocks of $T$; even commutation with
Frobenius is unnecessary for this trace formula.
Put $x_n=q^{-n}$. If it agreed with $h_n$ for every
sufficiently large $n$, the Laurent polynomial $\sum_r c_rx^r$
would agree at infinitely many points with
\[
 24+\frac{16x^2}{(1+x)^4}.
\]
Clearing denominators gives a polynomial with infinitely many
roots, hence an identity. But the latter rational function has
a genuine pole at $x=-1$, while a Laurent polynomial does not.
This contradiction persists after any Tate twist. The more
general finite Weil-complex obstruction is \fid{F26}.

For the positive fixed-parameter construction, degree $-12$
and weight $-12$ require
\[
 b+2s=8,\qquad b+s+2=6,
\]
so $b=0,s=4$. The corresponding actual summand is
\[
 \Q_\ell(6)^{85}\subset H_c^{-12}(\mathcal Y),
 \qquad a_4=35+40+10=85.
\]
The three terms count monomials of respective total degrees
$(4,0),(2,1),(0,2)$ in four variables of weight one and four
variables of weight two. Twist this summand by $(-6)$ and
choose a 64-dimensional subspace $W$. It is an effective
pure weight-zero Weil object with Frobenius exactly one.

Choose a finite coefficient extension $E/\Q_\ell$ containing
$v_0=q^{1/2}$, or its inverse. The actual 64-dimensional tensor
representation in \fid{F26} is defined at this parameter because
$D(q)=q+2+q^{-1}\ne0$. Choose an $E$-linear identification
with $W\otimes E$ and transport its specified normalized
generators. These actual endomorphisms obey all the specialized
relations and commute with scalar Frobenius, preserving every
word action at this one parameter.

For $T=e_1e_2$ their extension traces are consequently
$\Tr(TF^n)=h_1$ for every $n$, rather than $h_n$.
Before the twist they would be $q^{-6n}h_1$.
Raising Frobenius to its $n$th power does not reevaluate fixed
matrix coefficients at $q^{n/2}$.
\end{proof}

Thus a finite effective action model at fixed $q$ exists,
whereas retaining the full extension-wise normalized trace
requires different data. The construction gives no canonical
stack correspondence, restriction of a pre-existing stack
action, or nonstandard basis comparison. Its finite-window
obstruction is conditional on the explicitly stated
all-extension trace requirement.


\subsection{An inverse trace need not be an invertible object}
The stack construction in \fid{F35} permits convergent reciprocal
traces. The next calculation identifies the additional obstruction
to interpreting such a reciprocal as a tensor inverse in the usual
derived Weil realization.

\begin{finding}[\fid{F38}: trace inversion and tensor inversion differ]\label{f:traceinverse}
Over $\F_q$, set
\[
 \mathcal Z=\A^1\times\mathbb G_m^2\times B\GL_2,
 \qquad \mathcal W=\mathbb P^1\times\mathcal Z.
\]
For every $Q=q^n$ their groupoid masses are
\[
 \#\mathcal Z(\F_Q)=\frac1{Q+1},\qquad
 \#\mathcal W(\F_Q)=1.
\]
The corresponding compact Euler traces converge absolutely.
Nevertheless $R\Gamma_c(\mathcal W)$ is not the unit: it has
nonzero cohomology in arbitrarily negative degrees, including
$H_c^{-1}(\mathcal W)=\Q_\ell(1)^2$.

In the usual unbounded derived category of Weil vector spaces,
with the algebraic tensor product,
\[
 H=R\Gamma(\mathbb P^1_{\bar{\F}_q},\Q_\ell)
   =\Q_\ell\oplus\Q_\ell(-1)[-2]
\]
has no tensor inverse. More generally, no tensor-invertible object
in this category has all-extension Euler trace
$D(q^n)=q^n+2+q^{-n}$.
Thus a strong monoidal realization with this specified image of
the projective-line object cannot extend through a tensor localization
making that object invertible. No faithfulness assumption is needed.
\end{finding}

\begin{proof}
Both stacks are smooth, of dimensions $-1$ and zero, and have
affine stabilizers. The torsor and column-count calculation in
\fid{F35} gives
\[
 \#\mathcal Z(\F_Q)
 =\frac{Q(Q-1)^2}{|\GL_2(\F_Q)|}
 =\frac1{Q+1}.
\]
Multiplication by $\#\mathbb P^1(\F_Q)=Q+1$ proves the mass
formula for $\mathcal W$.

Let $c_s=\lfloor s/2\rfloor+1$. Using the classifying-stack
calculation \cite[Example~7.4, p.~109]{SunStacks} and the compact
K\"unneth formula \cite[Theorem~11.0.14, p.~204]{LO2}, we obtain
\begin{equation}\label{eq:reciprocalstackgroups}
 H_c^j(\mathcal Z)=
 \bigoplus_{\substack{0\le b\le2,\ s\ge0\\j=-2-b-2s}}
 \Q_\ell(b+s+1)^{\oplus\binom2b c_s}.
\end{equation}
Indeed $\A^1$ contributes degree two and twist $-1$;
$\mathbb G_m^2$ contributes degree $4-b$ and twist $-(2-b)$;
and $B\GL_2$ contributes degree $-8-2s$ and twist $s+4$.
The sum of these degrees and twists is precisely the displayed
formula. The finite-field and finite-type hypotheses are the
same as those verified in \fid{F35}.

Put $K=R\Gamma_c(\mathcal Z)$. The projective bundle formula,
already used in \fid{F15}, gives the displayed decomposition of $H$.
Properness of $\mathbb P^1$ and K\"unneth give
\[
 R\Gamma_c(\mathcal W)=H\otimes K,\qquad
 H_c^j(\mathcal W)=H_c^j(\mathcal Z)
       \oplus H_c^{j-2}(\mathcal Z)(-1).
\]
Hence $H_c^0(\mathcal W)=\Q_\ell$ and
$H_c^{-1}(\mathcal W)=\Q_\ell(1)^2$.
Formula~\eqref{eq:reciprocalstackgroups} also makes the unbounded
negative support explicit. These are actual cohomology direct sums,
so their nonzero summands cannot cancel.

For $x=q^{-n}$ the compact Euler series of $K$ is
\[
 x\frac{(1-x)^2}{(1-x)(1-x^2)}
   =\frac{x}{1+x}
   =\sum_{j\ge0}(-1)^j x^{j+1}.
\]
Every positive real absolute eigenvalue moment converges: if $t>0$,
the ungrouped sum is
\[
 \sum_{b,s}\binom2b c_s q^{-t(b+s+1)}
 =\frac{q^{-t}(1+q^{-t})}{(1-q^{-t})^2}<\infty.
\]
Thus Sun's trace theorem applies, and this explicit convergence
also justifies all regrouping. Tensoring with the finite complex
$H$ preserves absolute convergence. Since its Euler trace is
$1+x^{-1}$, the Euler trace of $H\otimes K$ is exactly one.
After grouping, its only visible root is one with coefficient one;
the remaining actual cohomology cancels between degrees. The
uniqueness in \fid{F32} concerns that visible spectrum and is not
a faithfulness assertion for the trace.

For completeness, over a field every complex, including an
unbounded one, splits as its cohomology with zero differential
plus a contractible complex. Tensoring a contractible complex
remains contractible. With algebraic direct-sum totalization,
the resulting K\"unneth formula is
\[
 H^m(A\otimes B)=\bigoplus_{i+j=m}H^i(A)\otimes H^j(B).
\]
If $A\otimes B$ is the one-dimensional unit, both factors are
nonzero. Two independent total cohomology vectors in either
factor would give at least two in their tensor product.
Consequently each factor has one-dimensional total cohomology,
concentrated in a single degree. Forgetting the Weil action is
exact and conservative, so the same necessity holds for Weil
complexes. Truncation then identifies each factor with a
one-dimensional Weil representation in one degree. Conversely
such a representation, with its shift, has its dual with the
opposite shift as inverse.

In particular $H$ has no inverse. Directly, if $H^a(B)\ne0$,
then $H\otimes B$ has nonzero cohomology in both degrees $a$
and $a+2$. Strong monoidal functors preserve tensor inverses.
Therefore no strong monoidal realization sending a designated
object to this $H$ can factor through a localization in which
that object becomes invertible.

For the nonstandard normalization, under $v=Q^{\pm1/2}$ we have
\[
 D(Q)=(v+v^{-1})^2=Q^{-1}(Q+1)^2.
\]
The actual complexes
\[
 H_D=\Q_\ell(1)\otimes H^{\otimes2},\qquad
 K_D=\Q_\ell(-1)\otimes K^{\otimes2}
\]
have traces $D(Q)$ and $D(Q)^{-1}$, respectively. Their product
has trace one, but $H_D$ has total cohomology dimension four and
cannot have a tensor inverse. This conclusion does not depend
on choosing $H_D$: every invertible Weil complex has trace
$\pm\alpha^n$ for a single nonzero Frobenius scalar $\alpha$.
The minus sign cannot give the positive first and second values
of $D(q^n)$. The plus sign would require $D(q^2)=D(q)^2$, whereas
\[
 D(q)^2-D(q^2)=4(q+q^{-1}+1)>0.
\]
This proves the asserted all-extension obstruction.
\end{proof}

The statement permits formal Euler-class localizations, other
realizations, and completions with a different tensor product or
trace functional. It also permits multiplying an operator by the
numerical scalar $D(q)^{-1}$ at a fixed $q$. Scalar division does
not assert an object-level inverse or a reevaluation at every
$q^n$. The fixed realization and usual derived tensor product
are essential hypotheses of the obstruction.

\subsection{Actual pure classical multiplicity objects via geometric Satake}
\begin{finding}[\fid{F20}: Schur summands have pure multiplicity spaces; their cycle action need not be positive]\label{f:satake}
Let $G=\GL_r/\F_q$ be split, with $\ell\nmid q$, and use the symmetric geometric Satake category with $\overline\Q_\ell$ coefficients. For a dominant partition $\mu$ of length at most $r$, let $\IC_\mu$ be the natural unnormalized Weil intersection complex of its Schubert variety, perverse of weight
\[
 D_\mu=\dim\Gr_\mu=\langle2\rho,\mu\rangle
 =\sum_{i=1}^r(r+1-2i)\mu_i.
\]
For $\nu\vdash d$, a primitive Young idempotent $e_\nu\in\Q[S_d]$ acts, through the parity-corrected symmetric commutativity, on
$K=\IC_\mu^{\star d}$. Its image $P_{\nu,\mu}$ is a pure perverse Weil sheaf of weight $dD_\mu$. The canonical multiplicity Weil module
\[
 M_{\lambda;\nu,\mu}
 =\Hom_{\mathrm{Sat}_{\overline{\F}_q}}
 (\IC_\lambda,P_{\nu,\mu})
\]
is pure of weight $dD_\mu-D_\lambda$ and has dimension
\[
 \dim M_{\lambda;\nu,\mu}
 =[S_\lambda(\overline\Q_\ell^r):
   S_\nu(S_\mu(\overline\Q_\ell^r))].
\]
On the product Grassmannian for $\GL_m\times\GL_n$, apply the Schur idempotent for $\nu\vdash d$ to the $d$th convolution power of
$\IC_{(1),m}\boxtimes\IC_{(1),n}$. Its multiplicity module at
$\IC_{\lambda,m}\boxtimes\IC_{\mu,n}$ has dimension $g(\lambda,\mu,\nu)$ and is pure of weight
\[
 d(m+n-2)-D_{\lambda,m}-D_{\mu,n}.
\]
These are ordinary geometric realizations of the classical coefficients, including rectangular cases within the dimension bounds. They do not by themselves identify the nonstandard quantum bases.

Even in this construction the symmetric-group action on a geometric multiplicity space need not permute a basis. For $G=\GL_2$, $\mu=(1)$ and $d=2$, the determinant multiplicity line is $\overline\Q_\ell(-1)$ and the transposition acts by $-1$. Thus no basis of that line is an $S_2$-permutation basis.
\end{finding}
\begin{proof}
The literature inputs are geometric Satake, its finite-field convolution/symmetry, and ordinary weight theory. We use the published version of Richarz \cite[Theorem~3.1, p.~214; proof pp.~226--227; Proposition~4.1, p.~227; Theorem~5.1, p.~230; Lemma~5.4 and its proof, p.~232]{Richarz}. The tensor equivalence is also Mirkovi\'c--Vilonen, Theorem~7.3 \cite{MV}. These theorems are cited, not re-proved here.

The twisted product of $d$ Schubert varieties is projective and is smooth-locally the product for purposes of its intersection complex. Its natural Weil IC is normalized to restrict to $\overline\Q_\ell[dD_\mu]$ on the dense twisted orbit. With that normalization, rather than only a geometric IC isomorphism, it is the twisted external product of the $\IC_\mu$ and is pure of weight $dD_\mu$; this is the IC identification in the proof of Richarz's Lemma~5.4, together with BBD~5.3.2, second proof (which covers a nonaffine open immersion). Its proper convolution pushforward is pure of that same complex weight (BBD, Remark~5.4.9) and is perverse (Richarz, Theorem~3.1). This proves the assertion about $K$.

The symmetry constructed over $\F_q$, with the component-parity sign specified in the cited proof, gives an $S_d$ action by Weil endomorphisms. In the abelian perverse category the idempotent splits; its image is a direct summand, hence remains pure by BBD~5.3.1. Under the symmetric monoidal Satake equivalence, $\IC_\mu$ corresponds to $S_\mu(\overline\Q_\ell^r)$ and $K$ to its $d$fold tensor power. In Schur--Weyl decomposition, the primitive idempotent $e_\nu$ has rank one on the Specht factor $[\nu]$ and vanishes on the other factors. Its image therefore corresponds to one copy of the Schur functor $S_\nu(S_\mu(\overline\Q_\ell^r))$. A central isotypic idempotent would instead give $\dim[\nu]$ copies; that is why the primitive condition matters.

Here is the Frobenius argument for the multiplicity weight. Geometric semisimplicity gives the canonical evaluation isomorphism
\[
 \bigoplus_\lambda\IC_\lambda\otimes
 \Hom(\IC_\lambda,\overline P_{\nu,\mu})
 \longrightarrow\overline P_{\nu,\mu}.
\]
It is the canonical isotypic decomposition, not an arbitrarily chosen decomposition. Every support and natural $\IC_\lambda$ is defined over the split base $\F_q$, so Frobenius preserves each summand. Give its Hom factor the conjugation Frobenius induced by the two fixed Weil structures. Evaluation is then Weil-equivariant, and the isotypic summand is a pure perverse Weil sheaf of weight $dD_\mu$. This descent/subquotient formalism is BBD~5.1.2; purity of the summand is BBD~5.3.1. On its own smooth open Schubert stratum inside its support, $\IC_\lambda$ is the constant sheaf shifted by $D_\lambda$. Restriction to this open preserves both weight inequalities because $j^*=j^!$ (BBD~5.1.14). The Hom module shifted by $D_\lambda$ is therefore pure of weight $dD_\mu$, which means the module itself has weight $dD_\mu-D_\lambda$ by the shifted convention of BBD~5.1.8. This proves the all-eigenvalue purity statement without assuming arithmetic semisimplicity. The tensor equivalence identifies its dimension with the required multiplicity.

For the product group, external product gives the simple object corresponding to $V\boxtimes W$. Its weight is $(m-1)+(n-1)$. The same Schur-idempotent construction corresponds to $S_\nu(V\otimes W)$ as a $\GL_m\times\GL_n$ representation. The definition of the Kronecker coefficient and the just-proved Hom-weight calculation give the stated formula.

Finally consider $\GL_2$ and two minuscule steps. The twisted convolution variety is a smooth surface. Over the central lattice $t\mathcal O^2$, its fiber consists of the intermediate lattices between $t\mathcal O^2$ and $\mathcal O^2$, hence of lines in $\mathcal O^2/t\mathcal O^2$; it is $\PP^1$. The contribution to the perverse degree-zero stalk is $H^2(\PP^1)=\overline\Q_\ell(-1)$. To identify it as the determinant multiplicity, Satake gives
\[
 V\otimes V=\Sym^2V\oplus\wedge^2V,
\]
with determinant multiplicity one. Equivalently the non-open relevant fiber supplies this unique one-dimensional summand; the IC of the open two-dimensional stratum has no perverse degree-zero stalk contribution at the boundary point by the intermediate-extension condition. The corrected symmetric commutativity corresponds to the ordinary tensor swap, which sends
$e_1\otimes e_2-e_2\otimes e_1$ to its negative. The sole nonzero basis vector of a one-dimensional permutation representation would have to be fixed, so no such basis exists. The symmetric projector kills this line, while the alternating projector keeps it.
\end{proof}
This is a constructive reduction of the ordinary-geometry gate for \emph{classical} coefficients. Its purity uses the existing Weil/BBD theorems; it is not a new application of \fid{F03}, and it does not provide the path-hypersurface identities for a nonstandard canonical coefficient. The remaining problem is more specific: describe the actual projector image, its graded basis and its compatibility with the designated quantum operations, together with effective labels when a counting rule is required. The $\GL_2$ example prevents replacing that task by unsigned orbit-counting on the unprojected component basis.

The construction also provides an exact Grothendieck-framework dictionary. Convolution is a proper constructible pushforward; the Young idempotent is an actual projector in a characteristic-zero symmetric tensor category; its image is an effective object whose $K_0$ class is the Schur operation. The weights record the unnormalized IC shifts, so no half-Tate convention is hidden. The motivating language of projectors occurs in cote~16, scans~35, 37 and~39; the actual theorem here uses the modern Satake equivalence and the precisely cited SGA/BBD weight formalism, not an attribution to those leaves. The explicit symmetric-group action is the additional datum beyond virtual $\lambda$ identities.

\paragraph{\fid{F20}, motivic refinement: scalar Tate Frobenius.}
In the rational motivic Satake category over $\F_q$, these same Schur
projector images have coefficient motives
\[
 \mathbf1(-h)^{\oplus a},\qquad
 h=\frac{dD_\mu-D_\lambda}{2}
   =n(\lambda)-d\,n(\mu),\qquad
 n(\alpha)=\sum_i(i-1)\alpha_i,
\]
where $a=[S_\lambda V:S_\nu(S_\mu V)]$. For the product-group Kronecker
construction the coefficient is
\[
 \mathbf1\bigl(-n(\lambda)-n(\mu)\bigr)^{\oplus g(\lambda,\mu,\nu)}.
\]
Consequently the natural Frobenius in the \emph{specified} Weil
multiplicity spaces is the scalar $q^h$ and has no Jordan part.

\begin{proof}
The input is Richarz--Scholbach \cite[Theorem~4.1; Proposition~4.8;
Proposition~5.9; Definition~6.2; Lemma~6.3; Corollary~6.4;
Lemma~6.5]{RS}. These provide pure intersection motives, weight-exact
convolution, symmetric constraints compatible with the corrected
\(\ell\)-adic realization, and a semisimple abelian compact Satake
category generated by the $\IC_\lambda(n)$ and stable under convolution.
Thus the primitive rational Young idempotent acts motivically and its
image is pure of weight $w=dD_\mu$. Its simple summands have the form
$\IC_\lambda(n)$, of weight $D_\lambda-2n$. Purity forces
$n=(D_\lambda-w)/2=-h$. Compatibility with geometric Satake identifies
their multiplicity with $a$. Restriction of the corresponding isotypic
summand to its open stratum and removal of the shift $D_\lambda$ gives
exactly $\mathbf1(-h)^{\oplus a}$; Lemma~6.3 gives the coefficient-motive
Hom identification. Since $|\lambda|=d|\mu|$ and
$D_\alpha=(r-1)|\alpha|-2n(\alpha)$, the stated integer formula for
$h$ follows. The product-group proof replaces $D_\lambda$ by
$D_{\lambda,m}+D_{\mu,n}$ and $w$ by $d(m+n-2)$, giving
$h=n(\lambda)+n(\mu)$. Finally the realization of $\mathbf1(-h)$ is
$\overline\Q_\ell(-h)$, whose geometric Frobenius is exactly $q^h$.
\end{proof}
Here the coefficient motive is not the untwisted vector space
$\Hom_{\rm mot}(\IC_\lambda,P)$, which can vanish when $h\ne0$.
The relevant multiplicity vector space is
$\Hom_{\rm mot}(\IC_\lambda(-h),P)$; tensoring it with the Tate line
recovers the Weil multiplicity module. This is stronger than purity for
these natural classical objects, and uses the modified dual-group
weight coordinate of \cite[Theorem~6.8]{RS}. It neither chooses a basis
of the coefficient vector space nor computes its rank efficiently.
Arithmetic semisimplicity for arbitrary pure Weil modules is still
false, as in \fid{F05}.

\subsection{A precise positive rule for the Schur-projector gate}
\begin{finding}[\fid{F23}: signed orbits suffice, but need not exist in pure Satake spaces]\label{f:signedorbits}
Let $M$ be a characteristic-zero $S_d$-module with a specified basis
$B$ permuted up to signs. For a label orbit $O$, choose $b_O$, let
$H_O$ be its label stabilizer, and define
$\chi_O(h)$ by $hb_O=\chi_O(h)b_O$. Then
\[
 M\simeq\bigoplus_O\operatorname{Ind}_{H_O}^{S_d}\chi_O,
\]
and
\[
 \dim M^{S_d}=\#\{O:\chi_O=1\},\qquad
 \dim\Hom_{S_d}(\mathrm{sgn},M)
 =\#\{O:\chi_O=\mathrm{sgn}|_{H_O}\}.
\]
For a primitive Young idempotent $e_\nu$,
\begin{equation}\label{eq:orbitprojector}
 \dim e_\nu M
 =\sum_O\dim\Hom_{H_O}(\Res^{S_d}_{H_O}[\nu],\chi_O).
\end{equation}
If every $H_O$ is a Young subgroup $\prod_i S_{a_i}$ and $\chi_O$ is a
product of trivial or sign characters on its blocks, this last sum
counts chains from $\varnothing$ to $\nu$: at block $i$ add a horizontal
$a_i$-strip for the trivial character, or a vertical $a_i$-strip for sign.

These are actual positive formulas for the corresponding Schur
multiplicity whenever the specified signed-basis data exist on
\[
 M_{\lambda,\mu,d}
 =\Hom_{\GL(V)}(S_\lambda V,(S_\mu V)^{\otimes d}).
\]
They do not always exist: for $\GL_2$, $\mu=(1)$, $d=3$ and
$\lambda=(2,1)$ this multiplicity module admits no signed-permutation
basis, although its natural geometric realization is
$\overline\Q_\ell(-1)^2$ with scalar Frobenius.
\end{finding}
\begin{proof}
The sign cocycle implies that $\chi_O:H_O\to\{\pm1\}$ is a character.
The map $k[S_d]\otimes_{k[H_O]}k_{\chi_O}\to kO$ sending
$g\otimes1$ to $gb_O$ is well defined, equivariant and surjective.
Both spaces have dimension $[S_d:H_O]$, so it is an isomorphism.
The distinct orbit spans are a direct sum. Frobenius reciprocity gives
\eqref{eq:orbitprojector}, and the trivial/sign cases have a
one-dimensional Hom precisely when the indicated characters agree.
In $M=\bigoplus_\rho[\rho]\otimes W_\rho$, a primitive idempotent has
rank one on $[\nu]$, hence image dimension $\dim W_\nu$.
A central idempotent would have rank $\dim[\nu]\dim W_\nu$ instead.
When the stabilizer is the stated Young subgroup, the Frobenius
characteristic is a product of $h_{a_i}$ and $e_{a_i}$. The classical
Pieri rule \cite[Theorem~7.10, p.~16]{Pieri} expands this product by exactly
the horizontal/vertical chains just described. This proves the positive
chain formula, without presupposing the desired projector-image basis.

For the counterexample, Schur--Weyl duality gives
\[
 (\overline\Q_\ell^2)^{\otimes3}
 =S_{(3)}V\otimes[(3)]\oplus S_{(2,1)}V\otimes[(2,1)].
\]
Thus the indicated multiplicity is the faithful two-dimensional
standard representation of $S_3$. It is faithful because it is the
sum-zero subspace of the permutation representation on three letters,
whose complementary all-ones line is fixed. A signed-permutation basis
would embed $S_3$ into the group of $2\times2$ signed permutation
matrices, of order $2^2\cdot2!=8$. Lagrange's theorem excludes this.

Here is the geometric identification in full. The threefold minuscule
convolution variety for $\GL_2$ is an iterated $\PP^1$-bundle, smooth of
dimension three. Over the lattice
$t^2\F_q[[t]]e_1\oplus t\F_q[[t]]e_2$, the quotient has dimension
three and multiplication by $t$ is nilpotent of type $(2,1)$. The fiber
parametrizes its full stable flags $L\subset P$. Write
$I=\operatorname{im}N$, $K=\ker N$, of dimensions one and two.
Stability requires $L\subset K$, $I\subset P$, and $N(P)\subset L$.
If $P\ne K$, then $N(P)=I$ and $L=I$; if $P=K$, every line
$L\subset K$ works. Thus the reduced fiber is
\[
 \{L=I,\ P\supset I\}\ \cup\ \{P=K,\ L\subset K\}
 =\PP^1\cup_{\mathrm{pt}}\PP^1.
\]
Both components and the intersection are rational over $\F_q$.
Its $H^2$ is therefore $\overline\Q_\ell(-1)^2$. Since the source
dimension is three and $D_{(2,1)}=1$, proper base change identifies this
group with the degree $-1$ stalk of the pure perverse convolution.
The intermediate extension $\IC_{(3)}$ has no degree
$-\dim\Gr_{(2,1)}=-1$ stalk on this boundary stratum, by its strict
support condition. This stalk is exactly the $\IC_{(2,1)}$ multiplicity
module. Corrected Satake symmetry gives the standard $S_3$ action just
used, while Frobenius is scalar $q$. This proves every assertion about
the example.
\end{proof}
The obstruction is to signs $\{\pm1\}$, not arbitrary scalar-monomial
bases: the standard $S_3$ representation is induced from a nontrivial
character of its cyclic subgroup of order three. Pure Tate Frobenius
does not restrict the phases of this separate commuting group action.
The example also has classical positive tableau formulas; it refutes
only the asserted universal signed-orbit route.

\paragraph{Exact encoding conditions for a counting rule.}
For inputs of binary length $L$, suppose orbit representatives have
unique polynomial-length encodings with polynomial-time validity tests.
For the trivial/sign formula, suppose a polynomial-size stabilizer
generating set and its character values are computable in polynomial
time, including the cost of permutation/sign encoding when $d$ is
binary. The character test on generators makes the orbit label itself
a counted witness. For the mixed-Young formula, additionally require
polynomially many blocks and rows, with polynomial-bit block sizes and
parts, and a polynomial-time procedure computing the block sizes and
trivial/sign types from the orbit label. A witness consists of that label
and its chain of partitions.
All parts are bounded by $d$, so the witness is polynomial length.
For $\alpha\subset\beta$, a horizontal strip is checked by
$\beta_i\ge\alpha_i\ge\beta_{i+1}$; a vertical strip by
$\beta_i-\alpha_i\in\{0,1\}$. Sizes and partitions are checked by
integer arithmetic. These tests count exactly the displayed chains,
giving a conditional \(\#P\) rule. They do not provide a single
polytope, an Ehrhart family or PH1's oracle from purity alone.

\paragraph{A known binary family, with its exact small polytope.}
For $\dim V=2$, Hermite reciprocity and Clebsch--Gordan with swap parity
give the classical formula
\[
 \Sym^d(\Sym^2 V)\simeq\Sym^2(\Sym^dV)
       \simeq\bigoplus_{j=0}^{\lfloor d/2\rfloor}S_{(2d-2j,2j)}V.
\]
The exact sources are Goodman--Wallach, Theorem~4.1.18, p.~191,
and Donley--Kim, Proposition~3.1, equation~(3.4), p.~3, and
Proposition~8.1, equation~(8.2), p.~11 \cite{GW,DK}.
These known theorems are not re-proved here. The passage from $SL_2$
to $GL_2$ fixes total polynomial degree $2d$; in the summand
$S_{(2d-r,r)}$, a highest-weight vector is
$(x_1y_2-x_2y_1)^r x_1^{d-r}x_2^{d-r}$ and swap acts by $(-1)^r$.
For $\lambda=(2d-r,r)$, $0\le r\le d$, the specialized consequence is
\[
 a_\lambda(d[2])=\mathbf1_{2\mid r},\qquad
 a_{N\lambda}((Nd)[2])
 =\#(N P_{d,r}\cap\Z),\quad P_{d,r}=\{r/2\}.
\]
The point is given by $2x\le r$, $-2x\le-r$, so has an explicit
polynomial-bit description and separation test. This is a restricted
fixed-inner-partition Ehrhart formula, not general PH1. On odd-$r$
rays its ordinary support is unsaturated, but the zero odd residue
polynomial is permitted by VI Definition~1.2.4, so the VI saturation
index remains zero, exactly as distinguished in \fid{F18}.


\section{What is reached in the GCT cases requested}\label{sec:rectangles}
We keep three different notions separate: scalar nonnegativity of a representation multiplicity; a positive effective rule for it; and positivity of normalized canonical structure constants. The first is already part of semisimple representation theory and \fid{F12}. The others do not follow from it.

Write $n\times d=(d^n)$, $\lambda\vdash nd$, and
\[
 g_n(\lambda)=g(\lambda,(d^n),(d^n)),\qquad
 a_\lambda(d[n])=[S_\lambda V:\Sym^d(\Sym^nV)].
\]
Let $q^m_\lambda(d[n])$ denote the degree-$d$ multiplicity in the coordinate ring of the orbit closure of $X_{11}^{n-m}\operatorname{per}_m$, in IP's covariant convention. The known inequality is $0\le q^m_\lambda(d[n])\le a_\lambda(d[n])$ (IP, equation~(1.3)).

\subsection{Exact rectangular results already proved}
The attached Ikenmeyer--Panova paper \cite{IP}, matching arXiv 1512.03798v2, states:
\begin{itemize}
\item Theorem~1.4: $q^m_\lambda(d[n])>g_n(\lambda)$ implies $\operatorname{dc}(\operatorname{per}_m)>n$.
\item Theorem~1.6: if $n>3m^4$ and $g_n(\lambda)=0$, then $q^m_\lambda(d[n])=0$.
\item Theorem~1.8: $q^m_\lambda(d[n])>0$ requires $\ell(\lambda)\le m^2$ and $|\bar\lambda|\le md$, where $\bar\lambda$ removes the first row.
\item Theorem~1.10: put
$\mathcal X=\{(1),(1,1),(1^4),(1^6),(2,1),(3,1)\}$.
If $\bar\lambda\in\mathcal X$, then $a_\lambda(d[n])=0$. If $\bar\lambda\notin\mathcal X$ and there is $m\ge3$ with
\[
 \ell(\lambda)\le m^2,\quad |\bar\lambda|\le md,\quad
 d>3m^3,\quad n>3m^4,
\]
then $g_n(\lambda)>0$.
\end{itemize}
These are substantive existing positivity and no-occurrence results. They are not new consequences of our conditional sign theorem. Revised IV, Theorems~1.6/15.21(vii)--(viii), also supplies a positive highest-weight crystal rule when the two restricted factors have at most two rows (and the $\GL_4$ label at most four parts). That includes two-row rectangles, not arbitrary-height rectangles in the asymptotic regime.

\subsection{An obstruction for every nontrivial rectangular tensor square}
\begin{finding}[\fid{F28}: an isolated hook excludes all Young-stabilizer monomial models]\label{f:recthooks}
Let $M$ be a characteristic-zero $S_N$-module. If an internal hook
degree $1\le k\le N-2$ occurs in $M$, while its two neighboring
hook degrees do not, then $M$ has no scalar-monomial basis all of
whose label stabilizers are conjugate to Young subgroups.

Consequently, for \emph{every} rectangle $\rho=(a^b)$ with $a,b\ge2$,
the actual module $[\rho]\otimes[\rho]$ has no such basis, even when
arbitrary nonzero scalar phases are allowed. Thus the Young-stabilizer
version of \fid{F23} cannot describe the unprojected rectangular
Satake multiplicity module, including large-parameter rectangles.
This does not exclude non-Young stabilizers or other positive rules.
\end{finding}
\begin{proof}
For a Young subgroup $H=\prod_{i=1}^r S_{a_i}$ with nonempty blocks,
every one-dimensional character is a product of trivial and sign
characters. Set $s=\sum_{\text{sign blocks }i}(a_i-1)$; singleton
blocks contribute zero under either designation. If $U_N$ is the
standard sum-zero representation, then
\[
 \Res_H U_N=\bigoplus_i U_{a_i}\oplus k^{r-1},\qquad U_1=0.
\]
The classical identification $\wedge^jU_a=[(a-j,1^j)]$ follows here
from the cited Pieri rule: the oriented $j$-subset module
$\wedge^j k^a=\operatorname{Ind}_{S_j\times S_{a-j}}^{S_a}
(\mathrm{sgn}\boxtimes1)$ is the sum of the two adjacent hooks for $1\le j\le a-1$,
and $k^a=1\oplus U_a$ identifies their exterior powers inductively.
In the exterior power of the displayed restriction, trivial on block
$i$ forces degree zero, and sign forces degree $a_i-1$. The remaining
degree is chosen from the trivial $(r-1)$-space. Frobenius reciprocity
therefore gives the exact formula
\[
 [\operatorname{Ind}_H^{S_N}\chi:(N-j,1^j)]
       =\binom{r-1}{j-s}.
\]
Its positive hook support is the consecutive interval
$s,\ldots,s+r-1$, contained in $0,\ldots,N-1$ because
$0\le s\le N-r$.

A scalar-monomial basis decomposes into actual orbit modules
$\operatorname{Ind}_{H_O}^{S_N}\chi_O$, by the orbit-map proof in
\fid{F23}. Some summand must contain the given internal hook. If its
Young subgroup has $r\ge2$ blocks, the interval just computed includes
at least one neighboring hook, contradicting its absence from the
effective direct sum. If $r=1$, the orbit is the trivial or sign
representation of the whole group, whose hook degree is $0$ or
$N-1$, not internal. This proves the first assertion. There is no
cancellation in an effective representation decomposition.

The exact source input for rectangles is Ikenmeyer--Panova
\cite[Proposition~6.4, pp.~15--16]{IP}. With $m=\min(a,b)$ and $N=ab$,
it states
\begin{equation}\label{eq:recthookproduct}
 \sum_{j=0}^{N-1}g((N-j,1^j),\rho,\rho)t^j
       =\prod_{i=2}^{m}(1+t^{2i-1}).
\end{equation}
Transposition of both rectangular Specht factors justifies the
symmetric shorter-side convention. For $N\ge6$, degree three is
internal. The product has coefficient one at degree three and zero
at degrees two and four. Apply the isolated-hook assertion.
The only smaller rectangle is $(2,2)$ of size four. Its square is
$[(4)]\oplus[(2,2)]\oplus[(1^4)]$ by the already checked character
calculation, or Manivel's Theorem~1. Its hook degrees are $0,3$.
Every Young orbit with at least two blocks would force consecutive
hooks; all remaining whole-group orbits are one-dimensional and
cannot supply $[(2,2)]$. This proves the edge case as well.
\end{proof}

The actual multiplicity identification is explicit: for $\dim U=\dim V=b$,
Schur--Weyl duality gives
\[
 \Hom_{\GL(U)\times\GL(V)}
 (\det^aU\boxtimes\det^aV,(U\otimes V)^{\otimes ab})
 \simeq[(a^b)]\otimes[(a^b)]
\]
with its diagonal symmetric-group action. Its primitive Schur
projectors give precisely the rectangular Kronecker coefficients.
The obstruction thus concerns the requested actual family, not an
unrelated test representation. A virtual alternating induction formula
can cancel the forbidden hooks, but would not be the positive orbit
decomposition considered here.

The scope is sharp already at size four. The permutation module on
the three pairings of four letters is $1\oplus[(2,2)]$, so
$[(2,2)]^{\otimes2}$ is that module plus a sign line: it does have
a signed-permutation model. A pairing stabilizer is
$S_2\wr S_2$, of order eight, whereas a Young subgroup of $S_4$ has
order $24,6,4,2$ or $1$. The stabilizer is not Young.

\paragraph{A working ordinary monomial model in the known height-two family.}
Manivel \cite[Theorem~1, p.~1; \S\S2--3, pp.~2--4]{Manivel} proves
that $[(d,d)]^{\otimes2}$ is multiplicity free, with constituents the
partitions of $2d$ having at most four even parts, or exactly four odd
parts. In fact, for dimensions $(\dim U,\dim V,\dim W)=(2,2,4)$,
his proof identifies the highest-weight invariant algebra
\[
 \C[U\otimes V\otimes W]^{SL(U)\times SL(V)\times N_W}
 =\C[f_1,f_2,f_3,f_4]
\]
in the source's dual-polynomial convention. The four $W$-partition
labels are $(1^4)$, $(2)$, $(2,2)$ and $(2,2,2)$, with degrees $4,2,4,6$.
These are known source results, not consequences newly obtained from
purity. A monomial with label $\nu$, padded to four parts, necessarily
has exponents
\[
 a_1=\nu_4,\quad a_2=(\nu_1-\nu_2)/2,\quad
 a_3=(\nu_2-\nu_3)/2,\quad a_4=(\nu_3-\nu_4)/2.
\]
They are nonnegative by the partition inequalities, and integral
exactly when the padded parts have common parity. The two determinant
exponents are $d=2a_1+a_2+2a_3+3a_4=|\nu|/2$.
Thus this is an actual known monomial basis with a verified
multiplicity comparison, in a different representation-theoretic model.

Equivalently, for $\ell(\nu)\le4$, let
\[
 P_\nu=\{((\nu_1-\nu_4)/2,(\nu_2-\nu_4)/2,(\nu_3-\nu_4)/2)\}.
\]
Its coordinates are $(a_2+a_3+a_4,a_3+a_4,a_4)$, an integral
unimodular change from the last three exponents. The exact restricted
Ehrhart formula is
\[
 g(N\nu,(Nd,Nd),(Nd,Nd))=\#(NP_\nu\cap\Z^3).
\]
Use the empty polytope for length greater than four. The point has
denominator at most two and a polynomial-bit description and
separation test. It is the source's actual exponent rule, not a
point assigned after consulting a zero/nonzero oracle. Its torus
monomial basis, for $d\ge2$, is not a scalar-monomial $S_{2d}$ Young-orbit basis;
the latter is exactly what \fid{F28} excludes.


\subsection{Rectangular hooks in ordinary group cohomology}
This is an application of the known rectangular hook formula and
a known split-group motive decomposition, rather than a new
Kronecker rule.

\begin{finding}[\fid{F33}: a group-scheme realization with mixed degrees]\label{f:hookgeometry}
Let $n,d\ge1$, $N=nd$, $m=\min(n,d)$ and $0\le k<N$.
For $G=SL_{m,\F_q}$ and $\ell\nmid q$,
\[
 g((N-k,1^k),(d^n),(d^n))
   =\dim_{\Q_\ell}H^k(G_{\bar{\F}_q},\Q_\ell).
\]
More precisely there is an actual Weil decomposition
\begin{equation}\label{eq:hookweil}
 H^k(G_{\bar{\F}_q},\Q_\ell)=
 \bigoplus_{\substack{I\subseteq\{2,\ldots,m\}\\
                     \sum_{i\in I}(2i-1)=k}}
       \Q_\ell\left(-\sum_{i\in I}i\right).
\end{equation}
The degree can be mixed: for $m=8$,
\[
 H^{15}(SL_{8,\bar{\F}_q},\Q_\ell)
       =\Q_\ell(-8)\oplus\Q_\ell(-9).
\]
This realizes the numerical coefficients in ordinary cohomology;
it does not identify a nonstandard action or canonical basis.
\end{finding}

\begin{proof}
The known formula in IP, Proposition~6.4, is
\[
 \sum_{k=0}^{N-1}g((N-k,1^k),(d^n),(d^n))t^k
       =\prod_{i=2}^m(1+t^{2i-1}).
\]
For the geometric input, Merkurjev gives the canonical split
special-linear-group motive decomposition over a perfect field
\cite[\S5, Corollary~5.5, p.~8; notation on p.~7]{MerkurjevSL}.
In his notation a sequence
$1\le j_1<\cdots<j_a\le m-1$ contributes
$\mathbf1(d+a)[2d+a]$, with $d=\sum j_r$.
Reindexing by $I=\{j_1+1,\ldots,j_a+1\}$ and
$s_I=\sum_{i\in I}i$ gives
\[
 M(SL_m)_{\Q}
       \simeq\bigoplus_{I\subseteq\{2,\ldots,m\}}
                 \mathbf1(s_I)[2s_I-|I|].
\]
The displayed known motive theorem is invoked, not reproved.
Finite fields satisfy its perfect-field hypothesis.
Under ordinary cohomological $\ell$-adic realization, the summand
$\mathbf1(s)[a]$ contributes $\Q_\ell(-s)$ in degree $a$.
The decomposition is defined over $\F_q$, so Frobenius on that
summand is exactly $q^s$. This proves \eqref{eq:hookweil}.
Taking dimensions yields the product above, hence the required
coefficient identity.

For a summand in degree $k=2s_I-|I|$, its weight is
$2s_I=k+|I|$. Different subset cardinalities can therefore give
different weights in the same degree. At $m=8,k=15$, the only
subsets are $\{8\}$ and $\{2,3,4\}$. A two-element subset has
even degree, while four or more elements have minimum degree
$3+5+7+9=24$. Among three-element subsets the minimum 15 is
attained only by $\{2,3,4\}$. Their sums are eight and nine,
giving the stated Tate summands, weights 16 and 18, and traces
$q^{8r}+q^{9r}$ over $\F_{q^r}$.
\end{proof}

The corresponding hook coefficient is two for every rectangle
whose minimum side is eight. A single Tate twist cannot make the
two displayed weights equal. This is consistent with \fid{F20},
whose pure Satake multiplicity motive is a different realization
of the same integer. The affine group $G$ supplies ordinary
cohomology; no assertion of projectivity or purity of each
cohomological degree is needed.


\subsection{The open-orbit method has a precise higher-height limit}
\begin{finding}[\fid{F29}: the same flag-orbit mechanism stops at height two]\label{f:sphericity}
Let $n\ge3$, $E=\C^n\otimes\C^n$, $N=n^2$, and
$G=SL_n\times SL_n$. For every $2\le r\le N$, $G$ has no
open orbit on $\operatorname{Flag}(1,\ldots,r;E)$.
Equivalently, $G\times B_r$ has no open orbit on $E\otimes\C^r$,
where $B_r$ is a Borel subgroup. For $r=2$ the generic orbit has
codimension $n-2$. The same open-orbit proof underlying Manivel's
height-two formula therefore does not extend to these heights,
even after restricting the output length to two.
\end{finding}
\begin{proof}
On the dense locus of injections $(\C^r)^*\to E$, quotienting
by the appropriate dual Borel records the flag of successive images.
For each flag the compatible ordered injections form one free Borel
orbit. This is the usual frame bundle of the filtered tautological
bundle, locally a principal Borel bundle. It proves the equivalence
of the two open-orbit questions. The partial flag dimension is
\[
 rN-r(r+1)/2,\qquad \dim G=2(N-1).
\]
At $r=3$ the difference is $N-4>0$; successive increments of the
flag dimension are $N-r-1\ge0$ for $3\le r<N$. Thus no open
orbit exists for $r\ge3$. At full length the deficit is
$(N-1)(N-4)/2$.

For $r=2$, view $E$ as matrices, with action $X\mapsto gXh^t$.
The invertible first lines form a dense open orbit in $\mathbf P(E)$:
over $\C$ a scalar can normalize any invertible matrix's determinant.
A generic flag can therefore be written
$\C I\subset\C I+\C B$, with $B$ regular semisimple.
Its stabilizer satisfies $gh^t=cI$, $c^n=1$, and
$gBg^{-1}=aI+bB$, $b\ne0$. The affine transformation
$x\mapsto a+bx$ must permute the distinct eigenvalues of $B$.
Only finitely many such affine transformations exist, because the
images of two distinct eigenvalues determine $a,b$. The identity
component therefore has $c=1,a=0,b=1$ and is the centralizer of
$B$ in $SL_n$, a torus of dimension $n-1$. The generic orbit has
dimension $2N-n-1$ in a flag variety of dimension $2N-3$,
giving codimension $n-2>0$. Any open orbit would meet this dense
regular-semisimple locus, which is impossible.
\end{proof}
For $r=1$ the invertible matrix lines do form an open orbit, so the
boundary is real. The precise source premise being tested is Manivel
\cite[Lemma~1 and Corollary~1, \S3, p.~3]{Manivel}. Failure
of this open-orbit premise excludes its multiplicity-free algebra
mechanism; it does not exclude positive bases, useful generators,
or counting rules for multiplicity-bearing invariant algebras.

\subsection{An actual primitive motive for two-row rectangular coefficients}
The coefficient identity used in this subsection is known:
if $N=nd$, $0\le r\le\lfloor N/2\rfloor$, and $p_r(n,d)$ counts
partitions of $r$ in an $n\times d$ rectangle, then
\begin{equation}\label{eq:rectprimitiveknown}
 g_r:=g((N-r,r),(d^n),(d^n))=p_r(n,d)-p_{r-1}(n,d),
 \qquad p_{-1}=0.
\end{equation}
For $r\ge1$ the inspected primary locator is Pak--Panova,
\emph{Bounds on certain classes of Kronecker and $q$-binomial
coefficients}, author manuscript dated 13 October 2016,
Lemma~1.3, p.~2 \cite{PPBounds}; the source preserves the earlier
attributions. For $r=0$ the identity is the multiplicity one of the
trivial representation in a Specht tensor square. We apply this
known identity and classical Lefschetz theory; no historical
novelty is claimed for the underlying geometric interpretation.

\begin{finding}[\fid{F31}: a primitive Grassmannian coefficient object]\label{f:grassprimitive}
Let $n,d\ge1$, $N=nd$, $0\le r\le\lfloor N/2\rfloor$, and
$X=\Gr(n,n+d)$ over $\F_q$. Write $L=c_1(\mathcal O_X(1))$
for the Pl\"ucker class in $H^2(X_{\bar{\F}_q},\Q_\ell(1))$,
where $\ell\nmid q$. The actual Weil-module cokernel
\begin{equation}\label{eq:rectprimitivecoker}
 C_r=\operatorname{coker}\bigl(
 L:H^{2r-2}(X_{\bar{\F}_q},\Q_\ell)(-1)
       \longrightarrow H^{2r}(X_{\bar{\F}_q},\Q_\ell)\bigr)
 \simeq\Q_\ell(-r)^{\oplus g_r}
\end{equation}
is pure of weight $2r$, with scalar geometric Frobenius $q^r$;
the source is zero for $r=0$.

For $V_j=H^{2j}(X_{\bar{\F}_q},\Q_\ell)(j)$ and $r\ge1$, put
\[
 A_r=L^{N-2r+1}:V_r\longrightarrow V_{N-r+1},\qquad
 B_r=L^{N-2r+2}:V_{r-1}\longrightarrow V_{N-r+1}.
\]
Then $B_r$ is invertible and the Frobenius-equivariant projector
\begin{equation}\label{eq:rectprimitiveprojector}
 E_r=\id_{V_r}-L B_r^{-1}A_r
\end{equation}
has image $\ker A_r$ and kernel $L V_{r-1}$. For $r=0$ set
$E_0=\id$. Untwisting identifies its image with $C_r$.
It is the realization of an actual rational Chow-motive summand
$\mathsf P_r\simeq(\mathbb L^r)^{\oplus g_r}$, where
$\mathbb L$ denotes the Lefschetz motive. Its matrix in the
Schubert basis need not be nonnegative or integral.
\end{finding}

\begin{proof}
We first specify the arithmetic Tate decomposition. Let $\mathbb G_m$
act on the ambient space with weights $1,2,\ldots,n+d$.
The fixed $n$-planes are the rational coordinate planes indexed by
$I=\{i_1<\cdots<i_n\}$. The tangent weights are $j-i$ for
$i\in I$, $j\notin I$, so the positive-weight dimension is
\[
 a_I=\sum_{s=1}^n(d+s-i_s).
\]
The sequence $\lambda_s=d+s-i_s$ is a partition in the
$n\times d$ rectangle, and the inverse assignment is
$i_s=d+s-\lambda_s$. Thus $a_I=j$ occurs exactly $p_j(n,d)$
times. The motivic Bia\l ynicki--Birula decomposition in
Brosnan \cite[Theorems~3.3 and~3.5, author PDF pp.~5--6]{Brosnan}
applies to this smooth projective variety and split action.
Its fixed pieces are rational points, hence it gives, over $\F_q$,
\[
 h(X)\simeq\bigoplus_{j=0}^{N}(\mathbb L^j)^{\oplus p_j(n,d)},
 \qquad
 H^{2j}(X_{\bar{\F}_q},\Q_\ell)
    \simeq\Q_\ell(-j)^{\oplus p_j(n,d)},\quad H^{2j+1}=0.
\]
In particular Frobenius on $V_j$ is the identity. This actual
Tate decomposition supplies arithmetic splitting; purity alone
would not exclude a unipotent arithmetic extension.

The Pl\"ucker bundle is ample and defined over $\F_q$, and
$X$ is smooth projective of dimension $N$. BBD Theorem~5.4.10,
p.~144, applies to $X\to\Spec\F_q$ and the pure perverse
sheaf $\Q_\ell[N]$. Hard Lefschetz gives the stated isomorphism
$B_r$. Since $A_rL=B_r$, the map $L:V_{r-1}\to V_r$ is injective.
Set $T_r=L B_r^{-1}A_r$. Then
\[
 T_r^2=L B_r^{-1}(A_rL)B_r^{-1}A_r=T_r,\qquad
 A_r E_r=0,\quad E_r L=0,\quad
 E_r|_{\ker A_r}=\id.
\]
Also $E_rv=0$ implies $v=T_rv\in L V_{r-1}$.
Consequently
\[
 V_r=L V_{r-1}\oplus\ker A_r.
\]
The quotient map restricts to an isomorphism from the primitive
kernel. Its dimension is $p_r-p_{r-1}=g_r$ by the cited identity.
All maps are Frobenius equivariant, including $B_r^{-1}$, so
untwisting gives \eqref{eq:rectprimitivecoker}. More explicitly
the primitive subspace in the original degree is
\[
 \ker\bigl(L^{N-2r+1}:H^{2r}(X_{\bar{\F}_q},\Q_\ell)
 \longrightarrow
 H^{2N-2r+2}(X_{\bar{\F}_q},\Q_\ell)(N-2r+1)\bigr).
\]
The quotient-to-kernel splitting is canonical once $L$ is fixed.
The twist $(-1)$ in \eqref{eq:rectprimitivecoker} is necessary:
the rational divisor class is invariant in $H^2(1)$.

For the motivic assertion use rational coefficients in the displayed
Tate decomposition. Morphisms between sums of a fixed Tate motive
are rational matrices, and degree-zero morphisms between distinct
Tate powers vanish. After the displayed degree pieces are normalized
by their Tate powers, multiplication by $L$ and its powers are such
matrices, with rational entries given by algebraic correspondences.
Hard Lefschetz makes $B_r$'s determinant nonzero, so its inverse
is a rational matrix and an actual morphism in this Tate category.
Thus $E_r$ is an algebraic idempotent; the idempotent completion
gives its motive image with the stated rank. Brosnan's \S2,
author PDF pp.~3--4, specifies the correspondence and projector
category used here. This proves algebraicity in the cellular Tate
case; it does not assume the general standard conjecture on
Lefschetz inverses.

Finally take $X=\Gr(2,4)$ and $r=2$. In its integral Schubert
cycle lattice the add-one-box Pieri rule gives
\[
 L\sigma_1=\sigma_2+\sigma_{11},\qquad
 L\sigma_2=L\sigma_{11}=\sigma_{21}.
\]
These are the ordinary Schubert specializations of the Pieri rule
\cite[Theorem~7.10, p.~16]{Pieri}. Here $A_2=(1,1)$, $B_2=(2)$,
and
\[
 E_2=\frac12
 \begin{pmatrix}1&-1\\-1&1\end{pmatrix}.
\]
Its image is the primitive Tate line. The lattices
$\Z(\sigma_2+\sigma_{11})$ and
$\Z(\sigma_2-\sigma_{11})$ have sum of index two in
$\Z\sigma_2\oplus\Z\sigma_{11}$. Thus the prescribed
primitive-versus-Lefschetz decomposition does not split this
integral lattice, and its rational projector has negative entries.
This does not prohibit every integral projector onto the primitive
line: $\bigl(\begin{smallmatrix}1&0\\-1&0\end{smallmatrix}\bigr)$
is an integral idempotent with the same image and a different kernel.
\end{proof}

The two rectangular factors in \fid{F31} have arbitrary height and
width. Its restriction is the two-row output $(nd-r,r)$: it covers
single-part padded tails, not IP's general multipart tails. Within
IP's padded-tail range $|\bar\lambda|\le md$, it applies when
$r\le\min(md,\lfloor nd/2\rfloor)$; its own construction requires
no IP size threshold. The Satake coefficient in \fid{F20} for the
same two rectangles, using $\GL_n\times\GL_n$, has weight
$2nd(n-1)$, whereas \fid{F31} has weight $2r$. A Tate twist can
align weights, but their equal ranks do not choose an isomorphism,
an operation-compatible basis, or the nonstandard path comparison.
This finding uses classical hard Lefschetz and a known coefficient
identity; it does not derive that identity or purity anew from the
sign-only path argument.


% Required bibliography keys: PPS (arXiv:2511.02312v3, Lemma 1.3, p.2)
% and KM (Kahle--Michalek, EJC 25(1) (2018), P1.41).
\subsection{A positive plethysm ray whose required dimension exceeds the compact bound}
The compact input convention in the supplied VI is explicit: zero parts
are suppressed, and height is the number of positive parts
\cite[\S1.1, p.~9; opening of \S3.4, p.~53]{VI}.
In particular, the source deliberately omits the trailing zeros of an
outer weight even when the ambient group has very large rank.
Its stretching operation fixes the inner weight
\cite[equation~(1.5), p.~20]{VI}.
These conventions matter for the following elementary deduction from a
classical two-row identity.

\begin{finding}[\fid{F34}: a necessary dimension and an encoding obstruction]\label{f:binarydimension}
Fix $k\ge3$, take $H=\GL_2$, inner weight $\mu=(k)$,
outer weight $\lambda=(1)$, and output weight $\pi=(k-1,1)$.
Their stretching multiplicity is
\[
 c_k(N)=a^{((k-1)N,N)}_{(N),(k)}
       =[\Sym^N(\Sym^k\C^2):S_{((k-1)N,N)}\C^2].
\]
For every $N\ge0$,
\[
 c_k(N)=
 \#\{(x_2,\ldots,x_k)\in\Z_{\ge0}^{k-1}:
                          2x_2+\cdots+kx_k=N\}.
\]
It is the Ehrhart function of the rational simplex
\[
 P_k=\{x\in\R_{\ge0}^{k-1}:2x_2+\cdots+kx_k=1\},
\]
whose affine dimension is $k-2$. Every fixed bounded polytope with this
Ehrhart function has dimension at least $k-2$.

Consequently the uniform dimension bound in the supplied
VI, Hypothesis~1.6.4(2), p.~23,
\[
 m=\operatorname{poly}
   (\langle\mu\rangle,\operatorname{ht}(\lambda),
                         \operatorname{ht}(\pi)),
\]
is false under its stated compact binary input convention.
Moreover, the plethysm counting function with all nonzero parts of
all three partitions encoded in binary is not a $\#P$ function:
already this family can require exponentially many output bits
relative to the input length.
These assertions do not apply to a unary weight convention or to
a bound polynomial in the numerical $k$.
\end{finding}
\begin{proof}
At $N=0$ both counts are one, so assume $N\ge1$.
The classical two-row specialization, restated in
\cite[Lemma~1.3, p.~2]{PPS}, gives
\[
 c_k(N)=p_N(N,k)-p_{N-1}(N,k),
\]
where $p_r(N,k)$ counts partitions of $r$ in the $k$ by $N$ rectangle.
For $r=N,N-1$ the width bound is automatic. After conjugation, the
remaining condition is that all parts have size at most $k$. Thus
\[
 c_k(N)
 =[z^N](1-z)\prod_{i=1}^k(1-z^i)^{-1}
 =[z^N]\prod_{i=2}^k(1-z^i)^{-1},
\]
which proves the positive count and the asserted simplex formula.

The weight calculation also checks which plethysm stretch is involved.
A monomial of $\Sym^N(\Sym^k\C^2)$ has nonnegative multiplicities
$a_0,\ldots,a_k$ satisfying $\sum a_i=N$ and second-coordinate weight
$\sum ia_i=r$. At $r=N$ or $N-1$, the determined value of $a_0$ is
automatically nonnegative. The GL$_2$ highest-weight multiplicity is
the difference of these adjacent weight dimensions, since
$N\le kN/2$. The inner $k$ is fixed, exactly as in VI's definition.

Put $L_k=\operatorname{lcm}(2,\ldots,k)$. For $N=L_ks$, the map
\[
 (y_2,\ldots,y_k)\longmapsto
       ((L_k/2)y_2,\ldots,(L_k/k)y_k)
\]
injects the nonnegative solutions of $\sum y_i=s$ into the displayed
count. Hence
\[
 c_k(L_ks)\ge\binom{s+k-2}{k-2}.
\]
Conversely, choosing $x_2,\ldots,x_{k-1}$ determines $x_k$ if it
exists, so $c_k(N)\le(N+1)^{k-2}$.

For any fixed bounded polytope $Q$ of affine dimension $d$,
$\#(NQ\cap\Z^m)=O_Q(N^d)$. To see this without a further asymptotic
theorem, choose $d$ coordinate functions injective on the direction
space of its affine hull. They inject the lattice points of $NQ$
into an integer box with side lengths $O_Q(N)$.
Comparison along $N=L_ks$ therefore forces $d\ge k-2$.
The displayed simplex has this dimension and attains the bound.

Here $\langle\mu\rangle=O(\log k)$,
$\operatorname{ht}(\lambda)=1$ and
$\operatorname{ht}(\pi)=2$, while the original group $H$ has fixed
rank two. Any ambient dimension for the claimed PH1 polytope is at
least its affine dimension $k-2$, which eventually exceeds every
fixed polynomial in these compact parameters. The omitted trailing
zeros in the $\GL_{k+1}$ outer weight cannot be charged to the input:
VI explicitly excludes them.

For the output-length assertion choose an integer $m\ge2$ and set
$k=2m$, $N=m(m+1)$. In the positive count choose independently
\[
 x_{2i}\in\{0,1\}\quad(2\le i\le m),\qquad
 x_j=0\quad(j\text{ odd}),
\]
and set
\[
 x_2=\frac{N-\sum_{i=2}^m2i\,x_{2i}}2.
\]
This is a nonnegative integer because the maximum subtracted sum
is $m(m+1)-2=N-2$. Distinct choices give distinct solutions, so
\[
 a^{((2m-1)N,N)}_{(N),(2m)}\ge2^{m-1}.
\]
The three partitions together have four nonzero parts, each with
$O(\log m)$ bits. A $\#P$ function on inputs of length $L$ has at
most $2^{p(L)}$ accepting witnesses for some fixed polynomial $p$.
The lower bound $2^{m-1}$ on these $O(\log m)$-bit inputs contradicts
that requirement for every fixed $p$.
\end{proof}

\paragraph{Known precedent and the scope of this deduction.}
The universal ordinary Ehrhart-existence claim in the supplied VI
was already refuted by Kahle--Micha{\l}ek
\cite[Theorem~5 and Remark~6, p.~4]{KM}.
Their fixed-small-parameter example has outer weight $(4)$, inner
weight $(3)$ and output weight $(7,5)$; its stretching function
violates Ehrhart reciprocity. Their p.~2 explicitly identifies the
question with VI, Hypothesis~1.6.4. Their pp.~3--4 distinguish the
later inhomogeneous formulation, which their result does not
dispose of in every dimension.
We cite that prior result for the historical failure of the old
existence assertion. Finding~\fid{F34} isolates the additional
compact-dimension issue on a family that \emph{does} have an exact
positive simplex model; no priority claim is made.

In the effective representation ring the numbers $c_k(N)$ are actual
Schur multiplicities, and their ordinary pure Satake realizations
are compatible with \fid{F20}/\fid{F31}. Purity does not change the
growth degree of the multiplicity. An encoding that includes the
numerical rank $\dim\Sym^k\C^2=k+1$ can accommodate this simplex,
but is different from VI's compact input. The associated rectangular
Kronecker input explicitly lists $k$ rows, so its input length is
already at least $k$; the binary plethysm output-length obstruction
does not apply to that encoding.


% Required bibliography keys: PPS (arXiv:2511.02312v3) and KM.
\subsubsection{The actual polynomial highest-weight algebra on the same ray}
The positive rule in \fid{F34} can be realized inside its ordinary
highest-weight algebra. The following direct slice calculation is a
constructive supplement to that finding, without a priority claim for
the classical method.

\paragraph{Claim.}
For $k\ge2$ put $V=\Sym^k\C^2$ and choose its divided basis
\[
 v_i=x^{k-i}y^i/i!,\qquad 0\le i\le k.
\]
In the free symmetric algebra $A=\Sym(V)=\C[v_0,\ldots,v_k]$,
let $D(v_0)=0$, $D(v_i)=v_{i-1}$, and assign
$\deg v_i=1$, $\operatorname{ind}v_i=i$.
The ray algebra
\[
 B=\bigoplus_{N\ge0}
       \bigl(A_{\deg=N,\,\operatorname{ind}=N}\cap\ker D\bigr)
\]
is a polynomial ring on $k-1$ explicit generators $f_2,\ldots,f_k$
of degrees $2,\ldots,k$. Its actual degree-$N$ monomial basis is
indexed by $2a_2+\cdots+ka_k=N$, and
\[
 \dim B_N=[\Sym^N(\Sym^k\C^2):S_{((k-1)N,N)}\C^2]=c_k(N).
\]
The symbols $v_i$ are free algebra generators here; their displayed
binary forms specify a basis and action, rather than an identification
of $A$ with a subalgebra of $\C[x,y]$.

\begin{proof}
The raising operator $x\partial_y$ acts by $Dv_i=v_{i-1}$ in
the divided basis. A degree-$N$, index-$r$ monomial has torus
weight $(kN-r,r)$. In characteristic zero, the vectors of this
weight killed by the raising operator are its highest-weight
vectors. Since $k\ge2$, the weight $((k-1)N,N)$ is dominant,
and ordinary highest-weight decomposition proves the multiplicity
identity. This also agrees with the classical two-row specialization
in \cite[Lemma~1.3, p.~2]{PPS} used in \fid{F34}.

In $A[v_0^{-1}]$ put $t=v_1/v_0$, so $Dt=1$, and define
\[
 f_i=v_0^{i-1}\sum_{j=0}^i\frac{(-t)^j}{j!}v_{i-j},
       \qquad 2\le i\le k.
\]
These elements belong to $A$. Indeed the terms with $j\le i-1$
are
\[
 \frac{(-1)^j}{j!}v_0^{i-1-j}v_1^jv_{i-j},
\]
and the last term is $(-1)^iv_1^i/i!$.
Each has degree and index $i$. Differentiating the sum, the terms
from $D((-t)^j/j!)$ cancel those from $D(v_{i-j})$ after shifting
the index; the remaining endpoint has $D(v_0)=0$.
Thus $Df_i=0$. For example,
\[
 f_2=v_0v_2-\tfrac12v_1^2,\qquad
 f_3=v_0^2v_3-v_0v_1v_2+\tfrac13v_1^3.
\]

The formula for $f_i$ is triangular, with leading term
$v_0^{i-1}v_i$ and remaining terms involving only earlier generators.
The leading coefficient is a unit in the localization. Solving
successively for $v_1=tv_0,v_2,\ldots,v_k$ therefore gives an
invertible change of free coordinates,
\[
 A[v_0^{-1}]=\C[v_0,v_0^{-1},t,f_2,\ldots,f_k].
\]
In these coordinates $D=\partial/\partial t$. Characteristic zero
then gives
\[
 A[v_0^{-1}]^D=\C[v_0,v_0^{-1},f_2,\ldots,f_k].
\]
This assertion concerns the explicitly computed localization; no
description of the entire unlocalized invariant ring is assumed.

Every monomial $v_0^a\prod_{i=2}^k f_i^{a_i}$ in this Laurent
polynomial ring has degree minus index equal to $a$.
Algebraic independence and homogeneity therefore imply that its
component with degree equal to index has $a=0$ and is exactly
$\C[f_2,\ldots,f_k]$. All the $f_i$ already belong to $A$, so
intersecting with $A$ loses nothing. Hence
\begin{equation}\label{eq:actualrayalgebra}
 B=\C[f_2,\ldots,f_k],\qquad \deg f_i=i.
\end{equation}
Its displayed monomials form the asserted actual highest-weight
basis. The construction is defined over $\Q$ and then base changed
to $\C$; all divided bases, factorials and the derivative-kernel
argument use characteristic zero.
\end{proof}

This gives multiplication constants equal to one in this particular
ordinary algebra. Its $k-1$ generators and their formulas have size
polynomial in the numerical $k$, not in $\log k$. The construction
therefore retains both the exact positive simplex and the compact
dimension obstruction in \fid{F34}. It does not identify this
monomial basis with VIII's nonstandard multiplication or
comultiplication basis, and asserts no theta or Weil realization.


\subsection{Short witnesses for the actual marked rectangular tree rule}
Pak--Panova--Swanson \cite[Theorem~4.5, p.~10]{PPS} give a positive
marked-KOH-tree rule for
\[
 g((nk-r,r),(n^k),(n^k)),\qquad
 n,k\ge1,\quad 0\le r\le nk/2.
\]
Their rule is a known result. The following encoding deduction makes
its witness size explicit when the rectangular partitions are supplied
as lists of $k$ binary-encoded parts.

\begin{finding}[\fid{F37}: a bijective compression of marked KOH witnesses]\label{f:kohcompression}
Every source marked KOH tree has a unique compressed encoding of length
\[
 O\!\left(k^2\log(k+1)+k\log(nk+1)\right)
\]
with a polynomial-time verifier that does not expand long unary chains.
The compression preserves the source objects bijectively.
Consequently the equal-rectangle, two-row-output Kronecker counting
problem is in $\#P$ under the explicit-row, binary-part input convention.
No bound polynomial in $\log k+\log n+\log(r+1)$ is asserted for a
rectangle specified only by two binary integers, and no Ehrhart PH1
conclusion follows.
\end{finding}
\begin{proof}
We first state the source rules needed for the encoding
\cite[Definition~3.2, p.~6; Definition~4.4, p.~10]{PPS}.
A node has a label $(\mu,a,b)$ with $a\ge0$, $b\ge1$, $\mu\vdash b$.
Leaves have $b=1$, $\mu=(1)$. A nonleaf has one child for every
distinct part $j$ of $\mu$, in increasing $j$ order, with type
\[
 a_j=(a+2)j-2Q_j(\mu),\qquad b_j=m_j(\mu),\qquad
 Q_j(\mu)=\sum_s\min(j,\mu_s).
\]
All child $a_j$ must be nonnegative. The root type is $(n,k)$.
For depth-first leaf labels $a_1,\ldots,a_t$, put
$\sigma=nk-\sum_i a_i$. The leaf marks satisfy
\[
 0=u_1\le\cdots\le u_t=r-\sigma/2,\qquad
 u_{i+1}-u_i\le
 \min\left(\sum_{j\le i}a_j-2u_i,\ a_{i+1}\right).
\]

Call a nonleaf a repeat node when $\mu=(1^b)$, $b\ge2$.
It has exactly one child, with type
\[
 (a,b)\longmapsto(a-2(b-1),b).
\]
Every other nonleaf obeys
\[
 \sum_j b_j=\ell(\mu)\le b-1.
\]
Contract each maximal consecutive run of repeat nodes, storing its
length $h$ in binary before the next surviving node. The entire run
is the update
\[
 (a,b)\longmapsto(a-2h(b-1),b).
\]
Its endpoint being nonnegative guarantees every intermediate
$a$-value is nonnegative. Retain the surviving partition labels,
ordered children and all leaf marks. A leaf has incoming repeat
count zero, because repeats preserve $b\ge2$.

This compression is bijective. Starting from the root type, a
compressed record first performs the displayed run update, then
requires its surviving partition to be different from $(1^b)$
unless it is a leaf. Its ordered child types are determined by the
source formula. Expanding a count $h$ inserts exactly $h$ nodes
labeled $(1^b)$ with the uniquely determined successive $a$-values.
The surviving node is nonrepeat, so the inserted run is maximal.
Compression and expansion are inverse operations. Unary contraction
does not change depth-first leaf order, any leaf label, or any mark;
all marking inequalities are therefore preserved.

Let the compressed tree have $H$ internal vertices and $t$ leaves.
The $b$-labels on each contracted edge are unchanged. Summing the
positive losses at internal vertices telescopes:
\[
 H\le
 \sum_{v\ {\rm internal}}
       \left(b_v-\sum_{w\ {\rm child\ of}\ v}b_w\right)
 =k-t.
\]
Thus there are at most $k$ surviving vertices. Every surviving
partition has at most $k$ parts, each at most $k$, accounting for
$O(k^2\log(k+1))$ bits.

To bound all numerical labels, set
$\beta(\mu)=\sum_s(s-1)\mu_s$. Decreasing row lengths give
\[
 \sum_jm_j(\mu)Q_j(\mu)
 =\sum_{r,s}\min(\mu_r,\mu_s)
 =b+2\beta(\mu).
\]
Hence
\[
 \sum_j a_jb_j=ab-4\beta(\mu)\le ab.
\]
In a valid tree all these terms are nonnegative, so every vertex
satisfies $0\le ab\le nk$. Each $a$-label is at most $nk$, and
every contracted repeat count is at most $a/(2(b-1))\le nk$.
The same identity gives $\sigma\ge0$. Every mark lies between
zero and $r\le nk/2$. At most $k$ repeat counts and $k$ marks
therefore take $O(k\log(nk+1))$ bits. Tree punctuation and
partition-length fields fit in the stated bound as well.

Use one canonical preorder serialization, with canonical integer
fields, decreasing partition parts, and the fixed child and leaf
orders. The verifier checks the number of surviving nodes, the
partition conditions, the root and recursive child types, the
binary run update, zero incoming repeat counts at leaves, and
then all leaf marking inequalities.
It can reject any derived type violating $0\le ab\le nk$;
the preceding identity proves that no valid source object is lost.
All work uses polynomially many operations on integers of the
stated bitlengths, and no run is expanded.

If the explicit-row input has length $L$, then $L\ge k$ and it
includes the binary lengths of $n$ and $r$. The witness bound and
the verification time are polynomial in $L$. Canonical serialization,
with required zero padding to a fixed polynomial length if desired,
gives exactly one accepting string for every compressed object.
The bijection and the cited PPS counting theorem identify the
number of accepting strings with the Kronecker coefficient.
\end{proof}

\paragraph{A literal tree can be exponentially long.}
Take $n=2s$, $k=2$, $r=2s$. At $b=2$ the only source partitions
are $(1,1)$ and $(2)$. After $j$ repeats of $(1,1)$, choosing $(2)$
ends in one leaf labeled $4s-4j$, so $\sigma=4j$.
Its sole mark is zero, forcing $0=r-\sigma/2=2s-2j$.
Thus the unique marked tree has $s$ repeat vertices, one final
$(2)$-vertex and one leaf. It has $s+2$ vertices on an input of
$O(\log s)$ bits. Its compressed encoding has two vertices and
one binary repeat count $s$. This is why local checks alone would
not have justified the binary witness bound.

\paragraph{Encoding and polyhedral limits.}
The related plethysm inputs $(n)$ and $(k)$ can both have one part,
so a bound polynomial in the numerical $k$ is not automatically
polynomial in their compact input length. Finding~\fid{F34} gives
an explicit output-length obstruction to that stronger assertion.
The present rectangular input instead lists $k$ rows.

Nor does a short witness rule imply one ordinary dilation polytope.
Here the tree shape and its partition labels are discrete choices,
and the local relations include fixed offsets. Indeed the known
Kahle--Micha{\l}ek counterexample
\cite[Theorem~5 and Remark~6, p.~4; rectangular identification, p.~6]{KM}
already supplies an equal-rectangle, two-row-output ray that is
not an ordinary Ehrhart function. The positive marked-tree rule
and the verified compression still apply to it.
This finding settles a restricted effective-label gate; it supplies
no designated nonstandard canonical basis, theta functor,
structure-constant positivity, or orbit-closure obstruction.


\subsection{Occurrence versus multiplicity}
Let $\Omega_n=\overline{\GL_{n^2}\cdot\det_n}$, $Z_{n,m}$ the padded-permanent orbit closure, and $K_n(\lambda)$ the degree-$d$ multiplicity in $\C[\Omega_n]$. B\"urgisser--Ikenmeyer--Panova \cite[Theorem 1.4]{BIP}, in the attached arXiv 1604.06431v3, proves for $n\ge m^{25}$:
\[
 q^m_\lambda(d[n])>0\quad\Longrightarrow\quad K_n(\lambda)>0.
\]
Theorem~1.5 is the numbering in an older version, not this attachment. The paper's \S1(b) distinguishes
\[
 K_n(\lambda)\le sk_n(\lambda)\le g_n(\lambda),
\]
where $sk_n$ is the symmetric rectangular coefficient controlling the determinant \emph{orbit}, not automatically its closure.

If $Z_{n,m}\subseteq\Omega_n$, restriction is a graded equivariant surjection
$\C[\Omega_n]\twoheadrightarrow\C[Z_{n,m}]$. Semisimplicity of the characteristic-zero representation in each finite-dimensional degree implies
\[
 q^m_\lambda(d[n])\le K_n(\lambda).
\]
Thus a strict reverse multiplicity inequality remains a valid obstruction, as BIP \S1(c) explains. A theorem about nonzero support does not settle it. The ordinary Kronecker upper bound supplies a sufficient stronger inequality, not an equality with closure multiplicities.

\paragraph{Consequence of the present work.}
\fid{F03}--\fid{F04} establish purity only for data satisfying their trace and sign hypotheses. No such data have been identified here for the normalized coefficients in VIII, the full rectangular region, or the closure multiplicity comparison. Therefore no new general positivity or multiplicity obstruction is derived from them. This is the exact limit of the proved transfer, not a restatement of numerical nonnegativity as a new theorem. The permanent--determinant problem concerns algebraic complexity classes; an additional argument would be required before making a Boolean $P\ne NP$ claim.

\part{The findings in Grothendieck's language}
\section{An exact dictionary, not an attribution of new theorems}\label{sec:groth}
The new deductions above are restated here in the requested frameworks. A framework locator says where the language and operations live; it does not claim that a theorem about GCT occurs in Grothendieck's manuscripts. Where no archival proof or identification was verified, we say so explicitly.

\subsection{Schemes and group schemes over the integers}
The classical objects are group schemes $\GL_{r,\Z}$ and their polynomial representations on finite locally free modules. Tensor products and integral Schur constructions give maps such as
\[
 \GL(V)\times\GL(W)\longrightarrow\GL(V\otimes W),\qquad
 \GL(V)\longrightarrow\GL(S_\mu V).
\]
Representations of an affine group scheme are comodules of its \emph{commutative} coordinate Hopf algebra (SGA~3, I, \S4.2 and 4.7.2). Representations of diagonalizable groups correspond to graded modules (I, 4.7.3). These descriptions carry \eqref{eq:plethysm}, \fid{F11} and \fid{F13} into group-scheme language. Cote~90, scan~17, gives an authentic manuscript anchor for a flat affine group over a base, a group action, and equivariant quasi-coherent modules.

Over characteristic zero the representation categories of the reductive groups here are semisimple, so their actual classes form the positive monoid in the irreducible basis of $K_0$. This is the precise content of scalar nonnegativity. Over $\Z$ or in characteristic $p$, do not replace a filtration by a direct sum of simple modules. A divided Young symmetrizer is only an idempotent when its denominators are invertible; integral Schur functors must be constructed integrally. For the integral Specht-model representations used in the Kronecker and plethysm instances of \fid{F12}, tensor the free integral lattices with $\Z[1/|H|]$. The averaging idempotent then has a finite locally free direct-summand image, an effective $K_0$ class whose complex rank is the stated multiplicity. More generally this construction works over any coefficient ring carrying the representation and in which $|H|$ is invertible; an arbitrary complex representation need not descend to $\Z[1/|H|]$.

The nonstandard bialgebra preserves some comodule language but is not automatically a commutative affine group scheme. \fid{F11}'s matrix calculation proves why a reduction modulo $p$ supplies no automatic absolute Frobenius. This is a limitation of that construction, not a statement against separately defined quantum Frobenius maps.

\subsection{Representation rings, $\lambda$-rings and the $\gamma$-filtration}
For locally free sheaves, $\lambda^i[E]=[\wedge^iE]$ and Schur operations are the associated symmetric-function operations. On virtual classes the identity $\lambda_t(x-y)=\lambda_t(x)/\lambda_t(y)$ belongs to the group-completed ring. SGA~6, Expos\'e~V, \S\S2--3, contains the $\lambda$ and $\gamma$ formalism; cote~133, scan~18, explicitly displays symmetric/exterior generating functions. Cote~97, scan~4, concerns Chern classes of representations. These are the appropriate anchors for the formal calculation in \fid{F11}.

The representation ring example is an exact dictionary, not an analogy: $z$ is the standard one-dimensional comodule of $\mathbb G_m$; $z-1$ is a rank-zero virtual class; its exterior-power series gives $\lambda^2(z-1)=1-z$. The augmentation ideal and its $\gamma$-filtration retain characteristic-class information but introduce no order turning this virtual class into an actual representation. In particular the equality of two expressions in a $\lambda$-ring cannot replace the effective-simple-basis hypothesis of \fid{F07}.

For singular schemes, one must distinguish $K_0$ of vector bundles or perfect complexes from $G_0$ of coherent sheaves. The GCT graded pieces themselves are finite-dimensional representations, so their representation-$K_0$ classes are well-defined without identifying $K_0(X)$ and $G_0(X)$ on a singular orbit closure. Coherent Euler pushforward may involve alternating higher cohomology; an actual section module is more data than a virtual Euler class.

\subsection{Weil realizations and virtual versus effective traces}
In the trace formalism, a bounded complex of finite-dimensional Weil representations yields
\[
 Z(K,T)=\prod_i\det(1-TF\mid H^i(K))^{(-1)^{i+1}},\qquad
 \log Z(K,T)=\sum_{n\ge1}\Tr(F^n\mid K)\frac{T^n}{n}.
\]
The trace is an alternating trace. \fid{F02} bounds the surviving finite spectrum after equal roots have been grouped. \fid{F05}(b) says that $K_0$ and $Z(K,T)$ can erase identical spectra in adjacent degrees. An effective parity statement prevents this erasure; that is the exact role of (S2).

The geometric dictionary of \S\ref{sec:tracegeom} is the proof-carrying transfer of \fid{F01}--\fid{F04} and \fid{F17}. Character orthogonality is expressed by Artin--Schreier sheaves; the external tensor product realizes powers; compactly supported direct image realizes the sum; one-degree support fixes the alternating sign; Tate twist supplies the factor $q^n$; duality supplies the reciprocal eigenvalues. Framework locators are SGA~4, XVII, \S3 (compact supports), XVIII, 3.2.5 (smooth duality); SGA~5, III, 4.7 (Lefschetz--Verdier); and SGA~$4\frac12$, \emph{Rapport}, 3.2 (the trace formula). The exact ordinary weight estimates are Weil~II, 3.3.1 and 3.3.4--3.3.6. These general theorems are cited background; this run does not re-prove their machinery.

For a scheme over $\F_q$ and a closed point $x$ with residue field $\F_{q^{\deg x}}$, a Weil sheaf is pure of weight $w$ if every eigenvalue $\alpha$ at $x$ has
\[
 |\sigma(\alpha)|=q^{(\deg x)w/2}
\]
for every relevant complex embedding. For complexes the weight conditions are degree-shifted and involve duality; purity of a complex does not assert that every singular stalk has exact weight equal to its degree. The KL theorem uses the stronger stalkwise even/Tate statement \eqref{eq:KL}. This distinction translates \fid{F06} and prevents a false purity-to-coefficient inference.

\subsection{Orbit closures, Hilbert/Quot, and motives}
\fid{F15}'s kernel construction is the affine scheme-theoretic-image construction of EGA~I (1960), \S9.5. Its explicit $f=p$ example is already a group-scheme example over $\Spec\Z$. Projective closures with fixed Hilbert polynomial and coherent quotients belong to FGA, Bourbaki Expos\'e~221, \S3, Theorems~3.1--3.2, pp.~259--260. Fixing the Hilbert polynomial uses the chosen polarization, and a Quot functor specifies the coherent sheaf being quotiented. They do not turn coherent section multiplicities into intersection-cohomology multiplicities by representability alone.

For smooth projective schemes with algebraic correspondences, pure motives and their realizations provide another language for duality and projectors. The Tate realization in \fid{F16} is legitimate. By contrast a general singular orbit-closure intersection complex is not automatically a Chow motive furnished by SGA~4/5/6, and a quantum coordinate algebra acquires no such motive from its name. Cote~16, scans~35, 37 and~39, supplies a historical anchor for Lefschetz operators and primitive projectors; scans~47--49 discuss conditional correspondence/motive consequences, with conjectural letters interpreted in the letter's own notation. Their algebraic-projector and index assumptions cannot be replaced by the Euler sign $(-1)^i$.

\subsection{The constructive models and their exact comparison limits}
\paragraph{\fid{F18}: weight lattices and rational semigroups.}
Regard $g(\nu,(d,d),(d,d))$ as the coefficient of
$\det^d\boxtimes\det^d$ in the restriction of $S_\nu(V\otimes W)$,
where $V,W$ have rank two. This is an equality in the representation
$K_0$ of the group scheme $\GL(V)\times\GL(W)$ after passage to
characteristic zero. Characters of a diagonalizable group correspond
to gradings (SGA~3, I, 4.7.3). Integral difference-constraint arrays
would supply a particular saturated graded semigroup. The flooring
and dilation proof of \fid{F18} transfers unchanged to that weight
lattice. Its parity example requires rational-lattice or congruence
data instead. SGA~6's formal $\lambda$ identities do not supply those
data; cote~133, scans~18--21, is a vocabulary anchor rather than a
positive-rule theorem.

\paragraph{\fid{F19}, \fid{F21}, \fid{F25}: coherent algebra and its actual module.}
Here there is a specified scheme over a localization of
$\Spec\Z[v,v^{-1}]$: the finite flat center $S$ of \fid{F19}, with
the explicitly trivial bundle $E$ of ranks $1,1,2$.
The noncommutative algebra is $\Gamma(S,\End E)$, and the actual
right tensor module is $\Gamma(S,\mathcal Hom(E,\mathcal F))$ for
the displayed coherent branch sheaf $\mathcal F$. Thus the module
dictionary is an explicit action-preserving isomorphism, not just
an equality of characters. The matrices in \fid{F21} carry the proof
over to this sheaf language. Equivariant quasi-coherent modules are
the framework of cote~90, scan~17, and SGA~3, I, 4.7.2.

The branch modules $L_i$ are coherent but are not asserted to be
locally free or perfect over the singular center. Their natural
Grothendieck group is $G_0(C)$; one must not silently apply vector-bundle
$\lambda$ operations from $K_0(C)$ to them. The periodic resolution
in \fid{F25} computes their actual coherent derived Hom. Its
$R/(\Delta)$ records an order-two deformation invisible in the
constant coefficient realization. The regular algebra basis in that
finding is unrelated to any assertion that this singular coherent
category is semisimple. Its scalar $\lambda$ has compact-support motive
\[
 \mathbb L^{-2}(\mathbb L-1)^2(1+\mathbb L+\mathbb L^2)
\]
in the virtual Tate ring. The trace equality follows from
$\mathbb G_m^2\times\PP^2$ and the trace formula; its mixed signs are
cohomological Euler signs, not dimensions of one pure degree.

\paragraph{\fid{F22}: the \'etale topos and the structure sheaf retain different data.}
The finite maps $\Spec C_n\to\Spec C_m$ are universal
homeomorphisms. EGA~IV, 18.1.2, in the precise modern formulation
Stacks 04DZ/03SI, transports the entire ordinary coefficient-sheaf
topos over the trait. The proof in \fid{F22} explicitly identifies
the constant nearby-cycle triangle with
$\Q_\ell\to\Q_\ell^2\to\Q_\ell$; all three objects have scalar
weight-zero Frobenius. In contrast the conductor and crossed coherent
$\operatorname{Ext}^1$ retain $m$. This is an exact instance of the
difference between a coefficient topos and a ringed coherent space.
It does not forbid adding a separately specified nilpotent action to
a Weil sheaf. No manuscript statement is claimed to prove this
nonstandard deformation comparison.

\paragraph{\fid{F20}: effective Schur operations and Tate coefficient motives.}
On a split affine Grassmannian, the primitive Young projector is an
actual idempotent in a rational symmetric tensor category. Its image
is an effective object, and its class is the specified Schur operation
on the locally defined representation. The modern motivic Satake
theorems cited in \fid{F20} identify its coefficient motive as
$\mathbf1(-h)^{\oplus a}$, not merely an alternating virtual class.
The ordinary coefficient realization has scalar Frobenius $q^h$.
This supplies an actual motive in these classical cases; it does not
assert the existence of a Chow motive for every singular orbit closure.
The source of the motivic theorem is Richarz--Scholbach, not SGA~6
or cote~16. The latter's scans~35,37,39 supply the historical language
of algebraic projectors. Canonical bases and efficient rank labels
remain separate information even when the motive is Tate.

\paragraph{\fid{F23}: induction in the effective representation monoid.}
A signed orbit is the representation of a constant finite group
scheme induced from a one-dimensional stabilizer representation.
Over $\Z[1/d!]$ with the relevant integral lattices, finite-group
averaging splits invariant projectors; over characteristic zero
the representation category is semisimple. The direct-sum equality
of \fid{F23} is effective in representation $K_0$, and its Young
stabilizer specialization is the ordinary symmetric/exterior Pieri
operation. The proof by orbit modules and restriction of Schur
functors is unchanged in this language. The pure Tate $\GL_2$
example shows why the $S_d$ correspondence action is additional to
the weight and motive. Cotes~90/17 and~133/19--21 give the respective
module and finite-set-operation anchors; neither provides the missing
universal signed basis.

\paragraph{\fid{F24}: a signed class can still retain a largest eigenvalue.}
For a bounded Weil complex, group equal roots in its alternating trace.
Its visible multiplicities are integers, possibly of both signs.
The external tensor product maps this class to its power in the
group ring of eigenvalues. The augmentation-and-Vandermonde argument
in \fid{F24} proves that an even power retains a positive top-radius
coefficient. This transfers the proof exactly to the triangulated
$K_0$ trace formalism of SGA~4/5, without requiring one input degree.
The negative visible path-output spectrum and any reciprocal visible
duality are still hypotheses. A zero virtual class may erase actual
cohomology, and finite extensions can erase a signed top spectrum;
the two countermodels specify those limits.

\paragraph{\fid{F26}--\fid{F27}: localization, bounded complexes, and finite traces.}
Finite Weil cohomology gives rational resolvents and zeta factors
$\det(1-TF)$, including their alternating products. This is the
finite determinant mechanism underlying the trace formalism; the
non-Laurent rational grading function in \fid{F26} cannot be realized
by such a fixed finite trace in all extensions. Its recurrence proof
is independent of any particular choice of category. The generic
Fredholm determinant formulas visible in cote~1, scans~77--78,
do not identify it with a finite Weil object or a $\Xi$-function.

Bounded complexes in the Hecke category do realize the signed Euler
class $P$. Their four basis coordinates prove the one-parity
obstruction in \fid{F27}. The rank-one bimodule comes from the
finite flat correspondence
$\Spec A\times_{\Spec A^s}\Spec A$. Separately, the coefficient
Hecke algebra $H=R[c]/(c^2-\delta c)$ normalizes to $R^2$
(and its tensor square to $R^4$); this separates its character
points and yields the other positive cone.
The actual four-point Tate construction retains the rank-one action
after the stated localization. Algebraic localization, coherent
normalization, and a trace-compatible tensor realization are three
specified operations here; none is inferred from the others.


\subsection{Finding-by-finding transfer table}
{\small\begin{longtable}{p{.12\linewidth}p{.48\linewidth}p{.32\linewidth}}
\toprule IDs&Exact Grothendieck-framework form&Locator/anchor and limit\\\midrule\endhead
F01--F04, F17&Trace of external tensor powers on the path hypersurface; parity of $R\Gamma_c$; Tate duality; optional finite base extension for each $k$. The calculations of \S\ref{sec:main} apply after the trace-formula identification.&SGA~4 XVII \S3, XVIII 3.2.5; SGA~5 III 4.7; Weil~II 3.3.1. Cote~16 provides only the distinct projector/duality framework, not this theorem.\\
F05&The zero virtual class $[\Q_\ell(-1)]-[\Q_\ell(-1)]$ has invisible large eigenvalues. A pure arithmetic Weil module can be a nontrivial extension.&Triangulated $K_0$; BBD 5.3.8 and its following example, pp.~138--139. No archival proof of the new examples is asserted.\\
F06--F07&Effective graded simple classes yield nonnegative multiplicities; a positive phase and Lefschetz maps are additional. Euler classes alone do not suffice.&BBD 5.3.2, 5.4.5, 5.4.10; cote~16, scans~35, 37, 39 for operator language.\\
F08&The rank vector bundles on two points form a special $\lambda$-ring; an effective integral basis need not be its simple positive basis.&SGA~6 V \S\S2--3; cote~133, scan~18. The entire counterexample is proved here.\\
F09--F10&An algebra inclusion need not preserve the effective monoid of a split $K_0$ categorification; exclusive monomial pivots give a sufficient alternative.&Split Grothendieck groups, effective objects, grading shifts; cote~90, scan~17 for equivariant-module context. We do not attribute the nonstandard inclusion claim to the manuscripts.\\
F11&Virtual exterior powers fail effectivity; noncommutative $p$th powers fail additivity, so do not define a scheme Frobenius.&SGA~3 I 4.2, 4.7.2--3; SGA~6 V \S\S2--3; cotes~97/4 and~133/18.\\
F12--F13&For the integral Specht models, invariant projectors are direct summands over $\Z[1/|H|]$; positive restriction matrices compose in representation $K_0$.&SGA~3 I 4.7.2--3; cote~90/17. Actual theta/GCT functor is still absent.\\
F14&Frobenius-stable subquotients inherit a weight; endomorphism spectra are ratios.&Weil representations; BBD 5.3.1 for perverse subquotient weight bounds. Realization/stability is still required.\\
F15&Kernel-defined orbit schemes over $\Z$; coherent polarization and linearization retained in $K_0/G_0$; finite-field base change is separate.&EGA~I \S9.5; FGA~221, 3.1--3.2; cote~90/17 for group actions.\\
F16&The centered polynomial of $\mathbf1\oplus\mathbb L^{\oplus2}\oplus\mathbb L^2$ realizes the rank-one product coefficient; its action comparison fails.&Projective bundle/K\"unneth calculations; cote~16 correspondence language. Pure Tate pieces, not an ungraded pure weight-zero module.\\
F18&Homogeneous weight-lattice difference arrays force ordinary support saturation; a rational half-point retains parity.&SGA~3 I4.7.3 for gradings; exact proof F18; VI Definition1.2.4 is a distinct residuewise notion.\\
F19, F21, F25&A finite flat center, a vector bundle endomorphism algebra and the actual coherent right tensor module; coherent Ext measures the collision.&Cote90/17; affine scheme/module dictionary; explicit complete proofs in these findings. Branch modules belong to coherent $G_0$.\\
F20&Primitive Schur projectors on convolution have effective pure Tate coefficient motives with fixed twists.&Richarz--Scholbach 4.1,4.8,5.9,6.2--6.5; BBD weight formalism; cote16 projector vocabulary only.\\
F22&Universal homeomorphisms preserve the ordinary coefficient topos while conductor and coherent Ext vary.&EGA~IV18.1.2; Stacks04DZ,03SI,0A4E.\\
F23, F28&Effective induced representation classes yield orbit/Pieri formulas when their stabilizer hypotheses hold; an isolated rectangular hook excludes that Young-orbit model.&SGA~3 representations; cote133/19--21 only a finite-set-operation anchor. IP Proposition6.4 supplies the hook input.\\
F24&Even external tensor powers of visible signed Euler spectra retain a positive largest-radius coefficient.&SGA~4/5 trace formalism; full proof F24, with powered-spectrum visibility and visible duality.\\
F26--F27&Finite Weil resolvents are rational; non-Laurent extension evaluations are excluded. Bounded Hecke complexes and normalized character schemes have different effective cones.&Split/triangulated $K_0$; finite normalization; cote1/77--78 gives generic Fredholm vocabulary, not this finite trace identity.\\\bottomrule
\end{longtable}}

\paragraph{\fid{F28}: rectangular projectors in representation $K_0$.}
The unprojected multiplicity module is the actual constant-group
representation $[(a^b)]\otimes[(a^b)]$. A proposed monomial orbit
basis would be an effective sum of induced one-dimensional Young
subgroup representations. Exterior powers of the standard
representation compute their consecutive hook supports. The
isolated hook supplied by IP Proposition~6.4 contradicts such an
effective sum; this proof carries verbatim to representation $K_0$.
A virtual sum may cancel those hooks and is a different assertion.
The known Manivel algebra instead has a torus-graded monomial basis,
with the displayed determinant-character comparison. This belongs
to the affine-group and weight-grading framework of SGA~3, I,
4.7.2--3 and cote~90, scan~17. Cote~133's operations on finite sets
do not identify these two different positive structures.


\paragraph{\fid{F29}: flag schemes and orbit dimensions.}
The action is that of the split group scheme $SL_n\times SL_n$ on
the partial flag scheme of its tensor product representation,
base-changed to $\C$. These flag schemes are iterated relative
Grassmannians, hence instances of the locally free Quot framework
in FGA, Expos\'e~221, \S3, Theorems~3.1--3.2. Their filtered
frame bundles give the Borel torsors used in the proof; the group
actions belong to SGA~3, I, \S4. The complete dimension and
stabilizer calculation proves only the absence of an open orbit for
this action. Cote~90, scan~17, anchors equivariant affine-module
language, not an attribution of this orbit theorem to the archive.

\paragraph{\fid{F30}: vector bundles and coherent commutants under specialization.}
In the language of schemes, $M_R$ is a vector bundle on $\Spec R$
with an action of a finite coherent algebra. Images of the regular
idempotents above are coherent direct summands and hence vector
subbundles. The regular mixing is an automorphism of the underlying
bundle; it does not commute with the whole algebra action. The coherent
commutant is the kernel of the finitely many linear commutation
equations for the generators and is the scalar sheaf in this case.
Its formation does not commute with the special fiber: that fiber
has a second classical projector. This is an explicit base-change
phenomenon for coherent endomorphisms, compatible with Grothendieck's
scheme and module language. It supplies neither a Frobenius action nor
a weight statement, and purity alone does not decide which algebra
equivariance a proposed comparison should retain.

The framework is the affine module/sheaf equivalence of EGA~I,
\S9.5, and equivariant modules in SGA~3, I, 4.7.2--3; cote~90,
scan~17, is the audited archive anchor. Every needed idempotent
image, scalar commutant and special-fiber calculation is proved
in \fid{F30}; no general base-change theorem is asserted without
its hypotheses.

\paragraph{\fid{F31} in Grothendieck's language.}
The integral Grassmannian is the projective Quot scheme of
rank-$d$ locally free quotients of $\mathcal O_{\Z}^{n+d}$,
with its universal sequence and Pl\"ucker line. FGA, Expos\'e~221,
\S3, Theorems~3.1--3.2, pp.~259--260, supplies the representability
framework. Its split $\GL_{n+d}$ action and the cocharacter used
above belong to SGA~3, I, 4.7.2--3. The cokernel is an actual
Weil object and the projector an actual rational correspondence,
with effective class $g_r[\mathbb L^r]$. The inverse exists here
because of the full matrix Hom spaces of the Tate motive, while
the integral Schubert cone need not be preserved. Cote~16,
scans~35, 37 and~39, gives the audited Lefschetz/primitive-projector
vocabulary; these leaves are not asserted to contain this
Kronecker identification or its proof.


\paragraph{\fid{F32}: completed Tate classes and the limits of visible duality.}
For a countable direct sum of Tate pieces whose Euler trace has a
convergent nonnegative absolute moment, the theorem specifies exactly
the discrete completion being used. It is not an element of ordinary
finite $K_0$ merely by notation. Tensor multiplication is legitimate
because the absolute-moment estimate proves convergence and finite
grouping at each nonzero root. The formal binomial expansion of
$(1+\mathbb L^{-1})^{-4}$ in the actual trace has such a realization
at the level of convergent Euler spectra; its inverse Tate degrees
are unbounded. Dualization sends them to unbounded positive degrees,
outside this bounded-spectrum class. Thus an effective or signed
completed class still needs a specified duality and trace functor.
The finite trace framework is SGA~4\,$\frac12$, \emph{Rapport},
3.2, and SGA~5, III, 4.7. Cote~1, scans~77--78, records the
audited generic Fredholm trace/determinant vocabulary; neither it
nor finite $K_0$ supplies this countable extension. Its convergence
and reciprocal-finiteness proof is given in \fid{F32} itself.

\paragraph{\fid{F35} in Grothendieck's language.}
The ordinary split group schemes and quotient presentation are
explicit. Positive-dimensional stabilizers are retained, so this
construction uses the Artin-stack extension of the six operations,
not the bounded constructible pushforward theorem for schemes.
Its rational mass and infinite alternating Tate trace are connected
by the precisely cited Laszlo--Olsson/Sun formalism. The effective
cohomology groups in \eqref{eq:stackcompactgroups} have positive
dimensions, while their grouped Euler coefficients have both signs.
The duality exists but does not preserve the chosen convergent
geometric-Frobenius trace class. The audited manuscript determinant
and projector vocabulary is only a framework anchor; no specific
stack, nonstandard action, or proof here is attributed to those leaves.

\paragraph{\fid{F38} in Grothendieck's language.}
The projective-line decomposition and the stack's actual Tate
cohomology distinguish a positive object from its virtual Euler
class. A rational mass can invert an Euler trace while leaving
unbounded nonzero cohomology. Thus localization of a coefficient
ring or of a convergent trace algebra does not itself provide an
invertible geometric object. The split $\GL_2$ group scheme,
projective bundle formula, and modern stack six operations are
the precise frameworks used here. No archive leaf is claimed to
contain this construction or to identify it with a nonstandard
canonical basis. A different categorical localization requires a
specified new realization and its own comparison theorem.

\paragraph{\fid{F39} in Grothendieck's language.}
The singular piece is the actual affine scheme
$\Spec k[\epsilon]/\epsilon^2$, and matrix Morita equivalence
expresses its coherent category without changing the defect.
Smoothness here asks for a perfect diagonal bimodule, a perfect-complex
condition distinct from finite-dimensionality or ordinary $\ell$-adic
purity. The displayed resolution proves its failure directly. EGA~I's
affine scheme language and the perfect-complex framework of SGA~6
provide the classical dictionary; the smooth proper dg and
noncommutative-motive extension is explicitly the modern cited theory.
The audited cote~16, scans~48--49, concerns conditional semisimplicity
and Weil statements in its own motivic framework; it is not a source
for this dg construction. A new realization and basis comparison is
required before its Frobenius can address the GCT coefficients.


\paragraph{\fid{F34}: effective representation classes and growth of graded algebras.}
The actual highest-weight multiplicities are ranks of effective Schur
summands, in the characteristic-zero representation $K_0$ of $\GL_2$.
Their exact positive simplex has dimension $k-2$; purity cannot alter
that growth degree. The polynomial bound fails only for the specified
compact encoding. Locally free Schur operations and their virtual
extensions use the SGA~3 representation and SGA~6 $\lambda$-ring
frameworks. Cote~133, scans~18--21, anchors those operations, not the
computational bound. The dimension and output-length arguments are
proved explicitly in \fid{F34}; they are not archival theorems.

\paragraph{\fid{F37}: effective multiplicity objects with compressed labels.}
The known rectangular identity identifies the marked-tree count with
the rank of the primitive Tate summand constructed in \fid{F31}.
Compression is a bijection on the finite indexing sets for each input,
so it changes neither that effective representation class nor its rank.
It gives a polynomial witness encoding under the listed-row convention;
it does not identify the trees with a canonical cycle basis or make their
recursion into a theta correspondence. Cote~133, scans~19--21, is a
framework anchor for finite-set operations, and cote~16, scans~35,37,39,
for primitive projectors. The exact counting input is PPS, Theorem~4.5;
the new encoding argument is the proof of \fid{F37} itself.



\paragraph{\fid{F36}: coherent Frobenius, normalization and descent.}

This is a calculation of finite maps, coherent modules, relative
Frobenius and derived tensor products on an explicit singular scheme.
The relative map from collision order $mp$ to order $m$ is a finite
universal homeomorphism, like those in \fid{F22}; its coherent cokernel
is nevertheless nonzero. Ordinary constant \'etale traces lose this
length. The normalization has unit coherent Frobenius and preserves
the original algebra action after restriction, while pushforward
through the collision does not preserve the unit property.


The affine module/sheaf language is EGA~I, \S9.5; the action and
equivariant-module language is SGA~3, I, 4.7.2--3, with the audited
cote~90, scan~17, as its archival anchor. The derived calculations
are the explicit resolutions in the proof, not an invocation of
flat base change at the singular point. No archival identification
of coherent absolute Frobenius with a nonstandard Weil realization
is asserted.


\paragraph{\fid{F34}, constructive supplement: the actual graded ray scheme.}
The actual scheme $\Spec B$ is affine $(k-1)$-space with positive
grading $2,\ldots,k$; $\operatorname{Proj}B$ is the corresponding
weighted projective scheme. The degree-$N$ graded coordinate piece
is the highest-weight multiplicity space just proved, so its Hilbert
series is $\prod_{i=2}^k(1-z^i)^{-1}$. The exponent simplex has
dimension $k-2$. These are exact statements about the affine
coordinate ring and its grading, without an identification of
coherent Hilbert multiplicities with arbitrary constructible
cohomology. The ordinary affine and graded-scheme dictionary applies;
no arithmetic weight, designated nonstandard basis, or archival
claim is supplied by this coordinate computation alone.

The affine dictionary is EGA~I, \S9.5, and the highest-weight representation dictionary uses SGA~3, I, \S4. The separate source anchors for Schur operations are recorded above; no complexity bound is ascribed to those texts.


\paragraph{\fid{F40} in Grothendieck's language.}
Affine endomorphism algebras, idempotent corners and perfect
diagonals give the precise coherent dictionary. The explicit
finite diagonal resolution proves smoothness of a larger category,
while its projective objects retain the singular corner and its
actual module. This uses the perfect-complex language associated
with SGA~6 and the explicitly cited modern smooth dg and
noncommutative-motive theory. The exceptional decomposition
explains the unit Frobenius spectrum without turning the corner
subcategory into an admissible motivic summand. No unpublished
archive leaf is claimed to contain this construction. A relative
order, a realization of the designated morphisms and a comparison
of the required bases remain separate mathematical tasks.

Exact classical framework locators are EGA~I, \S9.5, and SGA~6, Expos\'e~V, \S\S2--3, the latter retained as background with the source-access qualification in the archival audit. Cote~16, scans~48--49, is a vocabulary anchor for conditional motivic semisimplicity and Weil statements, not for this smooth dg construction.


\paragraph{\fid{F33} in Grothendieck's language.}
The geometric object is the ordinary split group scheme $SL_m$.
Its rational Tate-motive decomposition supplies actual Weil
summands and scheme Frobenius, so the equality of dimensions
has a precise realization. The relevant frameworks are SGA~3's
split group schemes and the explicitly cited modern motive
theorem. The generator degree $2i-1$ and weight $2i$ remain
different gradings. No manuscript leaf is claimed to contain
this hook formula or a nonstandard basis comparison.


The group-scheme language is SGA~3, I, \S4; cote~90, scan~17, is the audited group-action anchor. The motive decomposition itself is Merkurjev, Corollary~5.5, and not an assertion from the archive.


\paragraph{\fid{F41} in Grothendieck's language.}
The effective Tate summands and their weights are those of the
actual stack calculation in \fid{F35}. Eigenspace restriction
belongs to ordinary Weil linear algebra, and a finite weight
window restores the finiteness premise of \fid{F26}.
The chosen fixed-parameter action demonstrates that a geometric
coefficient object and an extension-wise parameter comparison
are distinct requirements. No action on the whole stack, or
archive attribution for such an action, is asserted.

The finite trace framework is SGA~4\,$\frac12$, \emph{Rapport}, 3.2, and SGA~5, III, 4.7. The infinite stack input uses the precisely cited modern extension; no bounded scheme theorem is applied to the unbounded object.


\paragraph{\fid{F42} in Grothendieck's language.}
The finite-resolvent obstruction concerns one finite realization
and its fixed Frobenius, whereas the Tate coefficient objects
use tensor powers with growing spectra. The fixed-parameter
construction is an exact action in the category of point Weil
objects, with named generators and tensor compatibility; it is
not an extension-wise quantum grading comparison. Scalar
division, an idempotent image, and inversion of a tensor object
are distinct operations here. The remaining path comparison
must preserve the actual operations and control its boundary
complex; no such comparison is attributed to an archive leaf.

The exact trace-framework locators are SGA~4\,$\frac12$, \emph{Rapport}, 3.2, and SGA~5, III, 4.7; the tensor/action language uses SGA~3, I, 4.7.2--3. Cote~90, scan~17, supplies the audited action vocabulary, not the comparison axiom.

\section{Archival audit and the Fable leads}\label{sec:archive}
The original Montpellier PDFs add a cover. For the documents inspected here, a bottom-stamped archive scan $s$ is one-based PDF page $s+1$. Catalogue titles and scanned contents are distinguished. The following checks were visual readings of the specified leaves, not certifications of full manuscript proofs.

{\small\begin{longtable}{p{.19\linewidth}p{.36\linewidth}p{.37\linewidth}}
\toprule Cote and scan&Verified content&Consequence for this investigation\\\midrule\endhead
90, scan 17 (PDF 18)&Flat affine group over a base; group action; quasi-coherent equivariant modules.&Authentic anchor for the representation/group-scheme dictionary.\\
97, scan 4 (PDF 5)&Chern classes of regular representations and products over characters.&Characteristic-class anchor, not a plethysm-positivity theorem.\\
133, scan 18 (PDF 19)&Symmetric and exterior generating operations, including $\sigma_t(E)\lambda_{-t}(E)=1$.&Direct language anchor for the $\lambda$-ring calculation.\\
133, scans 19--21 (PDF 20--22)&Virtual finite $G$-sets and operations; scan 20 explicitly names plethysm.&These finite-set operations must not be identified with a proved representation-theoretic positive rule.\\
16, scans 35, 37, 39 (PDF 36, 38, 40)&Handwritten Lefschetz/primitive-summand/projector formulas, including hyperplane-section comparisons.&Relevant standard-conjecture framework. No identification with Alp\"oge's path sign is proved.\\
16, scans 34, 36, 38, 40&Reused typescript backs; scan 34 starts Karoubi's \emph{Cohomologie des cat\'egories de Banach}.&The Fable range ``34--40'' must be refined; these backs are not Grothendieck's handwritten Hodge-index proof.\\
16, scans 47--49 (PDF 48--50)&Letter to Coates: correspondences/duality; on scan 48, $C(X)$ plus $I(X\times X)$ implies correspondence-ring semisimplicity; on scan 49, deduction of Weil from $C$ and $A$, in the letter's notation.&Conditional standard-conjecture statements. The semisimplicity passage is scan 48, not scan 50.\\
16, scan 50 (PDF 51)&Separate handwritten endomorphism-algebra lemma.&Not the location of the preceding semisimplicity sentence.\\
1, archival scans 77--78 (PDF78--79), handwritten folios79--80&Exterior-power/trace-norm estimates; scan78 has a Fredholm determinant product and trace sum.&Generic operator-determinant framework is visible; the heavily revised manuscript proof is not certified.\\
1, archival scans109--110 (PDF110--111), handwritten folios111--112&Entire functions, order/genus and canonical products.&No $\Xi$ identity or RH-equivalent positive operator was verified.\\
1, literal archival scans79--80 and111--112 (PDF80--81,112--113)&Corrected determinant calculations, convex/supporting functions, and entire-function or positive-sequence convergence material.&Both numbering interpretations were checked; these are not interchangeable leaves.\\\bottomrule
\end{longtable}}
Cote~133's catalogue title concerns Witt groups, quadratic forms and formal groups; ``plethysm'' is a content lead, not its official title. Cote~60 was found in the catalogue and inspected at thumbnail level, but no precise theorem from it is used.

The Fable suggestions are treated as follows. Suggestion~3's norm estimate is proved and extended in \fid{F01}; its literal positive-Fourier-operator identification fails by \fid{F05}(e). Suggestion~4's identification of the cohomological sign with a polarized evanescent intersection form remains unproved: parity in a trace formula and a positive-definite form are different structures. Suggestion~8 is made precise by \fid{F13}, conditional on an actual functor and endpoint identification. No assertion about the unprovided earlier runs or their theorem IDs is imported.

Some published-source checks have a narrower scope than the archive checks. The original SGA~6 scan could not be visually retrieved here; Expos\'e~V, \S\S2--3, is carried as background. The specific locator V, Lemma~2.2.1, p.~315, is independently cited by Altman--Kleiman \cite[\S2.10]{AK}; it is a locator check, not a re-proof of SGA~6. The modern transcription consulted as a finding aid is not used as technical authority. All elementary $\lambda$ identities used in the paper have their displayed direct proofs.

\part{Unproved bridges, attempted routes, and current state}
\section{What a sufficient nonstandard package would have to contain}\label{sec:gates}
There are no unchecked affirmative ``new theorems'' in this part. The following is a sufficient research specification, not a quotation from Mulmuley and not a proved construction.

\begin{enumerate}
\item \gate{R: realization and normalization.} For each relevant nonstandard coefficient construct a finite-field Weil object or a defined categorical replacement, with explicit Frobenius, trace formula, degrees and weights. Identify $v$ with the correct grading, and specify signs, half-powers and rational rescalings.
\item \gate{S: spectral sign and duality.} Verify the pointwise input bound, the finite-spectrum identification, the output sign for unbounded $k$, and controlled background terms. F24 allows signed visible inputs at even $k$; its finite-extension version requires retained top-radius visibility. Supply either the Tate self-duality of \fid{F03} or the dual system of \fid{F17}. These inputs imply the stated purity theorem by the checked proof.
\item \gate{E: effective basis and operations.} Identify the normalized coefficient with a graded multiplicity of actual simple objects and prove compatibility with multiplication/comultiplication. Purity alone does not give this. The fixed inherited ambient cone is ruled out by \fid{F09}; F27 restores the actual rank-one character action in a different localized cone. Neither identifies VIII's basis.
\item \gate{L: Lefschetz and defects.} Prove the appropriate Lefschetz maps for unimodality. In general plethysm, define the nonflatness correction and prove the quantitative error bounds of the intended version of VIII 2.9/3.6. A spectral error bound such as \fid{F04} is not automatically that coefficient-vector error bound.
\item \gate{C: canonical and comparison data.} Prove basis existence/globalization, cell and quasi-cell claims, and the required specialization. For general plethysm retain VII's duality/reciprocity hypotheses; for Kronecker do not re-open the already proved revised-IV duality.
\item \gate{P: effective labels under a stated encoding.} Specify the input convention before asking for polynomial witnesses. Finding~\fid{F37} settles the equal-rectangle, two-row-output witness problem with explicitly listed rows. Finding~\fid{F34} excludes a general \(\#P\) claim under fully compact binary plethysm encoding. The attached homogeneous PH1 is already refuted by the cited Kahle--Micha{\l}ek theorem; it is not an open requirement to be proved. Any replacement using inhomogeneous polytopes or another counting model needs its own exact formulation, size bounds and algorithm. SH or the appropriate index bound remains separate.
\item \gate{O: actual obstructions.} Establish a strict multiplicity comparison for the orbit closures in an appropriate asymptotic family. Positive formulas and efficient decisions do not themselves produce that inequality.
\end{enumerate}
The spectral theorems remove an estimate once the explicit trace and sign hypotheses are supplied. F32 extends the upper-radius argument to a specified countable completion, but bounded visible reciprocal duality then forces finiteness. F19--F21 and F27 construct specified coherent or ordinary motivic realizations. F36 constructs action-compatible coherent Frobenius models and computes their descent defect. F35 closes the scalar stack-realization question for the actual normalized word trace; its dual does not satisfy the needed second bound. F38 excludes the proposed object-level tensor inverse in the usual derived Weil realization, and F39 excludes the direct nonsmooth collision category from the smooth proper noncommutative Weil framework. F40 repairs that gate by an explicit smooth proper corner category, retaining the exact specialized action; its whole-category TP spectrum is pure but does not yet realize the designated coefficient morphisms. These results narrow particular realizations; the designated based operations and path-output comparison remain open. Replacing all seven by the single phrase ``nonstandard RH implies positivity'' hides independent claims.

For completeness, an effective-label claim has a simple formal implication: if the labels $\alpha$ in \eqref{eq:comparison} and witnesses for $m_\lambda^\alpha,n_\pi^\alpha$ have polynomial-length unique encodings and polynomial-time tests, then triples consisting of a label and the two witnesses count the sum of products. A nondeterministic machine guesses and checks the triple, giving a \(\#P\) function. This is just the definition of \(\#P\); neither the encodings nor a replacement for the refuted PH1 follows from the abstract multiplicity equality. No complexity-class collapse or lower bound is asserted.

\section{Approaches attempted and their exact outcomes}
{\small\begin{longtable}{p{.23\linewidth}p{.40\linewidth}p{.29\linewidth}}
\toprule Approach&Proved outcome or failure point&Next unresolved input\\\midrule\endhead
Literal source formulation&No finite-field nonstandard RH statement in the supplied versions; actual signed/near-positive targets located.&Choose a precise R/E target.\\
Path orthogonality and repeated Parseval&F01 gives the bound, but the constant trace example defeats purity without spectral signs.&S1/S2 realization.\\
Tensor powers and inverse roots&F02--F04 remove the fixed exponent loss; F17 allows separate finite extensions and a dual system.&Prove the concentration/error and duality hypotheses.\\
Fourier-positive operator&The bare Fourier-after-nonnegative-multiplication claim is false, F05(e).&A different operator and a proved identity to the coefficient.\\
Ordinary ambient purity&Stable subquotients inherit purity, F14. The actual inherited positive Hecke cone has only scalar image, F09.&A new effective realization/heart or different character map.\\
Change of rank-one positive model&F16 loses the actual action; F27 restores it in a localized four-character mixed Tate cone.&Identify the designated basis and operations; the repaired object is not pure of one weight.\\
Positive supports/crystals&Monomial pivots suffice, F10. Global positive disjoint supports and local disjointness alone fail.&Global exclusive pivots or another valid basis argument.\\
Orbit/Hilbert/Quot geometry&Integral models and ordinary IC purity exist; polarized coherent multiplicities are not determined by underlying Weil cohomology, F15.&A coherent-to-constructible multiplicity comparison.\\
Determinants and Gram forms&Positive projectors already express the ranks, F12; chosen-basis coefficient positivity fails in F08.&Image basis, independence and operation compatibility.\\
Theta and interlacing chains&F13 gives path/polytope counts for specified exact functors.&Identify the actual theta chain and endpoints with GCT restriction.\\
$\lambda$-rings and motives&Formal operations and Tate examples translate exactly; virtual effectivity and algebraic projectors remain extra.&Effective classes and the required realization, not only ring identities.\\\bottomrule
\end{longtable}}

\section{States of the five requested tasks}
\begin{enumerate}
\item \textbf{Statement needed.} The source-location and target audit is completed as A01 and \S\ref{sec:targets}. The demand for exact nonstandard Frobenius objects/weights cannot be met from those papers because they are not specified. The replacement package in \S\ref{sec:gates} is explicitly our sufficient specification.
\item \textbf{Template.} The standard geometric dictionary and every missing role are stated in \S\ref{sec:KL}. The actual rank-three coherent algebra and right tensor module are supplied by F19/F21; F20 supplies ordinary pure Tate classical coefficients. The designated nonstandard based realization remains open at R/E.
\item \textbf{Sign-only mechanism.} F01--F04, F17, F24 and the admissible countable extension F32 are proved and checked; omitted-hypothesis versions are refuted by the displayed countermodels. F26 excludes the finite normalized all-extension word dictionary. F35 supplies a genuine stack scalar realization with explicit cohomology, while F38 separates its reciprocal traces from actual tensor inverses. The specified main transfer remains open at R/S/E; a scalar mass identity is not a construction of the operators, path output, dual bound or designated basis.
\item \textbf{Positivity reached.} F31 gives actual primitive Tate motives for equal rectangles with two-row output. F37 gives polynomially verifiable compressed witnesses for the known marked-tree rule under the listed-row encoding. F34 proves the separate compact dimension and output-length obstructions. The supplied homogeneous PH1 is refuted by a cited prior theorem, including an equal-rectangle example. General designated canonical-basis positivity and strict orbit-closure multiplicity inequalities remain open at E/L/C/O; replacement effective rules require the revised P gate.
\item \textbf{Theta positivity.} F12--F13 and F23 state the exact projector/path consequences; F18 and F28 delimit integral-interlacing and Young-orbit versions. The theta/GCT functor and basis identification remain open. F08 refutes the generic determinant-to-basis-positivity shortcut.
\end{enumerate}

Research is continuing on the remaining gates under the scheduled-review rule. The present checkpoint is not a final return. No result here proves or refutes the intended full nonstandard positivity program, or $P\ne NP$.

\appendix
\section{Verification and reproducibility}
The independent spectral audit reconstructed F01, reproved F02 by Ces\`aro orthogonality rather than poles, and checked the two-limit argument of F03. The initial all-conjugate duality gap was recorded and repaired before the final theorem. Later full-paper reviews checked each proof, including the actual Hecke-cone example, the positive-pivot countermodels, the categorical and theta statements, and the scheme/base-change examples. Exact dates, reviewer runs, model identity, exposed effort information and check scopes are in \texttt{CHECK\_EVENTS.tsv}.

The file \texttt{check\_path\_bound.py} checks 54 small scalar/vector examples, including complex weights, for the exact path identity and the estimate. These floating-point checks are supplementary; the proof is F01. The optional exact Murnaghan--Nakayama script and character certificate compute a small rectangular coefficient and check character orthogonality; they are not used in any proof here and are not a Lean certificate. The source manifest records the unchanged attachments by SHA-256. Source reads and whole-proof checks are distinguished from cited background theorems and from formal verification.

The dated ledger records actual changes of gates. Scheduled checkpoints are due only at 45-minute intervals; scheduled progress reviews begin only after approximately three hours. The metadata records which such deadlines, if any, occurred; no review or checkpoint is invented to imply elapsed work that did not occur.

\begin{thebibliography}{99}
\bibitem{Alpoge} L.~Alp\"oge, \emph{Remark}, supplied four-page note, 20 November 2026; \url{https://alpo.ge/remark.pdf}. Main locators: p.~3 footnote~3 and steps~3--5; p.~4 steps~7--8. Supplied formalization entry points: \texttt{ALPOGE\_Main.lean}, \texttt{ALPOGE\_Statements.lean}; scan record \texttt{ALPOGE\_LEAN\_SCAN.tsv}. Results accepted as stipulated; missing imported development not executed.
\bibitem{IV} J.~Blasiak, K.~D.~Mulmuley and M.~Sohoni, \emph{Geometric Complexity Theory IV: nonstandard quantum group for the Kronecker problem}, arXiv:cs/0703110v4, 2013. \url{https://arxiv.org/abs/cs/0703110v4}. Exact statement locators are given in the text; attached source \texttt{Apr13KroneckerGCT4.tex}.
\bibitem{VI} K.~D.~Mulmuley, \emph{Geometric Complexity Theory VI: the flip via saturated and positive integer programming in representation theory and algebraic geometry}, arXiv:0704.0229v4, 2009. \url{https://arxiv.org/abs/0704.0229v4}.
\bibitem{VII} K.~D.~Mulmuley, \emph{Geometric Complexity Theory VII: Nonstandard quantum group for the plethysm problem}, arXiv:0709.0749v2, 2008. \url{https://arxiv.org/abs/0709.0749v2}.
\bibitem{VIII} K.~D.~Mulmuley, \emph{Geometric Complexity Theory VIII: On canonical bases for the nonstandard quantum groups}, arXiv:0709.0751v2, 2008. \url{https://arxiv.org/abs/0709.0751v2}.
\bibitem{IP} C.~Ikenmeyer and G.~Panova, \emph{Rectangular Kronecker coefficients and plethysms in geometric complexity theory}, arXiv:1512.03798v2, 2017; especially Theorems~1.4, 1.6, 1.8, 1.10. \url{https://arxiv.org/abs/1512.03798v2}.
\bibitem{BIP} P.~B\"urgisser, C.~Ikenmeyer and G.~Panova, \emph{No occurrence obstructions in geometric complexity theory}, arXiv:1604.06431v3, 2018; Theorem~1.4, \S1(c), Proposition~3.3, equation~(4.1), Proposition~4.5, Theorem~4.7. \url{https://arxiv.org/abs/1604.06431v3}.
\bibitem{KL} D.~Kazhdan and G.~Lusztig, \emph{Schubert varieties and Poincar\'e duality}, in \emph{Geometry of the Laplace Operator}, Proc. Sympos. Pure Math.~36, 1980, pp.~185--203; Theorem~4.3. The exact normalized form used here is checked through Haines's primary proof below.
\bibitem{Haines} T.~J.~Haines, \emph{A proof of the Kazhdan--Lusztig purity theorem via the decomposition theorem of BBD}; Theorems~2.0.1, 5.0.4, Proposition~3.0.2, \S4. \url{https://www.math.umd.edu/~tjh/KL_purity1.pdf}.
\bibitem{WeilII} P.~Deligne, \emph{La conjecture de Weil II}, Publ. Math. IH\'ES~52 (1980), pp.~137--252; Definitions~1.2.2 and~6.2.2--6.2.4, Theorem~3.3.1 and Corollaries~3.3.4--3.3.6. \url{https://www.numdam.org/item/PMIHES_1980__52__137_0/}.
\bibitem{BBD} A.~Beilinson, J.~Bernstein and P.~Deligne, \emph{Faisceaux pervers}, Ast\'erisque~100 (1982). Exact locators: 5.1.2, 5.1.8, 5.1.14, 5.3.1--5.3.2 (including the second proof), 5.3.8 and the following example, 5.4.4--5.4.5, Remark~5.4.9 (p.~143), 5.4.10 (p.~144), 6.2.5. \url{https://www.numdam.org/item/AST_1982__100__1_0/}.
\bibitem{Molev} A.~I.~Molev, \emph{Gelfand--Tsetlin bases for classical Lie algebras}, Theorem~2.1, equation~(2.2), patterns~(2.4); author-hosted preprint, printed pp.~7--9. \url{https://www.maths.usyd.edu.au/u/pubs/publist/preprints/2002/molev-18.pdf}.
\bibitem{Milne} J.~S.~Milne, \emph{Lectures on \'Etale Cohomology}, v2.21, 2013; Example~16.3, p.~109; Theorem~23.2, pp.~139--140. \url{https://www.jmilne.org/math/CourseNotes/LEC.pdf}.
\bibitem{EGA} A.~Grothendieck and J.~Dieudonn\'e, \emph{\'El\'ements de g\'eom\'etrie alg\'ebrique I: Le langage des sch\'emas}, Publ. Math. IH\'ES~4 (1960), \S9.5. Framework locator, with the affine instance proved in F15. \url{https://www.numdam.org/item/PMIHES_1960__4__5_0/}.
\bibitem{SGA3} M.~Demazure and A.~Grothendieck (eds.), \emph{Sch\'emas en groupes (SGA 3)}, Expos\'e~I by M.~Demazure, \S4.2, 4.7.2--4.7.3. Checked in the Gille--Polo re-edition, respecting explicitly marked editorial additions. \url{https://wstein.org/sga/SGA3/Exp1-20jan2010.pdf}.
\bibitem{SGA4} M.~Artin, A.~Grothendieck and J.-L.~Verdier (eds.), \emph{Th\'eorie des topos et cohomologie \'etale des sch\'emas (SGA 4)}, Expos\'es~XVII, \S3, and XVIII, 3.2.5. Background framework references.
\bibitem{SGA45} P.~Deligne, \emph{Cohomologie \'etale (SGA $4\frac12$)}, LNM~569, 1977; \emph{Rapport sur la formule des traces}, 3.2. Background trace-formula reference.
\bibitem{SGA5} A.~Grothendieck (ed.), \emph{Cohomologie \(\ell\)-adique et fonctions L (SGA 5)}, LNM~589, 1977; Expos\'e~III, 4.7. Background Lefschetz--Verdier reference.
\bibitem{SGA6} P.~Berthelot, A.~Grothendieck and L.~Illusie (eds.), \emph{Th\'eorie des intersections et th\'eor\`eme de Riemann--Roch (SGA 6)}, LNM~225, 1971; Expos\'e~V by P.~Berthelot, \S\S2--3. Original scan not visually verified in this run; used as background, not a claimed fresh source-proof check.
\bibitem{FGA} A.~Grothendieck, \emph{Techniques de construction et th\'eor\`emes d'existence en g\'eom\'etrie alg\'ebrique IV: les sch\'emas de Hilbert}, S\'eminaire Bourbaki, Expos\'e~221 (1961), \S3, Theorems~3.1--3.2, pp.~259--260. \url{https://www.numdam.org/article/SB_1960-1961__6__249_0.pdf}.
\bibitem{AK} A.~B.~Altman and S.~L.~Kleiman, \emph{Foundations of the theory of Fano schemes}, Compositio Math.~34 (1977), pp.~3--47; \S2.10, used only to corroborate the SGA~6 locator. \url{https://www.numdam.org/item/CM_1977__34_1_3_0/}.
\bibitem{Archive} Universit\'e de Montpellier, \emph{Archives Grothendieck}, cotes~1, 16, 90, 97, 133, scans specified in \S\ref{sec:archive}. Catalogue: \url{https://grothendieck.umontpellier.fr/archives-grothendieck/}; original files \url{https://grothendieck.umontpellier.fr/16.pdf}, \url{https://grothendieck.umontpellier.fr/90.pdf}, \url{https://grothendieck.umontpellier.fr/97.pdf}, \url{https://grothendieck.umontpellier.fr/133.pdf}. Cote~1 was checked at eight selected leaves, with both scan and handwritten-folio numbering as recorded in the text.
\bibitem{Fable} Claude Fable 5.2, effort max as reported in its supplied file, \emph{FABLE\_PROGRAM\_SUGGESTIONS\_20261120.md} and \emph{\ldots\_EXPLAINED.md}, suggestions~3, 4 and~8. Unchecked leads, not cited theorems.
\bibitem{Richarz} T.~Richarz, \emph{A new approach to the geometric Satake equivalence}, Documenta Mathematica 19 (2014), 209--246. Theorem~3.1, Proposition~4.1, Theorem~5.1, Lemma~5.4, and their proofs; the numbering differs from arXiv 1207.5314v2. \url{https://ems.press/content/serial-article-files/26224?nt=1}.
\bibitem{MV} I.~Mirkovi\'c and K.~Vilonen, \emph{Geometric Langlands duality and representations of algebraic groups over commutative rings}, Annals of Mathematics 166 (2007), 95--143; arXiv math/0401222v5 includes the later Appendix~B correction. Theorem~7.3, Lemma~7.1, Corollary~7.4. \url{https://arxiv.org/abs/math/0401222}.
\bibitem{StacksNil} The Stacks Project, Theorem~59.45.2, tag~04DZ, and Proposition~59.45.4, tag~03SI, topological invariance of the \'etale site/topos; Lemma~59.91.9, tag~0A4E, finite base change. \url{https://stacks.math.columbia.edu/tag/03SI}.
\bibitem{RS} T.~Richarz and J.~Scholbach, \emph{The motivic Satake equivalence}, Math. Ann.~380 (2021), 1595--1653; arXiv:1909.08322v3. Theorem~4.1, Propositions~4.8 and~5.9, Definition~6.2, Lemma~6.3, Corollary~6.4, Lemma~6.5, Theorem~6.8. \url{https://arxiv.org/abs/1909.08322}.
\bibitem{GW} R.~Goodman and N.~R.~Wallach, \emph{Symmetry, Representations, and Invariants}, author-hosted Chapter~4, Theorem~4.1.18 (Hermite reciprocity), p.~191. \url{https://mathweb.ucsd.edu/~nwallach/chapter4.pdf}.
\bibitem{DK} R.~W.~Donley, Jr., and W.~G.~Kim, \emph{A rational theory of Clebsch--Gordan coefficients}, arXiv:1707.03022v2; Proposition~3.1, equation~(3.4), p.~3; Proposition~8.1, equation~(8.2), p.~11. \url{https://arxiv.org/abs/1707.03022}.
\bibitem{Pieri} A.~Postnikov, \emph{MIT 18.315 lecture notes}, scribed by C.~Lian, 2014; Theorem~7.10, p.~16. Used for the classical horizontal and vertical Pieri rules. \url{https://web.mit.edu/clian/www/315_notes.pdf}.
\bibitem{EW} B.~Elias and G.~Williamson, \emph{The Hodge theory of Soergel bimodules}, Annals of Mathematics~180 (2014), 1089--1136; \S3.5, Theorem~3.5, pp.~1107--1108; Theorem~1.1, p.~1093. \url{https://annals.math.princeton.edu/wp-content/uploads/annals-v180-n3-p06-p.pdf}.
\bibitem{Manivel} L.~Manivel, \emph{A note on certain Kronecker coefficients}, arXiv:0809.3710v1 (2008); Theorem~1, p.~1; \S3, pp.~3--4, for the four-generator invariant algebra. \url{https://arxiv.org/abs/0809.3710}.

\bibitem{BlasiakRep} J.~Blasiak, \emph{Representation theory of the nonstandard Hecke algebra}, arXiv:1201.2209v2 (2012); Theorem~5.1, p.~14; Proposition~5.2, pp.~14--15; Definition~6.1 and Remark~6.2, p.~23. PDF equation~(25) has a different number in the converted HTML. \url{https://arxiv.org/abs/1201.2209v2}.
\bibitem{PPBounds} I.~Pak and G.~Panova, \emph{Bounds on certain classes of Kronecker and $q$-binomial coefficients}, author manuscript dated 13 October 2016; Lemma~1.3, p.~2. \url{https://www.math.ucla.edu/~pak/papers/QBC23.pdf}.
\bibitem{Brosnan} P.~Brosnan, \emph{On motivic decompositions arising from the method of Bia\l ynicki--Birula}, Invent. Math.~161 (2005), 91--111; inspected author PDF \S2, pp.~3--4, Theorems~3.3 and~3.5, pp.~5--6. \url{https://math.umd.edu/~pbrosnan/Papers/bb.pdf}.


\bibitem{KM} T.~Kahle and M.~Micha{\l}ek, \emph{Plethysm and lattice point counting}, Electronic Journal of Combinatorics~25(1) (2018), P1.41; Theorem~5 and Remark~6, p.~4; formulation history, pp.~2--4; rectangular identification, p.~6. \url{https://www.combinatorics.org/ojs/index.php/eljc/article/download/v25i1p41/pdf/}.
\bibitem{PPS} I.~Pak, G.~Panova and J.~P.~Swanson, \emph{A combinatorial interpretation for certain plethysm and Kronecker coefficients}, arXiv:2511.02312v3, uploaded 27 July 2026, title-page date 28 July 2026; Lemma~1.3, p.~2; \S2.2, p.~5; Definition~3.2 and Proposition~3.4, p.~6; Lemma~4.2, p.~9; Definition~4.4 and Theorem~4.5, p.~10. \url{https://arxiv.org/abs/2511.02312v3}.
\bibitem{SunStacks} S.~Sun, \emph{L-series of Artin stacks over finite fields}, Algebra \& Number Theory~6 (2012), 47--122; Theorem~1.1, p.~49; Definition~4.1, p.~84; Theorem~4.2, p.~85; Theorem~4.8 and \S4.8.2, pp.~89--90; Example~7.4, p.~109. \url{https://msp.org/ant/2012/6-1/ant-v6-n1-p03-s.pdf}.
\bibitem{LO2} Y.~Laszlo and M.~Olsson, \emph{The six operations for sheaves on Artin stacks II: Adic coefficients}, Publ. Math. IH\'ES~107 (2008), 169--210; \S9.1, Lemmas~9.1.1--9.1.2, p.~198; Theorem~11.0.14, p.~204. \url{https://numdam.org/item/10.1007/s10240-008-0012-5.pdf}.
\bibitem{TabuadaWeil} G.~Tabuada, \emph{Noncommutative Weil conjecture}, inspected version arXiv:1808.00950v7 (2019); \S1, pp.~2--3, Theorem~1.5; \S5.1, pp.~13--14; \S6.1, pp.~15--16; \S7, equations~(7.1)--(7.4), p.~17; Example~12.3, p.~24; Theorem~2.12, p.~8, and proof equation~(12.6), p.~26; \S6.2, equation~(6.9), p.~16. \url{https://arxiv.org/abs/1808.00950v7}.

\bibitem{MerkurjevSL} A.~S.~Merkurjev, \emph{Motivic decomposition of certain special linear groups}, inspected author PDF, \S5, Corollary~5.3 and notation, p.~7; Corollary~5.5, p.~8. \url{https://www.math.uni-bielefeld.de/LAG/man/530.pdf}.
\bibitem{Run} This run, \texttt{GCT\_NSRH\_20261121}; findings A01 and F01--F42, with agent source memoranda, proof reviews, computations and \texttt{CHECK\_EVENTS.tsv} in the final package. This is the only current-run identifier used.
\end{thebibliography}
\end{document}
